\section{Results}
\label{sec: Results}

\subsection{Five-Circuit Coast-Down and $C_LA$ Results}
\label{sec: Five Circuit Results}

Table~\ref{tab: five circuit results} summarises the aerodynamic parameters extracted at each of the five selected circuits. All results use driver~14 (Alonso, AMR24) with the identical fitting methodology (\S\ref{sec: Methodology}).

\begin{table*}[h]
    \centering
    \caption{Aerodynamic parameters across five circuits (driver~14, AMR24, 2024). $\rho$ is the mean free-practice air density. $n_\text{segs}$ is the coast-down segment count; $n_\text{CLA}$ the qualifying push-lap count. $C_LA$ total uncertainty combines statistical SE of the median and $\mu$-systematic in quadrature (see \S\ref{sec: CLA Depends on Assumed mu}).}
    \label{tab: five circuit results}
    \begin{tabularx}{\linewidth}{l C C C C C c c}
        \toprule
        Circuit & $\rho$~(kg/m$^3$) & $\hat\alpha$~(N\,s$^2$/m$^2$) & Composite $2\hat\alpha/\rho$~(m$^2$) & $C_LA$~(m$^2$) & $C_DA$~(m$^2$) & $n_\text{segs}$ & $n_\text{CLA}$ \\
        \midrule
        Jeddah 2024      & 1.180 & $0.876 \pm 0.004$ & 1.485 & $2.66 \pm 0.42$ & $1.448 \pm 0.009^\dagger$ & 13 & 5 \\
        Spa 2024         & 1.145 & $0.812 \pm 0.007$ & 1.419 & $3.06 \pm 0.48$ & $1.376 \pm 0.014^\dagger$ &  5 & 9 \\
        Silverstone 2024 & 1.193 & $0.913 \pm 0.004$ & 1.532 & $3.14 \pm 0.51$ & $1.487 \pm 0.010^\dagger$ & 10 & 6 \\
        Yas Marina 2024  & 1.178 & $0.817 \pm 0.008$ & 1.386 & $3.53 \pm 0.55$ & $1.335 \pm 0.016^\dagger$ &  6 & 6 \\
        Suzuka 2024      & 1.223 & $1.071 \pm 0.008$ & 1.750 & $4.33 \pm 0.67$ & $1.689 \pm 0.017^\dagger$ &  8 & 4 \\
        \bottomrule
    \end{tabularx}
    \vspace{2pt}
    {\small $C_{rr} \approx 0.014$ at all circuits (by construction; see \S\ref{sec: Coast Down ODE Model}). $C_LA$ uncertainty: stat SE of median $\oplus$ $\mu$-systematic. $^\dagger$$C_DA$ $\sigma$ values are Jacobian lower bounds; bootstrap-corrected CIs and the conservative $13\times$ scaling used in the polar are derived in \S\ref{sec: Absolute CDA Interpretation}.}
\end{table*}

The composite $2\hat\alpha/\rho$ ranges from $1.386$~m$^2$ (Yas Marina) to $1.750$~m$^2$ (Suzuka), consistent with Suzuka running significantly more wing than any other circuit in the set. $C_LA$ increases monotonically from $2.66$~m$^2$ (Jeddah) to $4.33$~m$^2$ (Suzuka), as expected for a low-to-high downforce progression.

One anomaly is visible in the $C_DA$ column: Yas Marina ($1.335$~m$^2$) sits below Spa ($1.376$~m$^2$) despite having a higher $C_LA$, which would be physically inconsistent under pure wing-angle variation. Engine-braking contamination of the coast-down segment pool was hypothesised as the cause. A gear-stratified falsification test (see \texttt{scripts/yas\_marina\_gear\_strat.py}) was performed: of 6 passing segments, 2 are in gear~8; the gear-8-only pooled fit yields $\hat\alpha = 0.789$, a $-3.4\%$ residual relative to the all-segment value of $0.817$. This is well within the bootstrap 90\% CI half-width of $\pm 0.101$, so engine-braking contamination is \emph{not} supported by the data (Jacobian-to-bootstrap ratio at Yas Marina: $12.0\times$, consistent with other circuits). The origin of the anomalous $C_DA$ ordering at Yas Marina therefore remains unexplained by the available data; the circuit is retained in the polar as a geometric outlier with an unresolved cause. Excluding Yas Marina shifts the polar slope from $3.02$ to $4.80$, but the confidence interval includes zero in both cases and the qualitative interpretation (positive slope, indeterminate efficiency) is unchanged; see \S\ref{sec: Aero Polar}.

\subsubsection{Reference Circuit Validation (Jeddah)}
\label{sec: Jeddah Validation}

Jeddah is used as the reference circuit for validation tests because it has the largest coast-down segment count (13 of 13 passing $R^2 \geq 0.90$).

\textbf{Gear-stratified falsification.} Nine of the thirteen segments are in gear~8. The gear-8-only pooled fit yields $\hat\alpha = 0.867$ ($-1.0\%$ relative to the all-segment value of $0.876$), not statistically distinguishable from zero within the bootstrap 90\% CI ($\pm 3.1\%$). Full interpretation and physical reasoning are in \S\ref{sec: Absolute CDA Interpretation}.

\textbf{Inter-driver spread.} Driver~18 (Stroll) has three passing segments at Jeddah; the independent pooled fit yields $\hat\alpha = 0.968$, a $10.4\%$ spread relative to driver~14. Because driver~14 is used at \emph{every} circuit in this analysis, this spread is not a source of cross-circuit bias; it is reported as a bound on intra-team setup variation (balance preferences, wing asymmetry, mechanical setup).

\textbf{Bootstrap uncertainty.} A segment-resampling bootstrap ($n = 1000$ replicates; \S\ref{sec: Absolute CDA Interpretation}) gives a 90\% CI of $[0.847, 0.902]$ on $\hat\alpha$, a half-width of $\pm 0.027$~N\,s$^2$/m$^2$ --- 6.8 times the Jacobian-derived bound of $\pm 0.004$. This is lower than the $\sim\!13\times$ ratio observed at Monza in earlier analysis (the Monza ratio was re-evaluated to account for a factor-of-two correction to the Jacobian covariance formula; the corrected Jacobian bound scales by $\sqrt{2}$ while the bootstrap estimate is unchanged, reducing the ratio from the previously reported $19\times$). The conservative $13\times$ scaling is retained for the polar uncertainty (see \S\ref{sec: Aero Polar}; full treatment in \S\ref{sec: Absolute CDA Interpretation}) to avoid under-stating $C_DA$ uncertainty at circuits where the ratio was not independently calibrated.

\subsection{Aero Polar}
\label{sec: Aero Polar}

Figure~\ref{fig: aero polar}(a) plots $C_LA$ against $C_DA$ for all five circuits. An OLS regression gives:

\begin{equation}
    C_LA = 3.02 \cdot C_DA - 1.09
    \label{eq: ols polar}
\end{equation}

Note that all five $C_DA$ values share a common cross-session mass inconsistency (FP mass $\sim\!855$~kg vs.\ qualifying mass $803$~kg), estimated to inflate each absolute composite by $\sim\!0.10$~m$^2$ (\S\ref{sec: Absolute CDA Interpretation}); this offset is common-mode across circuits and cancels in the polar slope, but bounds absolute accuracy.

The slope estimates the incremental aerodynamic efficiency $\Delta C_LA/\Delta C_DA$---how much additional downforce each unit of additional drag area delivers as wing angle increases---and is distinct from the total $L/D$ area ratio ($C_LA/C_DA \approx 2.5$ at Jeddah). The positive slope is physically consistent with wing-angle variation: adding downforce wing area increases both $C_LA$ and $C_DA$ simultaneously. A Monte Carlo simulation ($N = 50\,000$ replicates) propagates the $C_DA$ uncertainty (Jacobian $\alpha_\text{std}$ scaled by the conservative $13\times$ bootstrap ratio) and the total $C_LA$ uncertainty (statistical SE of median combined with $\mu$ systematic in quadrature) to give a 90\% confidence interval on the slope of $[-2.3,\, +6.0]$~m$^2$/m$^2$ ($1.08\,\sigma$ from zero; MC standard deviation $= 2.80$). (The $13\times$ scaling derivation and its Jeddah calibration are given in full in \S\ref{sec: Absolute CDA Interpretation}.)

\textbf{Robustness check: Deming regression.} OLS is inconsistent when the regressor ($C_DA$) carries non-negligible measurement error relative to its range. After the $13\times$ bootstrap correction the $C_DA$ uncertainties are $0.12$--$0.22$~m$^2$ per circuit, which is $33$--$62\%$ of the total $C_DA$ spread of $0.35$~m$^2$; errors-in-variables attenuation bias is therefore substantial in this dataset. Deming regression (orthogonal regression accounting for independent errors in both axes; see \texttt{scripts/deming\_regression.py}) with per-circuit variance ratio $\lambda_i = (\sigma_{C_DA,i}/\sigma_{C_LA,i})^2$ gives a slope of $6.95$~m$^2$/m$^2$ (bootstrap 90\% CI: $[-9.3,\, +8.8]$~m$^2$/m$^2$; $10\,000$ resamples), compared to the OLS estimate of $3.02$. The upward shift is consistent with the expected direction of attenuation correction. The OLS result~(Eq.~\ref{eq: ols polar}) is retained as the primary reported value for continuity with the polar figure. Two features of the Deming result deserve note. First, the CI widened substantially relative to OLS --- from $[-2.3, +6.0]$ to $[-9.3, +8.8]$~m$^2$/m$^2$ --- because with only $n=5$ points, the Deming closed-form amplifies the leverage of individual outlier resamples; this widening is expected and is not a deficiency of the method. Second, the bootstrap distribution is left-skewed (the 5th-percentile tail extends further from the point estimate than the 95th), driven by resamples that over-weight Yas Marina. Both effects are artefacts of the small sample size, and neither changes the qualitative conclusion: the confidence interval includes zero and the slope is poorly determined at this sample size regardless of estimator.

Note that OLS on five data points yields 3~degrees of freedom for the slope; on four points (Yas Marina excluded), 2~DOF. At these sample sizes the $t_{0.05}$ multiplier ($2.35$ and $2.92$, respectively) materially widens the frequentist interval relative to a Gaussian approximation, so the MC 90\% CIs should be read as approximate.

The confidence interval includes zero, so a precise aerodynamic efficiency cannot be stated. The improvement over earlier analysis is qualitative: the five-point polar using a single driver and data-driven circuit selection recovers the correct physical sign (positive slope), whereas the prior three-circuit analysis including Monza---where qualifying $C_LA$ was structurally unreliable (only 3 valid laps)---produced a spurious negative slope ($-2.1$, $0.43\,\sigma$).

\textbf{Robustness check.} Yas Marina is flagged as a geometric outlier in the polar (open diamond, panel~(a)): its fitted $C_DA = 1.335$~m$^2$ is the lowest of the five circuits yet its $C_LA = 3.53$~m$^2$ is the second highest, a physically inconsistent ordering under pure wing-angle variation. Engine-braking contamination was hypothesised as the cause but was ruled out by the gear-stratified falsification in \S\ref{sec: Five Circuit Results}; the origin of the anomaly is currently unexplained. Figure~\ref{fig: aero polar}(b) repeats the polar excluding Yas Marina. The slope increases to $4.80$~m$^2$/m$^2$ (MC 90\% CI: $[-2.9, +9.1]$; $0.99\,\sigma$ from zero), consistent with the five-circuit result within uncertainty. The sign and approximate magnitude of the slope are robust to this exclusion; the wider CI in the four-circuit case reflects the loss of one data point without a proportional reduction in $C_LA$ uncertainty.

The binding uncertainty in both fits is the $\mu$ systematic ($\pm 0.4$--$0.7$~m$^2$ per circuit), which contributes $\sim\!70\%$ of the MC slope variance. Reducing this---either through independent tyre-friction data or by extending the analysis to circuits where the friction limit is more tightly constrained---would be the highest-leverage improvement to the polar.

\begin{figure}[h]
    \centering
    %% Creator: Matplotlib, PGF backend
%%
%% To include the figure in your LaTeX document, write
%%   \input{<filename>.pgf}
%%
%% Make sure the required packages are loaded in your preamble
%%   \usepackage{pgf}
%%
%% Also ensure that all the required font packages are loaded; for instance,
%% the lmodern package is sometimes necessary when using math font.
%%   \usepackage{lmodern}
%%
%% Figures using additional raster images can only be included by \input if
%% they are in the same directory as the main LaTeX file. For loading figures
%% from other directories you can use the `import` package
%%   \usepackage{import}
%%
%% and then include the figures with
%%   \import{<path to file>}{<filename>.pgf}
%%
%% Matplotlib used the following preamble
%%   \def\mathdefault#1{#1}
%%   \everymath=\expandafter{\the\everymath\displaystyle}
%%   \IfFileExists{scrextend.sty}{
%%     \usepackage[fontsize=10.000000pt]{scrextend}
%%   }{
%%     \renewcommand{\normalsize}{\fontsize{10.000000}{12.000000}\selectfont}
%%     \normalsize
%%   }
%%   \usepackage{mathpazo}
%%   \usepackage{amsmath}
%%   \usepackage[T1]{fontenc}
%%   \makeatletter\@ifpackageloaded{underscore}{}{\usepackage[strings]{underscore}}\makeatother
%%
\begingroup%
\makeatletter%
\begin{pgfpicture}%
\pgfpathrectangle{\pgfpointorigin}{\pgfqpoint{6.200000in}{3.801389in}}%
\pgfusepath{use as bounding box, clip}%
\begin{pgfscope}%
\pgfsetbuttcap%
\pgfsetmiterjoin%
\definecolor{currentfill}{rgb}{1.000000,1.000000,1.000000}%
\pgfsetfillcolor{currentfill}%
\pgfsetlinewidth{0.000000pt}%
\definecolor{currentstroke}{rgb}{1.000000,1.000000,1.000000}%
\pgfsetstrokecolor{currentstroke}%
\pgfsetdash{}{0pt}%
\pgfpathmoveto{\pgfqpoint{0.000000in}{0.000000in}}%
\pgfpathlineto{\pgfqpoint{6.200000in}{0.000000in}}%
\pgfpathlineto{\pgfqpoint{6.200000in}{3.801389in}}%
\pgfpathlineto{\pgfqpoint{0.000000in}{3.801389in}}%
\pgfpathlineto{\pgfqpoint{0.000000in}{0.000000in}}%
\pgfpathclose%
\pgfusepath{fill}%
\end{pgfscope}%
\begin{pgfscope}%
\pgfsetbuttcap%
\pgfsetmiterjoin%
\definecolor{currentfill}{rgb}{1.000000,1.000000,1.000000}%
\pgfsetfillcolor{currentfill}%
\pgfsetlinewidth{0.000000pt}%
\definecolor{currentstroke}{rgb}{0.000000,0.000000,0.000000}%
\pgfsetstrokecolor{currentstroke}%
\pgfsetstrokeopacity{0.000000}%
\pgfsetdash{}{0pt}%
\pgfpathmoveto{\pgfqpoint{0.584940in}{1.184369in}}%
\pgfpathlineto{\pgfqpoint{3.025000in}{1.184369in}}%
\pgfpathlineto{\pgfqpoint{3.025000in}{3.517571in}}%
\pgfpathlineto{\pgfqpoint{0.584940in}{3.517571in}}%
\pgfpathlineto{\pgfqpoint{0.584940in}{1.184369in}}%
\pgfpathclose%
\pgfusepath{fill}%
\end{pgfscope}%
\begin{pgfscope}%
\pgfpathrectangle{\pgfqpoint{0.584940in}{1.184369in}}{\pgfqpoint{2.440060in}{2.333202in}}%
\pgfusepath{clip}%
\pgfsetbuttcap%
\pgfsetroundjoin%
\definecolor{currentfill}{rgb}{0.501961,0.501961,0.501961}%
\pgfsetfillcolor{currentfill}%
\pgfsetfillopacity{0.150000}%
\pgfsetlinewidth{1.003750pt}%
\definecolor{currentstroke}{rgb}{0.501961,0.501961,0.501961}%
\pgfsetstrokecolor{currentstroke}%
\pgfsetstrokeopacity{0.150000}%
\pgfsetdash{}{0pt}%
\pgfsys@defobject{currentmarker}{\pgfqpoint{0.584940in}{1.409229in}}{\pgfqpoint{3.025000in}{3.588434in}}{%
\pgfpathmoveto{\pgfqpoint{0.584940in}{2.504701in}}%
\pgfpathlineto{\pgfqpoint{0.584940in}{1.409229in}}%
\pgfpathlineto{\pgfqpoint{0.597202in}{1.420180in}}%
\pgfpathlineto{\pgfqpoint{0.609463in}{1.431131in}}%
\pgfpathlineto{\pgfqpoint{0.621725in}{1.442082in}}%
\pgfpathlineto{\pgfqpoint{0.633986in}{1.453032in}}%
\pgfpathlineto{\pgfqpoint{0.646248in}{1.463983in}}%
\pgfpathlineto{\pgfqpoint{0.658510in}{1.474934in}}%
\pgfpathlineto{\pgfqpoint{0.670771in}{1.485885in}}%
\pgfpathlineto{\pgfqpoint{0.683033in}{1.496836in}}%
\pgfpathlineto{\pgfqpoint{0.695294in}{1.507786in}}%
\pgfpathlineto{\pgfqpoint{0.707556in}{1.518737in}}%
\pgfpathlineto{\pgfqpoint{0.719818in}{1.529688in}}%
\pgfpathlineto{\pgfqpoint{0.732079in}{1.540639in}}%
\pgfpathlineto{\pgfqpoint{0.744341in}{1.551589in}}%
\pgfpathlineto{\pgfqpoint{0.756603in}{1.562540in}}%
\pgfpathlineto{\pgfqpoint{0.768864in}{1.573491in}}%
\pgfpathlineto{\pgfqpoint{0.781126in}{1.584442in}}%
\pgfpathlineto{\pgfqpoint{0.793387in}{1.595392in}}%
\pgfpathlineto{\pgfqpoint{0.805649in}{1.606343in}}%
\pgfpathlineto{\pgfqpoint{0.817911in}{1.617294in}}%
\pgfpathlineto{\pgfqpoint{0.830172in}{1.628245in}}%
\pgfpathlineto{\pgfqpoint{0.842434in}{1.639196in}}%
\pgfpathlineto{\pgfqpoint{0.854695in}{1.650146in}}%
\pgfpathlineto{\pgfqpoint{0.866957in}{1.661097in}}%
\pgfpathlineto{\pgfqpoint{0.879219in}{1.672048in}}%
\pgfpathlineto{\pgfqpoint{0.891480in}{1.682999in}}%
\pgfpathlineto{\pgfqpoint{0.903742in}{1.693949in}}%
\pgfpathlineto{\pgfqpoint{0.916003in}{1.704900in}}%
\pgfpathlineto{\pgfqpoint{0.928265in}{1.715851in}}%
\pgfpathlineto{\pgfqpoint{0.940527in}{1.726802in}}%
\pgfpathlineto{\pgfqpoint{0.952788in}{1.737753in}}%
\pgfpathlineto{\pgfqpoint{0.965050in}{1.748703in}}%
\pgfpathlineto{\pgfqpoint{0.977311in}{1.759654in}}%
\pgfpathlineto{\pgfqpoint{0.989573in}{1.770605in}}%
\pgfpathlineto{\pgfqpoint{1.001835in}{1.781556in}}%
\pgfpathlineto{\pgfqpoint{1.014096in}{1.792506in}}%
\pgfpathlineto{\pgfqpoint{1.026358in}{1.803457in}}%
\pgfpathlineto{\pgfqpoint{1.038619in}{1.814408in}}%
\pgfpathlineto{\pgfqpoint{1.050881in}{1.825359in}}%
\pgfpathlineto{\pgfqpoint{1.063143in}{1.836310in}}%
\pgfpathlineto{\pgfqpoint{1.075404in}{1.847260in}}%
\pgfpathlineto{\pgfqpoint{1.087666in}{1.858211in}}%
\pgfpathlineto{\pgfqpoint{1.099928in}{1.869162in}}%
\pgfpathlineto{\pgfqpoint{1.112189in}{1.880113in}}%
\pgfpathlineto{\pgfqpoint{1.124451in}{1.891063in}}%
\pgfpathlineto{\pgfqpoint{1.136712in}{1.902014in}}%
\pgfpathlineto{\pgfqpoint{1.148974in}{1.912965in}}%
\pgfpathlineto{\pgfqpoint{1.161236in}{1.923916in}}%
\pgfpathlineto{\pgfqpoint{1.173497in}{1.934867in}}%
\pgfpathlineto{\pgfqpoint{1.185759in}{1.945817in}}%
\pgfpathlineto{\pgfqpoint{1.198020in}{1.956768in}}%
\pgfpathlineto{\pgfqpoint{1.210282in}{1.967719in}}%
\pgfpathlineto{\pgfqpoint{1.222544in}{1.978670in}}%
\pgfpathlineto{\pgfqpoint{1.234805in}{1.989620in}}%
\pgfpathlineto{\pgfqpoint{1.247067in}{2.000571in}}%
\pgfpathlineto{\pgfqpoint{1.259328in}{2.011522in}}%
\pgfpathlineto{\pgfqpoint{1.271590in}{2.022473in}}%
\pgfpathlineto{\pgfqpoint{1.283852in}{2.033424in}}%
\pgfpathlineto{\pgfqpoint{1.296113in}{2.044374in}}%
\pgfpathlineto{\pgfqpoint{1.308375in}{2.055325in}}%
\pgfpathlineto{\pgfqpoint{1.320636in}{2.066276in}}%
\pgfpathlineto{\pgfqpoint{1.332898in}{2.077227in}}%
\pgfpathlineto{\pgfqpoint{1.345160in}{2.088177in}}%
\pgfpathlineto{\pgfqpoint{1.357421in}{2.099128in}}%
\pgfpathlineto{\pgfqpoint{1.369683in}{2.110079in}}%
\pgfpathlineto{\pgfqpoint{1.381945in}{2.121030in}}%
\pgfpathlineto{\pgfqpoint{1.394206in}{2.131981in}}%
\pgfpathlineto{\pgfqpoint{1.406468in}{2.142931in}}%
\pgfpathlineto{\pgfqpoint{1.418729in}{2.153882in}}%
\pgfpathlineto{\pgfqpoint{1.430991in}{2.164833in}}%
\pgfpathlineto{\pgfqpoint{1.443253in}{2.175784in}}%
\pgfpathlineto{\pgfqpoint{1.455514in}{2.171170in}}%
\pgfpathlineto{\pgfqpoint{1.467776in}{2.166472in}}%
\pgfpathlineto{\pgfqpoint{1.480037in}{2.161775in}}%
\pgfpathlineto{\pgfqpoint{1.492299in}{2.157077in}}%
\pgfpathlineto{\pgfqpoint{1.504561in}{2.152379in}}%
\pgfpathlineto{\pgfqpoint{1.516822in}{2.147682in}}%
\pgfpathlineto{\pgfqpoint{1.529084in}{2.142984in}}%
\pgfpathlineto{\pgfqpoint{1.541345in}{2.138286in}}%
\pgfpathlineto{\pgfqpoint{1.553607in}{2.133589in}}%
\pgfpathlineto{\pgfqpoint{1.565869in}{2.128891in}}%
\pgfpathlineto{\pgfqpoint{1.578130in}{2.124194in}}%
\pgfpathlineto{\pgfqpoint{1.590392in}{2.119496in}}%
\pgfpathlineto{\pgfqpoint{1.602653in}{2.114798in}}%
\pgfpathlineto{\pgfqpoint{1.614915in}{2.110101in}}%
\pgfpathlineto{\pgfqpoint{1.627177in}{2.105403in}}%
\pgfpathlineto{\pgfqpoint{1.639438in}{2.100705in}}%
\pgfpathlineto{\pgfqpoint{1.651700in}{2.096008in}}%
\pgfpathlineto{\pgfqpoint{1.663962in}{2.091310in}}%
\pgfpathlineto{\pgfqpoint{1.676223in}{2.086613in}}%
\pgfpathlineto{\pgfqpoint{1.688485in}{2.081915in}}%
\pgfpathlineto{\pgfqpoint{1.700746in}{2.077217in}}%
\pgfpathlineto{\pgfqpoint{1.713008in}{2.072520in}}%
\pgfpathlineto{\pgfqpoint{1.725270in}{2.067822in}}%
\pgfpathlineto{\pgfqpoint{1.737531in}{2.063124in}}%
\pgfpathlineto{\pgfqpoint{1.749793in}{2.058427in}}%
\pgfpathlineto{\pgfqpoint{1.762054in}{2.053729in}}%
\pgfpathlineto{\pgfqpoint{1.774316in}{2.049032in}}%
\pgfpathlineto{\pgfqpoint{1.786578in}{2.044334in}}%
\pgfpathlineto{\pgfqpoint{1.798839in}{2.039636in}}%
\pgfpathlineto{\pgfqpoint{1.811101in}{2.034939in}}%
\pgfpathlineto{\pgfqpoint{1.823362in}{2.030241in}}%
\pgfpathlineto{\pgfqpoint{1.835624in}{2.025543in}}%
\pgfpathlineto{\pgfqpoint{1.847886in}{2.020846in}}%
\pgfpathlineto{\pgfqpoint{1.860147in}{2.016148in}}%
\pgfpathlineto{\pgfqpoint{1.872409in}{2.011451in}}%
\pgfpathlineto{\pgfqpoint{1.884670in}{2.006753in}}%
\pgfpathlineto{\pgfqpoint{1.896932in}{2.002055in}}%
\pgfpathlineto{\pgfqpoint{1.909194in}{1.997358in}}%
\pgfpathlineto{\pgfqpoint{1.921455in}{1.992660in}}%
\pgfpathlineto{\pgfqpoint{1.933717in}{1.987962in}}%
\pgfpathlineto{\pgfqpoint{1.945978in}{1.983265in}}%
\pgfpathlineto{\pgfqpoint{1.958240in}{1.978567in}}%
\pgfpathlineto{\pgfqpoint{1.970502in}{1.973870in}}%
\pgfpathlineto{\pgfqpoint{1.982763in}{1.969172in}}%
\pgfpathlineto{\pgfqpoint{1.995025in}{1.964474in}}%
\pgfpathlineto{\pgfqpoint{2.007287in}{1.959777in}}%
\pgfpathlineto{\pgfqpoint{2.019548in}{1.955079in}}%
\pgfpathlineto{\pgfqpoint{2.031810in}{1.950381in}}%
\pgfpathlineto{\pgfqpoint{2.044071in}{1.945684in}}%
\pgfpathlineto{\pgfqpoint{2.056333in}{1.940986in}}%
\pgfpathlineto{\pgfqpoint{2.068595in}{1.936289in}}%
\pgfpathlineto{\pgfqpoint{2.080856in}{1.931591in}}%
\pgfpathlineto{\pgfqpoint{2.093118in}{1.926893in}}%
\pgfpathlineto{\pgfqpoint{2.105379in}{1.922196in}}%
\pgfpathlineto{\pgfqpoint{2.117641in}{1.917498in}}%
\pgfpathlineto{\pgfqpoint{2.129903in}{1.912800in}}%
\pgfpathlineto{\pgfqpoint{2.142164in}{1.908103in}}%
\pgfpathlineto{\pgfqpoint{2.154426in}{1.903405in}}%
\pgfpathlineto{\pgfqpoint{2.166687in}{1.898708in}}%
\pgfpathlineto{\pgfqpoint{2.178949in}{1.894010in}}%
\pgfpathlineto{\pgfqpoint{2.191211in}{1.889312in}}%
\pgfpathlineto{\pgfqpoint{2.203472in}{1.884615in}}%
\pgfpathlineto{\pgfqpoint{2.215734in}{1.879917in}}%
\pgfpathlineto{\pgfqpoint{2.227995in}{1.875219in}}%
\pgfpathlineto{\pgfqpoint{2.240257in}{1.870522in}}%
\pgfpathlineto{\pgfqpoint{2.252519in}{1.865824in}}%
\pgfpathlineto{\pgfqpoint{2.264780in}{1.861127in}}%
\pgfpathlineto{\pgfqpoint{2.277042in}{1.856429in}}%
\pgfpathlineto{\pgfqpoint{2.289304in}{1.851731in}}%
\pgfpathlineto{\pgfqpoint{2.301565in}{1.847034in}}%
\pgfpathlineto{\pgfqpoint{2.313827in}{1.842336in}}%
\pgfpathlineto{\pgfqpoint{2.326088in}{1.837638in}}%
\pgfpathlineto{\pgfqpoint{2.338350in}{1.832941in}}%
\pgfpathlineto{\pgfqpoint{2.350612in}{1.828243in}}%
\pgfpathlineto{\pgfqpoint{2.362873in}{1.823546in}}%
\pgfpathlineto{\pgfqpoint{2.375135in}{1.818848in}}%
\pgfpathlineto{\pgfqpoint{2.387396in}{1.814150in}}%
\pgfpathlineto{\pgfqpoint{2.399658in}{1.809453in}}%
\pgfpathlineto{\pgfqpoint{2.411920in}{1.804755in}}%
\pgfpathlineto{\pgfqpoint{2.424181in}{1.800057in}}%
\pgfpathlineto{\pgfqpoint{2.436443in}{1.795360in}}%
\pgfpathlineto{\pgfqpoint{2.448704in}{1.790662in}}%
\pgfpathlineto{\pgfqpoint{2.460966in}{1.785965in}}%
\pgfpathlineto{\pgfqpoint{2.473228in}{1.781267in}}%
\pgfpathlineto{\pgfqpoint{2.485489in}{1.776569in}}%
\pgfpathlineto{\pgfqpoint{2.497751in}{1.771872in}}%
\pgfpathlineto{\pgfqpoint{2.510012in}{1.767174in}}%
\pgfpathlineto{\pgfqpoint{2.522274in}{1.762476in}}%
\pgfpathlineto{\pgfqpoint{2.534536in}{1.757779in}}%
\pgfpathlineto{\pgfqpoint{2.546797in}{1.753081in}}%
\pgfpathlineto{\pgfqpoint{2.559059in}{1.748384in}}%
\pgfpathlineto{\pgfqpoint{2.571321in}{1.743686in}}%
\pgfpathlineto{\pgfqpoint{2.583582in}{1.738988in}}%
\pgfpathlineto{\pgfqpoint{2.595844in}{1.734291in}}%
\pgfpathlineto{\pgfqpoint{2.608105in}{1.729593in}}%
\pgfpathlineto{\pgfqpoint{2.620367in}{1.724895in}}%
\pgfpathlineto{\pgfqpoint{2.632629in}{1.720198in}}%
\pgfpathlineto{\pgfqpoint{2.644890in}{1.715500in}}%
\pgfpathlineto{\pgfqpoint{2.657152in}{1.710803in}}%
\pgfpathlineto{\pgfqpoint{2.669413in}{1.706105in}}%
\pgfpathlineto{\pgfqpoint{2.681675in}{1.701407in}}%
\pgfpathlineto{\pgfqpoint{2.693937in}{1.696710in}}%
\pgfpathlineto{\pgfqpoint{2.706198in}{1.692012in}}%
\pgfpathlineto{\pgfqpoint{2.718460in}{1.687314in}}%
\pgfpathlineto{\pgfqpoint{2.730721in}{1.682617in}}%
\pgfpathlineto{\pgfqpoint{2.742983in}{1.677919in}}%
\pgfpathlineto{\pgfqpoint{2.755245in}{1.673222in}}%
\pgfpathlineto{\pgfqpoint{2.767506in}{1.668524in}}%
\pgfpathlineto{\pgfqpoint{2.779768in}{1.663826in}}%
\pgfpathlineto{\pgfqpoint{2.792029in}{1.659129in}}%
\pgfpathlineto{\pgfqpoint{2.804291in}{1.654431in}}%
\pgfpathlineto{\pgfqpoint{2.816553in}{1.649733in}}%
\pgfpathlineto{\pgfqpoint{2.828814in}{1.645036in}}%
\pgfpathlineto{\pgfqpoint{2.841076in}{1.640338in}}%
\pgfpathlineto{\pgfqpoint{2.853337in}{1.635641in}}%
\pgfpathlineto{\pgfqpoint{2.865599in}{1.630943in}}%
\pgfpathlineto{\pgfqpoint{2.877861in}{1.626245in}}%
\pgfpathlineto{\pgfqpoint{2.890122in}{1.621548in}}%
\pgfpathlineto{\pgfqpoint{2.902384in}{1.616850in}}%
\pgfpathlineto{\pgfqpoint{2.914646in}{1.612152in}}%
\pgfpathlineto{\pgfqpoint{2.926907in}{1.607455in}}%
\pgfpathlineto{\pgfqpoint{2.939169in}{1.602757in}}%
\pgfpathlineto{\pgfqpoint{2.951430in}{1.598060in}}%
\pgfpathlineto{\pgfqpoint{2.963692in}{1.593362in}}%
\pgfpathlineto{\pgfqpoint{2.975954in}{1.588664in}}%
\pgfpathlineto{\pgfqpoint{2.988215in}{1.583967in}}%
\pgfpathlineto{\pgfqpoint{3.000477in}{1.579269in}}%
\pgfpathlineto{\pgfqpoint{3.012738in}{1.574571in}}%
\pgfpathlineto{\pgfqpoint{3.025000in}{1.569874in}}%
\pgfpathlineto{\pgfqpoint{3.025000in}{3.588434in}}%
\pgfpathlineto{\pgfqpoint{3.025000in}{3.588434in}}%
\pgfpathlineto{\pgfqpoint{3.012738in}{3.577483in}}%
\pgfpathlineto{\pgfqpoint{3.000477in}{3.566532in}}%
\pgfpathlineto{\pgfqpoint{2.988215in}{3.555581in}}%
\pgfpathlineto{\pgfqpoint{2.975954in}{3.544631in}}%
\pgfpathlineto{\pgfqpoint{2.963692in}{3.533680in}}%
\pgfpathlineto{\pgfqpoint{2.951430in}{3.522729in}}%
\pgfpathlineto{\pgfqpoint{2.939169in}{3.511778in}}%
\pgfpathlineto{\pgfqpoint{2.926907in}{3.500827in}}%
\pgfpathlineto{\pgfqpoint{2.914646in}{3.489877in}}%
\pgfpathlineto{\pgfqpoint{2.902384in}{3.478926in}}%
\pgfpathlineto{\pgfqpoint{2.890122in}{3.467975in}}%
\pgfpathlineto{\pgfqpoint{2.877861in}{3.457024in}}%
\pgfpathlineto{\pgfqpoint{2.865599in}{3.446074in}}%
\pgfpathlineto{\pgfqpoint{2.853337in}{3.435123in}}%
\pgfpathlineto{\pgfqpoint{2.841076in}{3.424172in}}%
\pgfpathlineto{\pgfqpoint{2.828814in}{3.413221in}}%
\pgfpathlineto{\pgfqpoint{2.816553in}{3.402270in}}%
\pgfpathlineto{\pgfqpoint{2.804291in}{3.391320in}}%
\pgfpathlineto{\pgfqpoint{2.792029in}{3.380369in}}%
\pgfpathlineto{\pgfqpoint{2.779768in}{3.369418in}}%
\pgfpathlineto{\pgfqpoint{2.767506in}{3.358467in}}%
\pgfpathlineto{\pgfqpoint{2.755245in}{3.347517in}}%
\pgfpathlineto{\pgfqpoint{2.742983in}{3.336566in}}%
\pgfpathlineto{\pgfqpoint{2.730721in}{3.325615in}}%
\pgfpathlineto{\pgfqpoint{2.718460in}{3.314664in}}%
\pgfpathlineto{\pgfqpoint{2.706198in}{3.303714in}}%
\pgfpathlineto{\pgfqpoint{2.693937in}{3.292763in}}%
\pgfpathlineto{\pgfqpoint{2.681675in}{3.281812in}}%
\pgfpathlineto{\pgfqpoint{2.669413in}{3.270861in}}%
\pgfpathlineto{\pgfqpoint{2.657152in}{3.259910in}}%
\pgfpathlineto{\pgfqpoint{2.644890in}{3.248960in}}%
\pgfpathlineto{\pgfqpoint{2.632629in}{3.238009in}}%
\pgfpathlineto{\pgfqpoint{2.620367in}{3.227058in}}%
\pgfpathlineto{\pgfqpoint{2.608105in}{3.216107in}}%
\pgfpathlineto{\pgfqpoint{2.595844in}{3.205157in}}%
\pgfpathlineto{\pgfqpoint{2.583582in}{3.194206in}}%
\pgfpathlineto{\pgfqpoint{2.571321in}{3.183255in}}%
\pgfpathlineto{\pgfqpoint{2.559059in}{3.172304in}}%
\pgfpathlineto{\pgfqpoint{2.546797in}{3.161353in}}%
\pgfpathlineto{\pgfqpoint{2.534536in}{3.150403in}}%
\pgfpathlineto{\pgfqpoint{2.522274in}{3.139452in}}%
\pgfpathlineto{\pgfqpoint{2.510012in}{3.128501in}}%
\pgfpathlineto{\pgfqpoint{2.497751in}{3.117550in}}%
\pgfpathlineto{\pgfqpoint{2.485489in}{3.106600in}}%
\pgfpathlineto{\pgfqpoint{2.473228in}{3.095649in}}%
\pgfpathlineto{\pgfqpoint{2.460966in}{3.084698in}}%
\pgfpathlineto{\pgfqpoint{2.448704in}{3.073747in}}%
\pgfpathlineto{\pgfqpoint{2.436443in}{3.062796in}}%
\pgfpathlineto{\pgfqpoint{2.424181in}{3.051846in}}%
\pgfpathlineto{\pgfqpoint{2.411920in}{3.040895in}}%
\pgfpathlineto{\pgfqpoint{2.399658in}{3.029944in}}%
\pgfpathlineto{\pgfqpoint{2.387396in}{3.018993in}}%
\pgfpathlineto{\pgfqpoint{2.375135in}{3.008043in}}%
\pgfpathlineto{\pgfqpoint{2.362873in}{2.997092in}}%
\pgfpathlineto{\pgfqpoint{2.350612in}{2.986141in}}%
\pgfpathlineto{\pgfqpoint{2.338350in}{2.975190in}}%
\pgfpathlineto{\pgfqpoint{2.326088in}{2.964239in}}%
\pgfpathlineto{\pgfqpoint{2.313827in}{2.953289in}}%
\pgfpathlineto{\pgfqpoint{2.301565in}{2.942338in}}%
\pgfpathlineto{\pgfqpoint{2.289304in}{2.931387in}}%
\pgfpathlineto{\pgfqpoint{2.277042in}{2.920436in}}%
\pgfpathlineto{\pgfqpoint{2.264780in}{2.909486in}}%
\pgfpathlineto{\pgfqpoint{2.252519in}{2.898535in}}%
\pgfpathlineto{\pgfqpoint{2.240257in}{2.887584in}}%
\pgfpathlineto{\pgfqpoint{2.227995in}{2.876633in}}%
\pgfpathlineto{\pgfqpoint{2.215734in}{2.865682in}}%
\pgfpathlineto{\pgfqpoint{2.203472in}{2.854732in}}%
\pgfpathlineto{\pgfqpoint{2.191211in}{2.843781in}}%
\pgfpathlineto{\pgfqpoint{2.178949in}{2.832830in}}%
\pgfpathlineto{\pgfqpoint{2.166687in}{2.821879in}}%
\pgfpathlineto{\pgfqpoint{2.154426in}{2.810929in}}%
\pgfpathlineto{\pgfqpoint{2.142164in}{2.799978in}}%
\pgfpathlineto{\pgfqpoint{2.129903in}{2.789027in}}%
\pgfpathlineto{\pgfqpoint{2.117641in}{2.778076in}}%
\pgfpathlineto{\pgfqpoint{2.105379in}{2.767125in}}%
\pgfpathlineto{\pgfqpoint{2.093118in}{2.756175in}}%
\pgfpathlineto{\pgfqpoint{2.080856in}{2.745224in}}%
\pgfpathlineto{\pgfqpoint{2.068595in}{2.734273in}}%
\pgfpathlineto{\pgfqpoint{2.056333in}{2.723322in}}%
\pgfpathlineto{\pgfqpoint{2.044071in}{2.712372in}}%
\pgfpathlineto{\pgfqpoint{2.031810in}{2.701421in}}%
\pgfpathlineto{\pgfqpoint{2.019548in}{2.690470in}}%
\pgfpathlineto{\pgfqpoint{2.007287in}{2.679519in}}%
\pgfpathlineto{\pgfqpoint{1.995025in}{2.668569in}}%
\pgfpathlineto{\pgfqpoint{1.982763in}{2.657618in}}%
\pgfpathlineto{\pgfqpoint{1.970502in}{2.646667in}}%
\pgfpathlineto{\pgfqpoint{1.958240in}{2.635716in}}%
\pgfpathlineto{\pgfqpoint{1.945978in}{2.624765in}}%
\pgfpathlineto{\pgfqpoint{1.933717in}{2.613815in}}%
\pgfpathlineto{\pgfqpoint{1.921455in}{2.602864in}}%
\pgfpathlineto{\pgfqpoint{1.909194in}{2.591913in}}%
\pgfpathlineto{\pgfqpoint{1.896932in}{2.580962in}}%
\pgfpathlineto{\pgfqpoint{1.884670in}{2.570012in}}%
\pgfpathlineto{\pgfqpoint{1.872409in}{2.559061in}}%
\pgfpathlineto{\pgfqpoint{1.860147in}{2.548110in}}%
\pgfpathlineto{\pgfqpoint{1.847886in}{2.537159in}}%
\pgfpathlineto{\pgfqpoint{1.835624in}{2.526208in}}%
\pgfpathlineto{\pgfqpoint{1.823362in}{2.515258in}}%
\pgfpathlineto{\pgfqpoint{1.811101in}{2.504307in}}%
\pgfpathlineto{\pgfqpoint{1.798839in}{2.493356in}}%
\pgfpathlineto{\pgfqpoint{1.786578in}{2.482405in}}%
\pgfpathlineto{\pgfqpoint{1.774316in}{2.471455in}}%
\pgfpathlineto{\pgfqpoint{1.762054in}{2.460504in}}%
\pgfpathlineto{\pgfqpoint{1.749793in}{2.449553in}}%
\pgfpathlineto{\pgfqpoint{1.737531in}{2.438602in}}%
\pgfpathlineto{\pgfqpoint{1.725270in}{2.427651in}}%
\pgfpathlineto{\pgfqpoint{1.713008in}{2.416701in}}%
\pgfpathlineto{\pgfqpoint{1.700746in}{2.405750in}}%
\pgfpathlineto{\pgfqpoint{1.688485in}{2.394799in}}%
\pgfpathlineto{\pgfqpoint{1.676223in}{2.383848in}}%
\pgfpathlineto{\pgfqpoint{1.663962in}{2.372898in}}%
\pgfpathlineto{\pgfqpoint{1.651700in}{2.361947in}}%
\pgfpathlineto{\pgfqpoint{1.639438in}{2.350996in}}%
\pgfpathlineto{\pgfqpoint{1.627177in}{2.340045in}}%
\pgfpathlineto{\pgfqpoint{1.614915in}{2.329094in}}%
\pgfpathlineto{\pgfqpoint{1.602653in}{2.318144in}}%
\pgfpathlineto{\pgfqpoint{1.590392in}{2.307193in}}%
\pgfpathlineto{\pgfqpoint{1.578130in}{2.296242in}}%
\pgfpathlineto{\pgfqpoint{1.565869in}{2.285291in}}%
\pgfpathlineto{\pgfqpoint{1.553607in}{2.274341in}}%
\pgfpathlineto{\pgfqpoint{1.541345in}{2.263390in}}%
\pgfpathlineto{\pgfqpoint{1.529084in}{2.252439in}}%
\pgfpathlineto{\pgfqpoint{1.516822in}{2.241488in}}%
\pgfpathlineto{\pgfqpoint{1.504561in}{2.230537in}}%
\pgfpathlineto{\pgfqpoint{1.492299in}{2.219587in}}%
\pgfpathlineto{\pgfqpoint{1.480037in}{2.208636in}}%
\pgfpathlineto{\pgfqpoint{1.467776in}{2.197685in}}%
\pgfpathlineto{\pgfqpoint{1.455514in}{2.186734in}}%
\pgfpathlineto{\pgfqpoint{1.443253in}{2.175867in}}%
\pgfpathlineto{\pgfqpoint{1.430991in}{2.180565in}}%
\pgfpathlineto{\pgfqpoint{1.418729in}{2.185263in}}%
\pgfpathlineto{\pgfqpoint{1.406468in}{2.189960in}}%
\pgfpathlineto{\pgfqpoint{1.394206in}{2.194658in}}%
\pgfpathlineto{\pgfqpoint{1.381945in}{2.199356in}}%
\pgfpathlineto{\pgfqpoint{1.369683in}{2.204053in}}%
\pgfpathlineto{\pgfqpoint{1.357421in}{2.208751in}}%
\pgfpathlineto{\pgfqpoint{1.345160in}{2.213448in}}%
\pgfpathlineto{\pgfqpoint{1.332898in}{2.218146in}}%
\pgfpathlineto{\pgfqpoint{1.320636in}{2.222844in}}%
\pgfpathlineto{\pgfqpoint{1.308375in}{2.227541in}}%
\pgfpathlineto{\pgfqpoint{1.296113in}{2.232239in}}%
\pgfpathlineto{\pgfqpoint{1.283852in}{2.236937in}}%
\pgfpathlineto{\pgfqpoint{1.271590in}{2.241634in}}%
\pgfpathlineto{\pgfqpoint{1.259328in}{2.246332in}}%
\pgfpathlineto{\pgfqpoint{1.247067in}{2.251029in}}%
\pgfpathlineto{\pgfqpoint{1.234805in}{2.255727in}}%
\pgfpathlineto{\pgfqpoint{1.222544in}{2.260425in}}%
\pgfpathlineto{\pgfqpoint{1.210282in}{2.265122in}}%
\pgfpathlineto{\pgfqpoint{1.198020in}{2.269820in}}%
\pgfpathlineto{\pgfqpoint{1.185759in}{2.274518in}}%
\pgfpathlineto{\pgfqpoint{1.173497in}{2.279215in}}%
\pgfpathlineto{\pgfqpoint{1.161236in}{2.283913in}}%
\pgfpathlineto{\pgfqpoint{1.148974in}{2.288610in}}%
\pgfpathlineto{\pgfqpoint{1.136712in}{2.293308in}}%
\pgfpathlineto{\pgfqpoint{1.124451in}{2.298006in}}%
\pgfpathlineto{\pgfqpoint{1.112189in}{2.302703in}}%
\pgfpathlineto{\pgfqpoint{1.099928in}{2.307401in}}%
\pgfpathlineto{\pgfqpoint{1.087666in}{2.312099in}}%
\pgfpathlineto{\pgfqpoint{1.075404in}{2.316796in}}%
\pgfpathlineto{\pgfqpoint{1.063143in}{2.321494in}}%
\pgfpathlineto{\pgfqpoint{1.050881in}{2.326191in}}%
\pgfpathlineto{\pgfqpoint{1.038619in}{2.330889in}}%
\pgfpathlineto{\pgfqpoint{1.026358in}{2.335587in}}%
\pgfpathlineto{\pgfqpoint{1.014096in}{2.340284in}}%
\pgfpathlineto{\pgfqpoint{1.001835in}{2.344982in}}%
\pgfpathlineto{\pgfqpoint{0.989573in}{2.349680in}}%
\pgfpathlineto{\pgfqpoint{0.977311in}{2.354377in}}%
\pgfpathlineto{\pgfqpoint{0.965050in}{2.359075in}}%
\pgfpathlineto{\pgfqpoint{0.952788in}{2.363772in}}%
\pgfpathlineto{\pgfqpoint{0.940527in}{2.368470in}}%
\pgfpathlineto{\pgfqpoint{0.928265in}{2.373168in}}%
\pgfpathlineto{\pgfqpoint{0.916003in}{2.377865in}}%
\pgfpathlineto{\pgfqpoint{0.903742in}{2.382563in}}%
\pgfpathlineto{\pgfqpoint{0.891480in}{2.387261in}}%
\pgfpathlineto{\pgfqpoint{0.879219in}{2.391958in}}%
\pgfpathlineto{\pgfqpoint{0.866957in}{2.396656in}}%
\pgfpathlineto{\pgfqpoint{0.854695in}{2.401353in}}%
\pgfpathlineto{\pgfqpoint{0.842434in}{2.406051in}}%
\pgfpathlineto{\pgfqpoint{0.830172in}{2.410749in}}%
\pgfpathlineto{\pgfqpoint{0.817911in}{2.415446in}}%
\pgfpathlineto{\pgfqpoint{0.805649in}{2.420144in}}%
\pgfpathlineto{\pgfqpoint{0.793387in}{2.424842in}}%
\pgfpathlineto{\pgfqpoint{0.781126in}{2.429539in}}%
\pgfpathlineto{\pgfqpoint{0.768864in}{2.434237in}}%
\pgfpathlineto{\pgfqpoint{0.756603in}{2.438934in}}%
\pgfpathlineto{\pgfqpoint{0.744341in}{2.443632in}}%
\pgfpathlineto{\pgfqpoint{0.732079in}{2.448330in}}%
\pgfpathlineto{\pgfqpoint{0.719818in}{2.453027in}}%
\pgfpathlineto{\pgfqpoint{0.707556in}{2.457725in}}%
\pgfpathlineto{\pgfqpoint{0.695294in}{2.462423in}}%
\pgfpathlineto{\pgfqpoint{0.683033in}{2.467120in}}%
\pgfpathlineto{\pgfqpoint{0.670771in}{2.471818in}}%
\pgfpathlineto{\pgfqpoint{0.658510in}{2.476515in}}%
\pgfpathlineto{\pgfqpoint{0.646248in}{2.481213in}}%
\pgfpathlineto{\pgfqpoint{0.633986in}{2.485911in}}%
\pgfpathlineto{\pgfqpoint{0.621725in}{2.490608in}}%
\pgfpathlineto{\pgfqpoint{0.609463in}{2.495306in}}%
\pgfpathlineto{\pgfqpoint{0.597202in}{2.500004in}}%
\pgfpathlineto{\pgfqpoint{0.584940in}{2.504701in}}%
\pgfpathlineto{\pgfqpoint{0.584940in}{2.504701in}}%
\pgfpathclose%
\pgfusepath{stroke,fill}%
}%
\begin{pgfscope}%
\pgfsys@transformshift{0.000000in}{0.000000in}%
\pgfsys@useobject{currentmarker}{}%
\end{pgfscope}%
\end{pgfscope}%
\begin{pgfscope}%
\pgfsetbuttcap%
\pgfsetroundjoin%
\definecolor{currentfill}{rgb}{0.000000,0.000000,0.000000}%
\pgfsetfillcolor{currentfill}%
\pgfsetlinewidth{0.803000pt}%
\definecolor{currentstroke}{rgb}{0.000000,0.000000,0.000000}%
\pgfsetstrokecolor{currentstroke}%
\pgfsetdash{}{0pt}%
\pgfsys@defobject{currentmarker}{\pgfqpoint{0.000000in}{-0.048611in}}{\pgfqpoint{0.000000in}{0.000000in}}{%
\pgfpathmoveto{\pgfqpoint{0.000000in}{0.000000in}}%
\pgfpathlineto{\pgfqpoint{0.000000in}{-0.048611in}}%
\pgfusepath{stroke,fill}%
}%
\begin{pgfscope}%
\pgfsys@transformshift{1.151382in}{1.184369in}%
\pgfsys@useobject{currentmarker}{}%
\end{pgfscope}%
\end{pgfscope}%
\begin{pgfscope}%
\definecolor{textcolor}{rgb}{0.000000,0.000000,0.000000}%
\pgfsetstrokecolor{textcolor}%
\pgfsetfillcolor{textcolor}%
\pgftext[x=1.151382in,y=1.087147in,,top]{\color{textcolor}{\rmfamily\fontsize{10.000000}{12.000000}\selectfont\catcode`\^=\active\def^{\ifmmode\sp\else\^{}\fi}\catcode`\%=\active\def%{\%}$\mathdefault{1.4}$}}%
\end{pgfscope}%
\begin{pgfscope}%
\pgfsetbuttcap%
\pgfsetroundjoin%
\definecolor{currentfill}{rgb}{0.000000,0.000000,0.000000}%
\pgfsetfillcolor{currentfill}%
\pgfsetlinewidth{0.803000pt}%
\definecolor{currentstroke}{rgb}{0.000000,0.000000,0.000000}%
\pgfsetstrokecolor{currentstroke}%
\pgfsetdash{}{0pt}%
\pgfsys@defobject{currentmarker}{\pgfqpoint{0.000000in}{-0.048611in}}{\pgfqpoint{0.000000in}{0.000000in}}{%
\pgfpathmoveto{\pgfqpoint{0.000000in}{0.000000in}}%
\pgfpathlineto{\pgfqpoint{0.000000in}{-0.048611in}}%
\pgfusepath{stroke,fill}%
}%
\begin{pgfscope}%
\pgfsys@transformshift{2.022832in}{1.184369in}%
\pgfsys@useobject{currentmarker}{}%
\end{pgfscope}%
\end{pgfscope}%
\begin{pgfscope}%
\definecolor{textcolor}{rgb}{0.000000,0.000000,0.000000}%
\pgfsetstrokecolor{textcolor}%
\pgfsetfillcolor{textcolor}%
\pgftext[x=2.022832in,y=1.087147in,,top]{\color{textcolor}{\rmfamily\fontsize{10.000000}{12.000000}\selectfont\catcode`\^=\active\def^{\ifmmode\sp\else\^{}\fi}\catcode`\%=\active\def%{\%}$\mathdefault{1.6}$}}%
\end{pgfscope}%
\begin{pgfscope}%
\pgfsetbuttcap%
\pgfsetroundjoin%
\definecolor{currentfill}{rgb}{0.000000,0.000000,0.000000}%
\pgfsetfillcolor{currentfill}%
\pgfsetlinewidth{0.803000pt}%
\definecolor{currentstroke}{rgb}{0.000000,0.000000,0.000000}%
\pgfsetstrokecolor{currentstroke}%
\pgfsetdash{}{0pt}%
\pgfsys@defobject{currentmarker}{\pgfqpoint{0.000000in}{-0.048611in}}{\pgfqpoint{0.000000in}{0.000000in}}{%
\pgfpathmoveto{\pgfqpoint{0.000000in}{0.000000in}}%
\pgfpathlineto{\pgfqpoint{0.000000in}{-0.048611in}}%
\pgfusepath{stroke,fill}%
}%
\begin{pgfscope}%
\pgfsys@transformshift{2.894282in}{1.184369in}%
\pgfsys@useobject{currentmarker}{}%
\end{pgfscope}%
\end{pgfscope}%
\begin{pgfscope}%
\definecolor{textcolor}{rgb}{0.000000,0.000000,0.000000}%
\pgfsetstrokecolor{textcolor}%
\pgfsetfillcolor{textcolor}%
\pgftext[x=2.894282in,y=1.087147in,,top]{\color{textcolor}{\rmfamily\fontsize{10.000000}{12.000000}\selectfont\catcode`\^=\active\def^{\ifmmode\sp\else\^{}\fi}\catcode`\%=\active\def%{\%}$\mathdefault{1.8}$}}%
\end{pgfscope}%
\begin{pgfscope}%
\definecolor{textcolor}{rgb}{0.000000,0.000000,0.000000}%
\pgfsetstrokecolor{textcolor}%
\pgfsetfillcolor{textcolor}%
\pgftext[x=1.804970in,y=0.891940in,,top]{\color{textcolor}{\rmfamily\fontsize{10.000000}{12.000000}\selectfont\catcode`\^=\active\def^{\ifmmode\sp\else\^{}\fi}\catcode`\%=\active\def%{\%}$C_{DA}$ (m$^2$)}}%
\end{pgfscope}%
\begin{pgfscope}%
\pgfsetbuttcap%
\pgfsetroundjoin%
\definecolor{currentfill}{rgb}{0.000000,0.000000,0.000000}%
\pgfsetfillcolor{currentfill}%
\pgfsetlinewidth{0.803000pt}%
\definecolor{currentstroke}{rgb}{0.000000,0.000000,0.000000}%
\pgfsetstrokecolor{currentstroke}%
\pgfsetdash{}{0pt}%
\pgfsys@defobject{currentmarker}{\pgfqpoint{-0.048611in}{0.000000in}}{\pgfqpoint{-0.000000in}{0.000000in}}{%
\pgfpathmoveto{\pgfqpoint{-0.000000in}{0.000000in}}%
\pgfpathlineto{\pgfqpoint{-0.048611in}{0.000000in}}%
\pgfusepath{stroke,fill}%
}%
\begin{pgfscope}%
\pgfsys@transformshift{0.584940in}{1.252992in}%
\pgfsys@useobject{currentmarker}{}%
\end{pgfscope}%
\end{pgfscope}%
\begin{pgfscope}%
\definecolor{textcolor}{rgb}{0.000000,0.000000,0.000000}%
\pgfsetstrokecolor{textcolor}%
\pgfsetfillcolor{textcolor}%
\pgftext[x=0.314107in, y=1.202750in, left, base]{\color{textcolor}{\rmfamily\fontsize{10.000000}{12.000000}\selectfont\catcode`\^=\active\def^{\ifmmode\sp\else\^{}\fi}\catcode`\%=\active\def%{\%}$\mathdefault{2.0}$}}%
\end{pgfscope}%
\begin{pgfscope}%
\pgfsetbuttcap%
\pgfsetroundjoin%
\definecolor{currentfill}{rgb}{0.000000,0.000000,0.000000}%
\pgfsetfillcolor{currentfill}%
\pgfsetlinewidth{0.803000pt}%
\definecolor{currentstroke}{rgb}{0.000000,0.000000,0.000000}%
\pgfsetstrokecolor{currentstroke}%
\pgfsetdash{}{0pt}%
\pgfsys@defobject{currentmarker}{\pgfqpoint{-0.048611in}{0.000000in}}{\pgfqpoint{-0.000000in}{0.000000in}}{%
\pgfpathmoveto{\pgfqpoint{-0.000000in}{0.000000in}}%
\pgfpathlineto{\pgfqpoint{-0.048611in}{0.000000in}}%
\pgfusepath{stroke,fill}%
}%
\begin{pgfscope}%
\pgfsys@transformshift{0.584940in}{1.596110in}%
\pgfsys@useobject{currentmarker}{}%
\end{pgfscope}%
\end{pgfscope}%
\begin{pgfscope}%
\definecolor{textcolor}{rgb}{0.000000,0.000000,0.000000}%
\pgfsetstrokecolor{textcolor}%
\pgfsetfillcolor{textcolor}%
\pgftext[x=0.314107in, y=1.545868in, left, base]{\color{textcolor}{\rmfamily\fontsize{10.000000}{12.000000}\selectfont\catcode`\^=\active\def^{\ifmmode\sp\else\^{}\fi}\catcode`\%=\active\def%{\%}$\mathdefault{2.5}$}}%
\end{pgfscope}%
\begin{pgfscope}%
\pgfsetbuttcap%
\pgfsetroundjoin%
\definecolor{currentfill}{rgb}{0.000000,0.000000,0.000000}%
\pgfsetfillcolor{currentfill}%
\pgfsetlinewidth{0.803000pt}%
\definecolor{currentstroke}{rgb}{0.000000,0.000000,0.000000}%
\pgfsetstrokecolor{currentstroke}%
\pgfsetdash{}{0pt}%
\pgfsys@defobject{currentmarker}{\pgfqpoint{-0.048611in}{0.000000in}}{\pgfqpoint{-0.000000in}{0.000000in}}{%
\pgfpathmoveto{\pgfqpoint{-0.000000in}{0.000000in}}%
\pgfpathlineto{\pgfqpoint{-0.048611in}{0.000000in}}%
\pgfusepath{stroke,fill}%
}%
\begin{pgfscope}%
\pgfsys@transformshift{0.584940in}{1.939228in}%
\pgfsys@useobject{currentmarker}{}%
\end{pgfscope}%
\end{pgfscope}%
\begin{pgfscope}%
\definecolor{textcolor}{rgb}{0.000000,0.000000,0.000000}%
\pgfsetstrokecolor{textcolor}%
\pgfsetfillcolor{textcolor}%
\pgftext[x=0.314107in, y=1.888986in, left, base]{\color{textcolor}{\rmfamily\fontsize{10.000000}{12.000000}\selectfont\catcode`\^=\active\def^{\ifmmode\sp\else\^{}\fi}\catcode`\%=\active\def%{\%}$\mathdefault{3.0}$}}%
\end{pgfscope}%
\begin{pgfscope}%
\pgfsetbuttcap%
\pgfsetroundjoin%
\definecolor{currentfill}{rgb}{0.000000,0.000000,0.000000}%
\pgfsetfillcolor{currentfill}%
\pgfsetlinewidth{0.803000pt}%
\definecolor{currentstroke}{rgb}{0.000000,0.000000,0.000000}%
\pgfsetstrokecolor{currentstroke}%
\pgfsetdash{}{0pt}%
\pgfsys@defobject{currentmarker}{\pgfqpoint{-0.048611in}{0.000000in}}{\pgfqpoint{-0.000000in}{0.000000in}}{%
\pgfpathmoveto{\pgfqpoint{-0.000000in}{0.000000in}}%
\pgfpathlineto{\pgfqpoint{-0.048611in}{0.000000in}}%
\pgfusepath{stroke,fill}%
}%
\begin{pgfscope}%
\pgfsys@transformshift{0.584940in}{2.282346in}%
\pgfsys@useobject{currentmarker}{}%
\end{pgfscope}%
\end{pgfscope}%
\begin{pgfscope}%
\definecolor{textcolor}{rgb}{0.000000,0.000000,0.000000}%
\pgfsetstrokecolor{textcolor}%
\pgfsetfillcolor{textcolor}%
\pgftext[x=0.314107in, y=2.232104in, left, base]{\color{textcolor}{\rmfamily\fontsize{10.000000}{12.000000}\selectfont\catcode`\^=\active\def^{\ifmmode\sp\else\^{}\fi}\catcode`\%=\active\def%{\%}$\mathdefault{3.5}$}}%
\end{pgfscope}%
\begin{pgfscope}%
\pgfsetbuttcap%
\pgfsetroundjoin%
\definecolor{currentfill}{rgb}{0.000000,0.000000,0.000000}%
\pgfsetfillcolor{currentfill}%
\pgfsetlinewidth{0.803000pt}%
\definecolor{currentstroke}{rgb}{0.000000,0.000000,0.000000}%
\pgfsetstrokecolor{currentstroke}%
\pgfsetdash{}{0pt}%
\pgfsys@defobject{currentmarker}{\pgfqpoint{-0.048611in}{0.000000in}}{\pgfqpoint{-0.000000in}{0.000000in}}{%
\pgfpathmoveto{\pgfqpoint{-0.000000in}{0.000000in}}%
\pgfpathlineto{\pgfqpoint{-0.048611in}{0.000000in}}%
\pgfusepath{stroke,fill}%
}%
\begin{pgfscope}%
\pgfsys@transformshift{0.584940in}{2.625464in}%
\pgfsys@useobject{currentmarker}{}%
\end{pgfscope}%
\end{pgfscope}%
\begin{pgfscope}%
\definecolor{textcolor}{rgb}{0.000000,0.000000,0.000000}%
\pgfsetstrokecolor{textcolor}%
\pgfsetfillcolor{textcolor}%
\pgftext[x=0.314107in, y=2.575221in, left, base]{\color{textcolor}{\rmfamily\fontsize{10.000000}{12.000000}\selectfont\catcode`\^=\active\def^{\ifmmode\sp\else\^{}\fi}\catcode`\%=\active\def%{\%}$\mathdefault{4.0}$}}%
\end{pgfscope}%
\begin{pgfscope}%
\pgfsetbuttcap%
\pgfsetroundjoin%
\definecolor{currentfill}{rgb}{0.000000,0.000000,0.000000}%
\pgfsetfillcolor{currentfill}%
\pgfsetlinewidth{0.803000pt}%
\definecolor{currentstroke}{rgb}{0.000000,0.000000,0.000000}%
\pgfsetstrokecolor{currentstroke}%
\pgfsetdash{}{0pt}%
\pgfsys@defobject{currentmarker}{\pgfqpoint{-0.048611in}{0.000000in}}{\pgfqpoint{-0.000000in}{0.000000in}}{%
\pgfpathmoveto{\pgfqpoint{-0.000000in}{0.000000in}}%
\pgfpathlineto{\pgfqpoint{-0.048611in}{0.000000in}}%
\pgfusepath{stroke,fill}%
}%
\begin{pgfscope}%
\pgfsys@transformshift{0.584940in}{2.968582in}%
\pgfsys@useobject{currentmarker}{}%
\end{pgfscope}%
\end{pgfscope}%
\begin{pgfscope}%
\definecolor{textcolor}{rgb}{0.000000,0.000000,0.000000}%
\pgfsetstrokecolor{textcolor}%
\pgfsetfillcolor{textcolor}%
\pgftext[x=0.314107in, y=2.918339in, left, base]{\color{textcolor}{\rmfamily\fontsize{10.000000}{12.000000}\selectfont\catcode`\^=\active\def^{\ifmmode\sp\else\^{}\fi}\catcode`\%=\active\def%{\%}$\mathdefault{4.5}$}}%
\end{pgfscope}%
\begin{pgfscope}%
\pgfsetbuttcap%
\pgfsetroundjoin%
\definecolor{currentfill}{rgb}{0.000000,0.000000,0.000000}%
\pgfsetfillcolor{currentfill}%
\pgfsetlinewidth{0.803000pt}%
\definecolor{currentstroke}{rgb}{0.000000,0.000000,0.000000}%
\pgfsetstrokecolor{currentstroke}%
\pgfsetdash{}{0pt}%
\pgfsys@defobject{currentmarker}{\pgfqpoint{-0.048611in}{0.000000in}}{\pgfqpoint{-0.000000in}{0.000000in}}{%
\pgfpathmoveto{\pgfqpoint{-0.000000in}{0.000000in}}%
\pgfpathlineto{\pgfqpoint{-0.048611in}{0.000000in}}%
\pgfusepath{stroke,fill}%
}%
\begin{pgfscope}%
\pgfsys@transformshift{0.584940in}{3.311700in}%
\pgfsys@useobject{currentmarker}{}%
\end{pgfscope}%
\end{pgfscope}%
\begin{pgfscope}%
\definecolor{textcolor}{rgb}{0.000000,0.000000,0.000000}%
\pgfsetstrokecolor{textcolor}%
\pgfsetfillcolor{textcolor}%
\pgftext[x=0.314107in, y=3.261457in, left, base]{\color{textcolor}{\rmfamily\fontsize{10.000000}{12.000000}\selectfont\catcode`\^=\active\def^{\ifmmode\sp\else\^{}\fi}\catcode`\%=\active\def%{\%}$\mathdefault{5.0}$}}%
\end{pgfscope}%
\begin{pgfscope}%
\definecolor{textcolor}{rgb}{0.000000,0.000000,0.000000}%
\pgfsetstrokecolor{textcolor}%
\pgfsetfillcolor{textcolor}%
\pgftext[x=0.258551in,y=2.350970in,,bottom,rotate=90.000000]{\color{textcolor}{\rmfamily\fontsize{10.000000}{12.000000}\selectfont\catcode`\^=\active\def^{\ifmmode\sp\else\^{}\fi}\catcode`\%=\active\def%{\%}$C_{LA}$ (m$^2$)}}%
\end{pgfscope}%
\begin{pgfscope}%
\pgfpathrectangle{\pgfqpoint{0.584940in}{1.184369in}}{\pgfqpoint{2.440060in}{2.333202in}}%
\pgfusepath{clip}%
\pgfsetbuttcap%
\pgfsetroundjoin%
\pgfsetlinewidth{0.803000pt}%
\definecolor{currentstroke}{rgb}{0.501961,0.501961,0.501961}%
\pgfsetstrokecolor{currentstroke}%
\pgfsetstrokeopacity{0.600000}%
\pgfsetdash{{2.960000pt}{1.280000pt}}{0.000000pt}%
\pgfpathmoveto{\pgfqpoint{0.584940in}{1.409229in}}%
\pgfpathlineto{\pgfqpoint{1.443253in}{2.175784in}}%
\pgfpathlineto{\pgfqpoint{1.627177in}{2.105403in}}%
\pgfpathlineto{\pgfqpoint{3.025000in}{1.569874in}}%
\pgfpathlineto{\pgfqpoint{3.025000in}{1.569874in}}%
\pgfusepath{stroke}%
\end{pgfscope}%
\begin{pgfscope}%
\pgfpathrectangle{\pgfqpoint{0.584940in}{1.184369in}}{\pgfqpoint{2.440060in}{2.333202in}}%
\pgfusepath{clip}%
\pgfsetbuttcap%
\pgfsetroundjoin%
\pgfsetlinewidth{0.803000pt}%
\definecolor{currentstroke}{rgb}{0.501961,0.501961,0.501961}%
\pgfsetstrokecolor{currentstroke}%
\pgfsetstrokeopacity{0.600000}%
\pgfsetdash{{2.960000pt}{1.280000pt}}{0.000000pt}%
\pgfpathmoveto{\pgfqpoint{0.584940in}{2.504701in}}%
\pgfpathlineto{\pgfqpoint{1.443253in}{2.175867in}}%
\pgfpathlineto{\pgfqpoint{1.663962in}{2.372898in}}%
\pgfpathlineto{\pgfqpoint{2.956851in}{3.527571in}}%
\pgfpathlineto{\pgfqpoint{2.956851in}{3.527571in}}%
\pgfusepath{stroke}%
\end{pgfscope}%
\begin{pgfscope}%
\pgfpathrectangle{\pgfqpoint{0.584940in}{1.184369in}}{\pgfqpoint{2.440060in}{2.333202in}}%
\pgfusepath{clip}%
\pgfsetrectcap%
\pgfsetroundjoin%
\pgfsetlinewidth{1.505625pt}%
\definecolor{currentstroke}{rgb}{0.000000,0.000000,0.000000}%
\pgfsetstrokecolor{currentstroke}%
\pgfsetdash{}{0pt}%
\pgfpathmoveto{\pgfqpoint{0.584940in}{1.768035in}}%
\pgfpathlineto{\pgfqpoint{3.025000in}{2.927283in}}%
\pgfpathlineto{\pgfqpoint{3.025000in}{2.927283in}}%
\pgfusepath{stroke}%
\end{pgfscope}%
\begin{pgfscope}%
\pgfpathrectangle{\pgfqpoint{0.584940in}{1.184369in}}{\pgfqpoint{2.440060in}{2.333202in}}%
\pgfusepath{clip}%
\pgfsetbuttcap%
\pgfsetroundjoin%
\pgfsetlinewidth{1.505625pt}%
\definecolor{currentstroke}{rgb}{0.129412,0.400000,0.674510}%
\pgfsetstrokecolor{currentstroke}%
\pgfsetdash{}{0pt}%
\pgfpathmoveto{\pgfqpoint{0.939397in}{1.707280in}}%
\pgfpathlineto{\pgfqpoint{1.781664in}{1.707280in}}%
\pgfusepath{stroke}%
\end{pgfscope}%
\begin{pgfscope}%
\pgfpathrectangle{\pgfqpoint{0.584940in}{1.184369in}}{\pgfqpoint{2.440060in}{2.333202in}}%
\pgfusepath{clip}%
\pgfsetbuttcap%
\pgfsetroundjoin%
\pgfsetlinewidth{1.505625pt}%
\definecolor{currentstroke}{rgb}{0.129412,0.400000,0.674510}%
\pgfsetstrokecolor{currentstroke}%
\pgfsetdash{}{0pt}%
\pgfpathmoveto{\pgfqpoint{1.360530in}{1.418248in}}%
\pgfpathlineto{\pgfqpoint{1.360530in}{1.996312in}}%
\pgfusepath{stroke}%
\end{pgfscope}%
\begin{pgfscope}%
\pgfpathrectangle{\pgfqpoint{0.584940in}{1.184369in}}{\pgfqpoint{2.440060in}{2.333202in}}%
\pgfusepath{clip}%
\pgfsetbuttcap%
\pgfsetroundjoin%
\definecolor{currentfill}{rgb}{0.129412,0.400000,0.674510}%
\pgfsetfillcolor{currentfill}%
\pgfsetlinewidth{1.003750pt}%
\definecolor{currentstroke}{rgb}{0.129412,0.400000,0.674510}%
\pgfsetstrokecolor{currentstroke}%
\pgfsetdash{}{0pt}%
\pgfsys@defobject{currentmarker}{\pgfqpoint{0.000000in}{-0.069444in}}{\pgfqpoint{0.000000in}{0.069444in}}{%
\pgfpathmoveto{\pgfqpoint{0.000000in}{-0.069444in}}%
\pgfpathlineto{\pgfqpoint{0.000000in}{0.069444in}}%
\pgfusepath{stroke,fill}%
}%
\begin{pgfscope}%
\pgfsys@transformshift{0.939397in}{1.707280in}%
\pgfsys@useobject{currentmarker}{}%
\end{pgfscope}%
\end{pgfscope}%
\begin{pgfscope}%
\pgfpathrectangle{\pgfqpoint{0.584940in}{1.184369in}}{\pgfqpoint{2.440060in}{2.333202in}}%
\pgfusepath{clip}%
\pgfsetbuttcap%
\pgfsetroundjoin%
\definecolor{currentfill}{rgb}{0.129412,0.400000,0.674510}%
\pgfsetfillcolor{currentfill}%
\pgfsetlinewidth{1.003750pt}%
\definecolor{currentstroke}{rgb}{0.129412,0.400000,0.674510}%
\pgfsetstrokecolor{currentstroke}%
\pgfsetdash{}{0pt}%
\pgfsys@defobject{currentmarker}{\pgfqpoint{0.000000in}{-0.069444in}}{\pgfqpoint{0.000000in}{0.069444in}}{%
\pgfpathmoveto{\pgfqpoint{0.000000in}{-0.069444in}}%
\pgfpathlineto{\pgfqpoint{0.000000in}{0.069444in}}%
\pgfusepath{stroke,fill}%
}%
\begin{pgfscope}%
\pgfsys@transformshift{1.781664in}{1.707280in}%
\pgfsys@useobject{currentmarker}{}%
\end{pgfscope}%
\end{pgfscope}%
\begin{pgfscope}%
\pgfpathrectangle{\pgfqpoint{0.584940in}{1.184369in}}{\pgfqpoint{2.440060in}{2.333202in}}%
\pgfusepath{clip}%
\pgfsetbuttcap%
\pgfsetroundjoin%
\definecolor{currentfill}{rgb}{0.129412,0.400000,0.674510}%
\pgfsetfillcolor{currentfill}%
\pgfsetlinewidth{1.003750pt}%
\definecolor{currentstroke}{rgb}{0.129412,0.400000,0.674510}%
\pgfsetstrokecolor{currentstroke}%
\pgfsetdash{}{0pt}%
\pgfsys@defobject{currentmarker}{\pgfqpoint{-0.069444in}{-0.000000in}}{\pgfqpoint{0.069444in}{0.000000in}}{%
\pgfpathmoveto{\pgfqpoint{0.069444in}{-0.000000in}}%
\pgfpathlineto{\pgfqpoint{-0.069444in}{0.000000in}}%
\pgfusepath{stroke,fill}%
}%
\begin{pgfscope}%
\pgfsys@transformshift{1.360530in}{1.418248in}%
\pgfsys@useobject{currentmarker}{}%
\end{pgfscope}%
\end{pgfscope}%
\begin{pgfscope}%
\pgfpathrectangle{\pgfqpoint{0.584940in}{1.184369in}}{\pgfqpoint{2.440060in}{2.333202in}}%
\pgfusepath{clip}%
\pgfsetbuttcap%
\pgfsetroundjoin%
\definecolor{currentfill}{rgb}{0.129412,0.400000,0.674510}%
\pgfsetfillcolor{currentfill}%
\pgfsetlinewidth{1.003750pt}%
\definecolor{currentstroke}{rgb}{0.129412,0.400000,0.674510}%
\pgfsetstrokecolor{currentstroke}%
\pgfsetdash{}{0pt}%
\pgfsys@defobject{currentmarker}{\pgfqpoint{-0.069444in}{-0.000000in}}{\pgfqpoint{0.069444in}{0.000000in}}{%
\pgfpathmoveto{\pgfqpoint{0.069444in}{-0.000000in}}%
\pgfpathlineto{\pgfqpoint{-0.069444in}{0.000000in}}%
\pgfusepath{stroke,fill}%
}%
\begin{pgfscope}%
\pgfsys@transformshift{1.360530in}{1.996312in}%
\pgfsys@useobject{currentmarker}{}%
\end{pgfscope}%
\end{pgfscope}%
\begin{pgfscope}%
\pgfpathrectangle{\pgfqpoint{0.584940in}{1.184369in}}{\pgfqpoint{2.440060in}{2.333202in}}%
\pgfusepath{clip}%
\pgfsetbuttcap%
\pgfsetroundjoin%
\pgfsetlinewidth{1.505625pt}%
\definecolor{currentstroke}{rgb}{0.454902,0.678431,0.819608}%
\pgfsetstrokecolor{currentstroke}%
\pgfsetdash{}{0pt}%
\pgfpathmoveto{\pgfqpoint{0.574940in}{1.983147in}}%
\pgfpathlineto{\pgfqpoint{1.770098in}{1.983147in}}%
\pgfusepath{stroke}%
\end{pgfscope}%
\begin{pgfscope}%
\pgfpathrectangle{\pgfqpoint{0.584940in}{1.184369in}}{\pgfqpoint{2.440060in}{2.333202in}}%
\pgfusepath{clip}%
\pgfsetbuttcap%
\pgfsetroundjoin%
\pgfsetlinewidth{1.505625pt}%
\definecolor{currentstroke}{rgb}{0.454902,0.678431,0.819608}%
\pgfsetstrokecolor{currentstroke}%
\pgfsetdash{}{0pt}%
\pgfpathmoveto{\pgfqpoint{1.046808in}{1.655565in}}%
\pgfpathlineto{\pgfqpoint{1.046808in}{2.310729in}}%
\pgfusepath{stroke}%
\end{pgfscope}%
\begin{pgfscope}%
\pgfpathrectangle{\pgfqpoint{0.584940in}{1.184369in}}{\pgfqpoint{2.440060in}{2.333202in}}%
\pgfusepath{clip}%
\pgfsetbuttcap%
\pgfsetroundjoin%
\definecolor{currentfill}{rgb}{0.454902,0.678431,0.819608}%
\pgfsetfillcolor{currentfill}%
\pgfsetlinewidth{1.003750pt}%
\definecolor{currentstroke}{rgb}{0.454902,0.678431,0.819608}%
\pgfsetstrokecolor{currentstroke}%
\pgfsetdash{}{0pt}%
\pgfsys@defobject{currentmarker}{\pgfqpoint{0.000000in}{-0.069444in}}{\pgfqpoint{0.000000in}{0.069444in}}{%
\pgfpathmoveto{\pgfqpoint{0.000000in}{-0.069444in}}%
\pgfpathlineto{\pgfqpoint{0.000000in}{0.069444in}}%
\pgfusepath{stroke,fill}%
}%
\begin{pgfscope}%
\pgfsys@transformshift{0.323519in}{1.983147in}%
\pgfsys@useobject{currentmarker}{}%
\end{pgfscope}%
\end{pgfscope}%
\begin{pgfscope}%
\pgfpathrectangle{\pgfqpoint{0.584940in}{1.184369in}}{\pgfqpoint{2.440060in}{2.333202in}}%
\pgfusepath{clip}%
\pgfsetbuttcap%
\pgfsetroundjoin%
\definecolor{currentfill}{rgb}{0.454902,0.678431,0.819608}%
\pgfsetfillcolor{currentfill}%
\pgfsetlinewidth{1.003750pt}%
\definecolor{currentstroke}{rgb}{0.454902,0.678431,0.819608}%
\pgfsetstrokecolor{currentstroke}%
\pgfsetdash{}{0pt}%
\pgfsys@defobject{currentmarker}{\pgfqpoint{0.000000in}{-0.069444in}}{\pgfqpoint{0.000000in}{0.069444in}}{%
\pgfpathmoveto{\pgfqpoint{0.000000in}{-0.069444in}}%
\pgfpathlineto{\pgfqpoint{0.000000in}{0.069444in}}%
\pgfusepath{stroke,fill}%
}%
\begin{pgfscope}%
\pgfsys@transformshift{1.770098in}{1.983147in}%
\pgfsys@useobject{currentmarker}{}%
\end{pgfscope}%
\end{pgfscope}%
\begin{pgfscope}%
\pgfpathrectangle{\pgfqpoint{0.584940in}{1.184369in}}{\pgfqpoint{2.440060in}{2.333202in}}%
\pgfusepath{clip}%
\pgfsetbuttcap%
\pgfsetroundjoin%
\definecolor{currentfill}{rgb}{0.454902,0.678431,0.819608}%
\pgfsetfillcolor{currentfill}%
\pgfsetlinewidth{1.003750pt}%
\definecolor{currentstroke}{rgb}{0.454902,0.678431,0.819608}%
\pgfsetstrokecolor{currentstroke}%
\pgfsetdash{}{0pt}%
\pgfsys@defobject{currentmarker}{\pgfqpoint{-0.069444in}{-0.000000in}}{\pgfqpoint{0.069444in}{0.000000in}}{%
\pgfpathmoveto{\pgfqpoint{0.069444in}{-0.000000in}}%
\pgfpathlineto{\pgfqpoint{-0.069444in}{0.000000in}}%
\pgfusepath{stroke,fill}%
}%
\begin{pgfscope}%
\pgfsys@transformshift{1.046808in}{1.655565in}%
\pgfsys@useobject{currentmarker}{}%
\end{pgfscope}%
\end{pgfscope}%
\begin{pgfscope}%
\pgfpathrectangle{\pgfqpoint{0.584940in}{1.184369in}}{\pgfqpoint{2.440060in}{2.333202in}}%
\pgfusepath{clip}%
\pgfsetbuttcap%
\pgfsetroundjoin%
\definecolor{currentfill}{rgb}{0.454902,0.678431,0.819608}%
\pgfsetfillcolor{currentfill}%
\pgfsetlinewidth{1.003750pt}%
\definecolor{currentstroke}{rgb}{0.454902,0.678431,0.819608}%
\pgfsetstrokecolor{currentstroke}%
\pgfsetdash{}{0pt}%
\pgfsys@defobject{currentmarker}{\pgfqpoint{-0.069444in}{-0.000000in}}{\pgfqpoint{0.069444in}{0.000000in}}{%
\pgfpathmoveto{\pgfqpoint{0.069444in}{-0.000000in}}%
\pgfpathlineto{\pgfqpoint{-0.069444in}{0.000000in}}%
\pgfusepath{stroke,fill}%
}%
\begin{pgfscope}%
\pgfsys@transformshift{1.046808in}{2.310729in}%
\pgfsys@useobject{currentmarker}{}%
\end{pgfscope}%
\end{pgfscope}%
\begin{pgfscope}%
\pgfpathrectangle{\pgfqpoint{0.584940in}{1.184369in}}{\pgfqpoint{2.440060in}{2.333202in}}%
\pgfusepath{clip}%
\pgfsetbuttcap%
\pgfsetroundjoin%
\pgfsetlinewidth{1.505625pt}%
\definecolor{currentstroke}{rgb}{0.992157,0.682353,0.380392}%
\pgfsetstrokecolor{currentstroke}%
\pgfsetdash{}{0pt}%
\pgfpathmoveto{\pgfqpoint{1.113921in}{2.033242in}}%
\pgfpathlineto{\pgfqpoint{1.947005in}{2.033242in}}%
\pgfusepath{stroke}%
\end{pgfscope}%
\begin{pgfscope}%
\pgfpathrectangle{\pgfqpoint{0.584940in}{1.184369in}}{\pgfqpoint{2.440060in}{2.333202in}}%
\pgfusepath{clip}%
\pgfsetbuttcap%
\pgfsetroundjoin%
\pgfsetlinewidth{1.505625pt}%
\definecolor{currentstroke}{rgb}{0.992157,0.682353,0.380392}%
\pgfsetstrokecolor{currentstroke}%
\pgfsetdash{}{0pt}%
\pgfpathmoveto{\pgfqpoint{1.530463in}{1.684488in}}%
\pgfpathlineto{\pgfqpoint{1.530463in}{2.381997in}}%
\pgfusepath{stroke}%
\end{pgfscope}%
\begin{pgfscope}%
\pgfpathrectangle{\pgfqpoint{0.584940in}{1.184369in}}{\pgfqpoint{2.440060in}{2.333202in}}%
\pgfusepath{clip}%
\pgfsetbuttcap%
\pgfsetroundjoin%
\definecolor{currentfill}{rgb}{0.992157,0.682353,0.380392}%
\pgfsetfillcolor{currentfill}%
\pgfsetlinewidth{1.003750pt}%
\definecolor{currentstroke}{rgb}{0.992157,0.682353,0.380392}%
\pgfsetstrokecolor{currentstroke}%
\pgfsetdash{}{0pt}%
\pgfsys@defobject{currentmarker}{\pgfqpoint{0.000000in}{-0.069444in}}{\pgfqpoint{0.000000in}{0.069444in}}{%
\pgfpathmoveto{\pgfqpoint{0.000000in}{-0.069444in}}%
\pgfpathlineto{\pgfqpoint{0.000000in}{0.069444in}}%
\pgfusepath{stroke,fill}%
}%
\begin{pgfscope}%
\pgfsys@transformshift{1.113921in}{2.033242in}%
\pgfsys@useobject{currentmarker}{}%
\end{pgfscope}%
\end{pgfscope}%
\begin{pgfscope}%
\pgfpathrectangle{\pgfqpoint{0.584940in}{1.184369in}}{\pgfqpoint{2.440060in}{2.333202in}}%
\pgfusepath{clip}%
\pgfsetbuttcap%
\pgfsetroundjoin%
\definecolor{currentfill}{rgb}{0.992157,0.682353,0.380392}%
\pgfsetfillcolor{currentfill}%
\pgfsetlinewidth{1.003750pt}%
\definecolor{currentstroke}{rgb}{0.992157,0.682353,0.380392}%
\pgfsetstrokecolor{currentstroke}%
\pgfsetdash{}{0pt}%
\pgfsys@defobject{currentmarker}{\pgfqpoint{0.000000in}{-0.069444in}}{\pgfqpoint{0.000000in}{0.069444in}}{%
\pgfpathmoveto{\pgfqpoint{0.000000in}{-0.069444in}}%
\pgfpathlineto{\pgfqpoint{0.000000in}{0.069444in}}%
\pgfusepath{stroke,fill}%
}%
\begin{pgfscope}%
\pgfsys@transformshift{1.947005in}{2.033242in}%
\pgfsys@useobject{currentmarker}{}%
\end{pgfscope}%
\end{pgfscope}%
\begin{pgfscope}%
\pgfpathrectangle{\pgfqpoint{0.584940in}{1.184369in}}{\pgfqpoint{2.440060in}{2.333202in}}%
\pgfusepath{clip}%
\pgfsetbuttcap%
\pgfsetroundjoin%
\definecolor{currentfill}{rgb}{0.992157,0.682353,0.380392}%
\pgfsetfillcolor{currentfill}%
\pgfsetlinewidth{1.003750pt}%
\definecolor{currentstroke}{rgb}{0.992157,0.682353,0.380392}%
\pgfsetstrokecolor{currentstroke}%
\pgfsetdash{}{0pt}%
\pgfsys@defobject{currentmarker}{\pgfqpoint{-0.069444in}{-0.000000in}}{\pgfqpoint{0.069444in}{0.000000in}}{%
\pgfpathmoveto{\pgfqpoint{0.069444in}{-0.000000in}}%
\pgfpathlineto{\pgfqpoint{-0.069444in}{0.000000in}}%
\pgfusepath{stroke,fill}%
}%
\begin{pgfscope}%
\pgfsys@transformshift{1.530463in}{1.684488in}%
\pgfsys@useobject{currentmarker}{}%
\end{pgfscope}%
\end{pgfscope}%
\begin{pgfscope}%
\pgfpathrectangle{\pgfqpoint{0.584940in}{1.184369in}}{\pgfqpoint{2.440060in}{2.333202in}}%
\pgfusepath{clip}%
\pgfsetbuttcap%
\pgfsetroundjoin%
\definecolor{currentfill}{rgb}{0.992157,0.682353,0.380392}%
\pgfsetfillcolor{currentfill}%
\pgfsetlinewidth{1.003750pt}%
\definecolor{currentstroke}{rgb}{0.992157,0.682353,0.380392}%
\pgfsetstrokecolor{currentstroke}%
\pgfsetdash{}{0pt}%
\pgfsys@defobject{currentmarker}{\pgfqpoint{-0.069444in}{-0.000000in}}{\pgfqpoint{0.069444in}{0.000000in}}{%
\pgfpathmoveto{\pgfqpoint{0.069444in}{-0.000000in}}%
\pgfpathlineto{\pgfqpoint{-0.069444in}{0.000000in}}%
\pgfusepath{stroke,fill}%
}%
\begin{pgfscope}%
\pgfsys@transformshift{1.530463in}{2.381997in}%
\pgfsys@useobject{currentmarker}{}%
\end{pgfscope}%
\end{pgfscope}%
\begin{pgfscope}%
\pgfpathrectangle{\pgfqpoint{0.584940in}{1.184369in}}{\pgfqpoint{2.440060in}{2.333202in}}%
\pgfusepath{clip}%
\pgfsetbuttcap%
\pgfsetroundjoin%
\pgfsetlinewidth{1.505625pt}%
\definecolor{currentstroke}{rgb}{0.956863,0.427451,0.262745}%
\pgfsetstrokecolor{currentstroke}%
\pgfsetdash{}{0pt}%
\pgfpathmoveto{\pgfqpoint{0.574940in}{2.303619in}}%
\pgfpathlineto{\pgfqpoint{1.711214in}{2.303619in}}%
\pgfusepath{stroke}%
\end{pgfscope}%
\begin{pgfscope}%
\pgfpathrectangle{\pgfqpoint{0.584940in}{1.184369in}}{\pgfqpoint{2.440060in}{2.333202in}}%
\pgfusepath{clip}%
\pgfsetbuttcap%
\pgfsetroundjoin%
\pgfsetlinewidth{1.505625pt}%
\definecolor{currentstroke}{rgb}{0.956863,0.427451,0.262745}%
\pgfsetstrokecolor{currentstroke}%
\pgfsetdash{}{0pt}%
\pgfpathmoveto{\pgfqpoint{0.868161in}{1.924910in}}%
\pgfpathlineto{\pgfqpoint{0.868161in}{2.682328in}}%
\pgfusepath{stroke}%
\end{pgfscope}%
\begin{pgfscope}%
\pgfpathrectangle{\pgfqpoint{0.584940in}{1.184369in}}{\pgfqpoint{2.440060in}{2.333202in}}%
\pgfusepath{clip}%
\pgfsetbuttcap%
\pgfsetroundjoin%
\definecolor{currentfill}{rgb}{0.956863,0.427451,0.262745}%
\pgfsetfillcolor{currentfill}%
\pgfsetlinewidth{2.007500pt}%
\definecolor{currentstroke}{rgb}{0.956863,0.427451,0.262745}%
\pgfsetstrokecolor{currentstroke}%
\pgfsetdash{}{0pt}%
\pgfsys@defobject{currentmarker}{\pgfqpoint{0.000000in}{-0.069444in}}{\pgfqpoint{0.000000in}{0.069444in}}{%
\pgfpathmoveto{\pgfqpoint{0.000000in}{-0.069444in}}%
\pgfpathlineto{\pgfqpoint{0.000000in}{0.069444in}}%
\pgfusepath{stroke,fill}%
}%
\begin{pgfscope}%
\pgfsys@transformshift{0.025109in}{2.303619in}%
\pgfsys@useobject{currentmarker}{}%
\end{pgfscope}%
\end{pgfscope}%
\begin{pgfscope}%
\pgfpathrectangle{\pgfqpoint{0.584940in}{1.184369in}}{\pgfqpoint{2.440060in}{2.333202in}}%
\pgfusepath{clip}%
\pgfsetbuttcap%
\pgfsetroundjoin%
\definecolor{currentfill}{rgb}{0.956863,0.427451,0.262745}%
\pgfsetfillcolor{currentfill}%
\pgfsetlinewidth{2.007500pt}%
\definecolor{currentstroke}{rgb}{0.956863,0.427451,0.262745}%
\pgfsetstrokecolor{currentstroke}%
\pgfsetdash{}{0pt}%
\pgfsys@defobject{currentmarker}{\pgfqpoint{0.000000in}{-0.069444in}}{\pgfqpoint{0.000000in}{0.069444in}}{%
\pgfpathmoveto{\pgfqpoint{0.000000in}{-0.069444in}}%
\pgfpathlineto{\pgfqpoint{0.000000in}{0.069444in}}%
\pgfusepath{stroke,fill}%
}%
\begin{pgfscope}%
\pgfsys@transformshift{1.711214in}{2.303619in}%
\pgfsys@useobject{currentmarker}{}%
\end{pgfscope}%
\end{pgfscope}%
\begin{pgfscope}%
\pgfpathrectangle{\pgfqpoint{0.584940in}{1.184369in}}{\pgfqpoint{2.440060in}{2.333202in}}%
\pgfusepath{clip}%
\pgfsetbuttcap%
\pgfsetroundjoin%
\definecolor{currentfill}{rgb}{0.956863,0.427451,0.262745}%
\pgfsetfillcolor{currentfill}%
\pgfsetlinewidth{2.007500pt}%
\definecolor{currentstroke}{rgb}{0.956863,0.427451,0.262745}%
\pgfsetstrokecolor{currentstroke}%
\pgfsetdash{}{0pt}%
\pgfsys@defobject{currentmarker}{\pgfqpoint{-0.069444in}{-0.000000in}}{\pgfqpoint{0.069444in}{0.000000in}}{%
\pgfpathmoveto{\pgfqpoint{0.069444in}{-0.000000in}}%
\pgfpathlineto{\pgfqpoint{-0.069444in}{0.000000in}}%
\pgfusepath{stroke,fill}%
}%
\begin{pgfscope}%
\pgfsys@transformshift{0.868161in}{1.924910in}%
\pgfsys@useobject{currentmarker}{}%
\end{pgfscope}%
\end{pgfscope}%
\begin{pgfscope}%
\pgfpathrectangle{\pgfqpoint{0.584940in}{1.184369in}}{\pgfqpoint{2.440060in}{2.333202in}}%
\pgfusepath{clip}%
\pgfsetbuttcap%
\pgfsetroundjoin%
\definecolor{currentfill}{rgb}{0.956863,0.427451,0.262745}%
\pgfsetfillcolor{currentfill}%
\pgfsetlinewidth{2.007500pt}%
\definecolor{currentstroke}{rgb}{0.956863,0.427451,0.262745}%
\pgfsetstrokecolor{currentstroke}%
\pgfsetdash{}{0pt}%
\pgfsys@defobject{currentmarker}{\pgfqpoint{-0.069444in}{-0.000000in}}{\pgfqpoint{0.069444in}{0.000000in}}{%
\pgfpathmoveto{\pgfqpoint{0.069444in}{-0.000000in}}%
\pgfpathlineto{\pgfqpoint{-0.069444in}{0.000000in}}%
\pgfusepath{stroke,fill}%
}%
\begin{pgfscope}%
\pgfsys@transformshift{0.868161in}{2.682328in}%
\pgfsys@useobject{currentmarker}{}%
\end{pgfscope}%
\end{pgfscope}%
\begin{pgfscope}%
\pgfpathrectangle{\pgfqpoint{0.584940in}{1.184369in}}{\pgfqpoint{2.440060in}{2.333202in}}%
\pgfusepath{clip}%
\pgfsetbuttcap%
\pgfsetroundjoin%
\pgfsetlinewidth{1.505625pt}%
\definecolor{currentstroke}{rgb}{0.843137,0.188235,0.152941}%
\pgfsetstrokecolor{currentstroke}%
\pgfsetdash{}{0pt}%
\pgfpathmoveto{\pgfqpoint{1.598585in}{2.851922in}}%
\pgfpathlineto{\pgfqpoint{3.035000in}{2.851922in}}%
\pgfusepath{stroke}%
\end{pgfscope}%
\begin{pgfscope}%
\pgfpathrectangle{\pgfqpoint{0.584940in}{1.184369in}}{\pgfqpoint{2.440060in}{2.333202in}}%
\pgfusepath{clip}%
\pgfsetbuttcap%
\pgfsetroundjoin%
\pgfsetlinewidth{1.505625pt}%
\definecolor{currentstroke}{rgb}{0.843137,0.188235,0.152941}%
\pgfsetstrokecolor{currentstroke}%
\pgfsetdash{}{0pt}%
\pgfpathmoveto{\pgfqpoint{2.410628in}{2.392887in}}%
\pgfpathlineto{\pgfqpoint{2.410628in}{3.310957in}}%
\pgfusepath{stroke}%
\end{pgfscope}%
\begin{pgfscope}%
\pgfpathrectangle{\pgfqpoint{0.584940in}{1.184369in}}{\pgfqpoint{2.440060in}{2.333202in}}%
\pgfusepath{clip}%
\pgfsetbuttcap%
\pgfsetroundjoin%
\definecolor{currentfill}{rgb}{0.843137,0.188235,0.152941}%
\pgfsetfillcolor{currentfill}%
\pgfsetlinewidth{1.003750pt}%
\definecolor{currentstroke}{rgb}{0.843137,0.188235,0.152941}%
\pgfsetstrokecolor{currentstroke}%
\pgfsetdash{}{0pt}%
\pgfsys@defobject{currentmarker}{\pgfqpoint{0.000000in}{-0.069444in}}{\pgfqpoint{0.000000in}{0.069444in}}{%
\pgfpathmoveto{\pgfqpoint{0.000000in}{-0.069444in}}%
\pgfpathlineto{\pgfqpoint{0.000000in}{0.069444in}}%
\pgfusepath{stroke,fill}%
}%
\begin{pgfscope}%
\pgfsys@transformshift{1.598585in}{2.851922in}%
\pgfsys@useobject{currentmarker}{}%
\end{pgfscope}%
\end{pgfscope}%
\begin{pgfscope}%
\pgfpathrectangle{\pgfqpoint{0.584940in}{1.184369in}}{\pgfqpoint{2.440060in}{2.333202in}}%
\pgfusepath{clip}%
\pgfsetbuttcap%
\pgfsetroundjoin%
\definecolor{currentfill}{rgb}{0.843137,0.188235,0.152941}%
\pgfsetfillcolor{currentfill}%
\pgfsetlinewidth{1.003750pt}%
\definecolor{currentstroke}{rgb}{0.843137,0.188235,0.152941}%
\pgfsetstrokecolor{currentstroke}%
\pgfsetdash{}{0pt}%
\pgfsys@defobject{currentmarker}{\pgfqpoint{0.000000in}{-0.069444in}}{\pgfqpoint{0.000000in}{0.069444in}}{%
\pgfpathmoveto{\pgfqpoint{0.000000in}{-0.069444in}}%
\pgfpathlineto{\pgfqpoint{0.000000in}{0.069444in}}%
\pgfusepath{stroke,fill}%
}%
\begin{pgfscope}%
\pgfsys@transformshift{3.222670in}{2.851922in}%
\pgfsys@useobject{currentmarker}{}%
\end{pgfscope}%
\end{pgfscope}%
\begin{pgfscope}%
\pgfpathrectangle{\pgfqpoint{0.584940in}{1.184369in}}{\pgfqpoint{2.440060in}{2.333202in}}%
\pgfusepath{clip}%
\pgfsetbuttcap%
\pgfsetroundjoin%
\definecolor{currentfill}{rgb}{0.843137,0.188235,0.152941}%
\pgfsetfillcolor{currentfill}%
\pgfsetlinewidth{1.003750pt}%
\definecolor{currentstroke}{rgb}{0.843137,0.188235,0.152941}%
\pgfsetstrokecolor{currentstroke}%
\pgfsetdash{}{0pt}%
\pgfsys@defobject{currentmarker}{\pgfqpoint{-0.069444in}{-0.000000in}}{\pgfqpoint{0.069444in}{0.000000in}}{%
\pgfpathmoveto{\pgfqpoint{0.069444in}{-0.000000in}}%
\pgfpathlineto{\pgfqpoint{-0.069444in}{0.000000in}}%
\pgfusepath{stroke,fill}%
}%
\begin{pgfscope}%
\pgfsys@transformshift{2.410628in}{2.392887in}%
\pgfsys@useobject{currentmarker}{}%
\end{pgfscope}%
\end{pgfscope}%
\begin{pgfscope}%
\pgfpathrectangle{\pgfqpoint{0.584940in}{1.184369in}}{\pgfqpoint{2.440060in}{2.333202in}}%
\pgfusepath{clip}%
\pgfsetbuttcap%
\pgfsetroundjoin%
\definecolor{currentfill}{rgb}{0.843137,0.188235,0.152941}%
\pgfsetfillcolor{currentfill}%
\pgfsetlinewidth{1.003750pt}%
\definecolor{currentstroke}{rgb}{0.843137,0.188235,0.152941}%
\pgfsetstrokecolor{currentstroke}%
\pgfsetdash{}{0pt}%
\pgfsys@defobject{currentmarker}{\pgfqpoint{-0.069444in}{-0.000000in}}{\pgfqpoint{0.069444in}{0.000000in}}{%
\pgfpathmoveto{\pgfqpoint{0.069444in}{-0.000000in}}%
\pgfpathlineto{\pgfqpoint{-0.069444in}{0.000000in}}%
\pgfusepath{stroke,fill}%
}%
\begin{pgfscope}%
\pgfsys@transformshift{2.410628in}{3.310957in}%
\pgfsys@useobject{currentmarker}{}%
\end{pgfscope}%
\end{pgfscope}%
\begin{pgfscope}%
\pgfpathrectangle{\pgfqpoint{0.584940in}{1.184369in}}{\pgfqpoint{2.440060in}{2.333202in}}%
\pgfusepath{clip}%
\pgfsetbuttcap%
\pgfsetroundjoin%
\definecolor{currentfill}{rgb}{0.129412,0.400000,0.674510}%
\pgfsetfillcolor{currentfill}%
\pgfsetlinewidth{1.003750pt}%
\definecolor{currentstroke}{rgb}{0.129412,0.400000,0.674510}%
\pgfsetstrokecolor{currentstroke}%
\pgfsetdash{}{0pt}%
\pgfsys@defobject{currentmarker}{\pgfqpoint{-0.055556in}{-0.055556in}}{\pgfqpoint{0.055556in}{0.055556in}}{%
\pgfpathmoveto{\pgfqpoint{0.000000in}{-0.055556in}}%
\pgfpathcurveto{\pgfqpoint{0.014734in}{-0.055556in}}{\pgfqpoint{0.028866in}{-0.049702in}}{\pgfqpoint{0.039284in}{-0.039284in}}%
\pgfpathcurveto{\pgfqpoint{0.049702in}{-0.028866in}}{\pgfqpoint{0.055556in}{-0.014734in}}{\pgfqpoint{0.055556in}{0.000000in}}%
\pgfpathcurveto{\pgfqpoint{0.055556in}{0.014734in}}{\pgfqpoint{0.049702in}{0.028866in}}{\pgfqpoint{0.039284in}{0.039284in}}%
\pgfpathcurveto{\pgfqpoint{0.028866in}{0.049702in}}{\pgfqpoint{0.014734in}{0.055556in}}{\pgfqpoint{0.000000in}{0.055556in}}%
\pgfpathcurveto{\pgfqpoint{-0.014734in}{0.055556in}}{\pgfqpoint{-0.028866in}{0.049702in}}{\pgfqpoint{-0.039284in}{0.039284in}}%
\pgfpathcurveto{\pgfqpoint{-0.049702in}{0.028866in}}{\pgfqpoint{-0.055556in}{0.014734in}}{\pgfqpoint{-0.055556in}{0.000000in}}%
\pgfpathcurveto{\pgfqpoint{-0.055556in}{-0.014734in}}{\pgfqpoint{-0.049702in}{-0.028866in}}{\pgfqpoint{-0.039284in}{-0.039284in}}%
\pgfpathcurveto{\pgfqpoint{-0.028866in}{-0.049702in}}{\pgfqpoint{-0.014734in}{-0.055556in}}{\pgfqpoint{0.000000in}{-0.055556in}}%
\pgfpathlineto{\pgfqpoint{0.000000in}{-0.055556in}}%
\pgfpathclose%
\pgfusepath{stroke,fill}%
}%
\begin{pgfscope}%
\pgfsys@transformshift{1.360530in}{1.707280in}%
\pgfsys@useobject{currentmarker}{}%
\end{pgfscope}%
\end{pgfscope}%
\begin{pgfscope}%
\pgfpathrectangle{\pgfqpoint{0.584940in}{1.184369in}}{\pgfqpoint{2.440060in}{2.333202in}}%
\pgfusepath{clip}%
\pgfsetbuttcap%
\pgfsetroundjoin%
\definecolor{currentfill}{rgb}{0.454902,0.678431,0.819608}%
\pgfsetfillcolor{currentfill}%
\pgfsetlinewidth{1.003750pt}%
\definecolor{currentstroke}{rgb}{0.454902,0.678431,0.819608}%
\pgfsetstrokecolor{currentstroke}%
\pgfsetdash{}{0pt}%
\pgfsys@defobject{currentmarker}{\pgfqpoint{-0.055556in}{-0.055556in}}{\pgfqpoint{0.055556in}{0.055556in}}{%
\pgfpathmoveto{\pgfqpoint{0.000000in}{-0.055556in}}%
\pgfpathcurveto{\pgfqpoint{0.014734in}{-0.055556in}}{\pgfqpoint{0.028866in}{-0.049702in}}{\pgfqpoint{0.039284in}{-0.039284in}}%
\pgfpathcurveto{\pgfqpoint{0.049702in}{-0.028866in}}{\pgfqpoint{0.055556in}{-0.014734in}}{\pgfqpoint{0.055556in}{0.000000in}}%
\pgfpathcurveto{\pgfqpoint{0.055556in}{0.014734in}}{\pgfqpoint{0.049702in}{0.028866in}}{\pgfqpoint{0.039284in}{0.039284in}}%
\pgfpathcurveto{\pgfqpoint{0.028866in}{0.049702in}}{\pgfqpoint{0.014734in}{0.055556in}}{\pgfqpoint{0.000000in}{0.055556in}}%
\pgfpathcurveto{\pgfqpoint{-0.014734in}{0.055556in}}{\pgfqpoint{-0.028866in}{0.049702in}}{\pgfqpoint{-0.039284in}{0.039284in}}%
\pgfpathcurveto{\pgfqpoint{-0.049702in}{0.028866in}}{\pgfqpoint{-0.055556in}{0.014734in}}{\pgfqpoint{-0.055556in}{0.000000in}}%
\pgfpathcurveto{\pgfqpoint{-0.055556in}{-0.014734in}}{\pgfqpoint{-0.049702in}{-0.028866in}}{\pgfqpoint{-0.039284in}{-0.039284in}}%
\pgfpathcurveto{\pgfqpoint{-0.028866in}{-0.049702in}}{\pgfqpoint{-0.014734in}{-0.055556in}}{\pgfqpoint{0.000000in}{-0.055556in}}%
\pgfpathlineto{\pgfqpoint{0.000000in}{-0.055556in}}%
\pgfpathclose%
\pgfusepath{stroke,fill}%
}%
\begin{pgfscope}%
\pgfsys@transformshift{1.046808in}{1.983147in}%
\pgfsys@useobject{currentmarker}{}%
\end{pgfscope}%
\end{pgfscope}%
\begin{pgfscope}%
\pgfpathrectangle{\pgfqpoint{0.584940in}{1.184369in}}{\pgfqpoint{2.440060in}{2.333202in}}%
\pgfusepath{clip}%
\pgfsetbuttcap%
\pgfsetroundjoin%
\definecolor{currentfill}{rgb}{0.992157,0.682353,0.380392}%
\pgfsetfillcolor{currentfill}%
\pgfsetlinewidth{1.003750pt}%
\definecolor{currentstroke}{rgb}{0.992157,0.682353,0.380392}%
\pgfsetstrokecolor{currentstroke}%
\pgfsetdash{}{0pt}%
\pgfsys@defobject{currentmarker}{\pgfqpoint{-0.055556in}{-0.055556in}}{\pgfqpoint{0.055556in}{0.055556in}}{%
\pgfpathmoveto{\pgfqpoint{0.000000in}{-0.055556in}}%
\pgfpathcurveto{\pgfqpoint{0.014734in}{-0.055556in}}{\pgfqpoint{0.028866in}{-0.049702in}}{\pgfqpoint{0.039284in}{-0.039284in}}%
\pgfpathcurveto{\pgfqpoint{0.049702in}{-0.028866in}}{\pgfqpoint{0.055556in}{-0.014734in}}{\pgfqpoint{0.055556in}{0.000000in}}%
\pgfpathcurveto{\pgfqpoint{0.055556in}{0.014734in}}{\pgfqpoint{0.049702in}{0.028866in}}{\pgfqpoint{0.039284in}{0.039284in}}%
\pgfpathcurveto{\pgfqpoint{0.028866in}{0.049702in}}{\pgfqpoint{0.014734in}{0.055556in}}{\pgfqpoint{0.000000in}{0.055556in}}%
\pgfpathcurveto{\pgfqpoint{-0.014734in}{0.055556in}}{\pgfqpoint{-0.028866in}{0.049702in}}{\pgfqpoint{-0.039284in}{0.039284in}}%
\pgfpathcurveto{\pgfqpoint{-0.049702in}{0.028866in}}{\pgfqpoint{-0.055556in}{0.014734in}}{\pgfqpoint{-0.055556in}{0.000000in}}%
\pgfpathcurveto{\pgfqpoint{-0.055556in}{-0.014734in}}{\pgfqpoint{-0.049702in}{-0.028866in}}{\pgfqpoint{-0.039284in}{-0.039284in}}%
\pgfpathcurveto{\pgfqpoint{-0.028866in}{-0.049702in}}{\pgfqpoint{-0.014734in}{-0.055556in}}{\pgfqpoint{0.000000in}{-0.055556in}}%
\pgfpathlineto{\pgfqpoint{0.000000in}{-0.055556in}}%
\pgfpathclose%
\pgfusepath{stroke,fill}%
}%
\begin{pgfscope}%
\pgfsys@transformshift{1.530463in}{2.033242in}%
\pgfsys@useobject{currentmarker}{}%
\end{pgfscope}%
\end{pgfscope}%
\begin{pgfscope}%
\pgfpathrectangle{\pgfqpoint{0.584940in}{1.184369in}}{\pgfqpoint{2.440060in}{2.333202in}}%
\pgfusepath{clip}%
\pgfsetbuttcap%
\pgfsetmiterjoin%
\definecolor{currentfill}{rgb}{0.000000,0.000000,0.000000}%
\pgfsetfillcolor{currentfill}%
\pgfsetfillopacity{0.000000}%
\pgfsetlinewidth{2.007500pt}%
\definecolor{currentstroke}{rgb}{0.956863,0.427451,0.262745}%
\pgfsetstrokecolor{currentstroke}%
\pgfsetdash{}{0pt}%
\pgfsys@defobject{currentmarker}{\pgfqpoint{-0.078567in}{-0.078567in}}{\pgfqpoint{0.078567in}{0.078567in}}{%
\pgfpathmoveto{\pgfqpoint{-0.000000in}{-0.078567in}}%
\pgfpathlineto{\pgfqpoint{0.078567in}{0.000000in}}%
\pgfpathlineto{\pgfqpoint{0.000000in}{0.078567in}}%
\pgfpathlineto{\pgfqpoint{-0.078567in}{0.000000in}}%
\pgfpathlineto{\pgfqpoint{-0.000000in}{-0.078567in}}%
\pgfpathclose%
\pgfusepath{stroke,fill}%
}%
\begin{pgfscope}%
\pgfsys@transformshift{0.868161in}{2.303619in}%
\pgfsys@useobject{currentmarker}{}%
\end{pgfscope}%
\end{pgfscope}%
\begin{pgfscope}%
\pgfpathrectangle{\pgfqpoint{0.584940in}{1.184369in}}{\pgfqpoint{2.440060in}{2.333202in}}%
\pgfusepath{clip}%
\pgfsetbuttcap%
\pgfsetroundjoin%
\definecolor{currentfill}{rgb}{0.843137,0.188235,0.152941}%
\pgfsetfillcolor{currentfill}%
\pgfsetlinewidth{1.003750pt}%
\definecolor{currentstroke}{rgb}{0.843137,0.188235,0.152941}%
\pgfsetstrokecolor{currentstroke}%
\pgfsetdash{}{0pt}%
\pgfsys@defobject{currentmarker}{\pgfqpoint{-0.055556in}{-0.055556in}}{\pgfqpoint{0.055556in}{0.055556in}}{%
\pgfpathmoveto{\pgfqpoint{0.000000in}{-0.055556in}}%
\pgfpathcurveto{\pgfqpoint{0.014734in}{-0.055556in}}{\pgfqpoint{0.028866in}{-0.049702in}}{\pgfqpoint{0.039284in}{-0.039284in}}%
\pgfpathcurveto{\pgfqpoint{0.049702in}{-0.028866in}}{\pgfqpoint{0.055556in}{-0.014734in}}{\pgfqpoint{0.055556in}{0.000000in}}%
\pgfpathcurveto{\pgfqpoint{0.055556in}{0.014734in}}{\pgfqpoint{0.049702in}{0.028866in}}{\pgfqpoint{0.039284in}{0.039284in}}%
\pgfpathcurveto{\pgfqpoint{0.028866in}{0.049702in}}{\pgfqpoint{0.014734in}{0.055556in}}{\pgfqpoint{0.000000in}{0.055556in}}%
\pgfpathcurveto{\pgfqpoint{-0.014734in}{0.055556in}}{\pgfqpoint{-0.028866in}{0.049702in}}{\pgfqpoint{-0.039284in}{0.039284in}}%
\pgfpathcurveto{\pgfqpoint{-0.049702in}{0.028866in}}{\pgfqpoint{-0.055556in}{0.014734in}}{\pgfqpoint{-0.055556in}{0.000000in}}%
\pgfpathcurveto{\pgfqpoint{-0.055556in}{-0.014734in}}{\pgfqpoint{-0.049702in}{-0.028866in}}{\pgfqpoint{-0.039284in}{-0.039284in}}%
\pgfpathcurveto{\pgfqpoint{-0.028866in}{-0.049702in}}{\pgfqpoint{-0.014734in}{-0.055556in}}{\pgfqpoint{0.000000in}{-0.055556in}}%
\pgfpathlineto{\pgfqpoint{0.000000in}{-0.055556in}}%
\pgfpathclose%
\pgfusepath{stroke,fill}%
}%
\begin{pgfscope}%
\pgfsys@transformshift{2.410628in}{2.851922in}%
\pgfsys@useobject{currentmarker}{}%
\end{pgfscope}%
\end{pgfscope}%
\begin{pgfscope}%
\pgfsetrectcap%
\pgfsetmiterjoin%
\pgfsetlinewidth{0.803000pt}%
\definecolor{currentstroke}{rgb}{0.000000,0.000000,0.000000}%
\pgfsetstrokecolor{currentstroke}%
\pgfsetdash{}{0pt}%
\pgfpathmoveto{\pgfqpoint{0.584940in}{1.184369in}}%
\pgfpathlineto{\pgfqpoint{0.584940in}{3.517571in}}%
\pgfusepath{stroke}%
\end{pgfscope}%
\begin{pgfscope}%
\pgfsetrectcap%
\pgfsetmiterjoin%
\pgfsetlinewidth{0.803000pt}%
\definecolor{currentstroke}{rgb}{0.000000,0.000000,0.000000}%
\pgfsetstrokecolor{currentstroke}%
\pgfsetdash{}{0pt}%
\pgfpathmoveto{\pgfqpoint{3.025000in}{1.184369in}}%
\pgfpathlineto{\pgfqpoint{3.025000in}{3.517571in}}%
\pgfusepath{stroke}%
\end{pgfscope}%
\begin{pgfscope}%
\pgfsetrectcap%
\pgfsetmiterjoin%
\pgfsetlinewidth{0.803000pt}%
\definecolor{currentstroke}{rgb}{0.000000,0.000000,0.000000}%
\pgfsetstrokecolor{currentstroke}%
\pgfsetdash{}{0pt}%
\pgfpathmoveto{\pgfqpoint{0.584940in}{1.184369in}}%
\pgfpathlineto{\pgfqpoint{3.025000in}{1.184369in}}%
\pgfusepath{stroke}%
\end{pgfscope}%
\begin{pgfscope}%
\pgfsetrectcap%
\pgfsetmiterjoin%
\pgfsetlinewidth{0.803000pt}%
\definecolor{currentstroke}{rgb}{0.000000,0.000000,0.000000}%
\pgfsetstrokecolor{currentstroke}%
\pgfsetdash{}{0pt}%
\pgfpathmoveto{\pgfqpoint{0.584940in}{3.517571in}}%
\pgfpathlineto{\pgfqpoint{3.025000in}{3.517571in}}%
\pgfusepath{stroke}%
\end{pgfscope}%
\begin{pgfscope}%
\pgfsetbuttcap%
\pgfsetmiterjoin%
\definecolor{currentfill}{rgb}{1.000000,1.000000,1.000000}%
\pgfsetfillcolor{currentfill}%
\pgfsetfillopacity{0.800000}%
\pgfsetlinewidth{1.003750pt}%
\definecolor{currentstroke}{rgb}{0.000000,0.000000,0.000000}%
\pgfsetstrokecolor{currentstroke}%
\pgfsetstrokeopacity{0.800000}%
\pgfsetdash{}{0pt}%
\pgfpathmoveto{\pgfqpoint{2.072417in}{1.271862in}}%
\pgfpathlineto{\pgfqpoint{2.951798in}{1.271862in}}%
\pgfpathquadraticcurveto{\pgfqpoint{2.980965in}{1.271862in}}{\pgfqpoint{2.980965in}{1.301029in}}%
\pgfpathlineto{\pgfqpoint{2.980965in}{1.626807in}}%
\pgfpathquadraticcurveto{\pgfqpoint{2.980965in}{1.655974in}}{\pgfqpoint{2.951798in}{1.655974in}}%
\pgfpathlineto{\pgfqpoint{2.072417in}{1.655974in}}%
\pgfpathquadraticcurveto{\pgfqpoint{2.043250in}{1.655974in}}{\pgfqpoint{2.043250in}{1.626807in}}%
\pgfpathlineto{\pgfqpoint{2.043250in}{1.301029in}}%
\pgfpathquadraticcurveto{\pgfqpoint{2.043250in}{1.271862in}}{\pgfqpoint{2.072417in}{1.271862in}}%
\pgfpathlineto{\pgfqpoint{2.072417in}{1.271862in}}%
\pgfpathclose%
\pgfusepath{stroke,fill}%
\end{pgfscope}%
\begin{pgfscope}%
\definecolor{textcolor}{rgb}{0.000000,0.000000,0.000000}%
\pgfsetstrokecolor{textcolor}%
\pgfsetfillcolor{textcolor}%
\pgftext[x=2.421575in, y=1.556468in, left, base]{\color{textcolor}{\rmfamily\fontsize{7.000000}{8.400000}\selectfont\catcode`\^=\active\def^{\ifmmode\sp\else\^{}\fi}\catcode`\%=\active\def%{\%}slope $= 3.02$}}%
\end{pgfscope}%
\begin{pgfscope}%
\definecolor{textcolor}{rgb}{0.000000,0.000000,0.000000}%
\pgfsetstrokecolor{textcolor}%
\pgfsetfillcolor{textcolor}%
\pgftext[x=2.072417in, y=1.440268in, left, base]{\color{textcolor}{\rmfamily\fontsize{7.000000}{8.400000}\selectfont\catcode`\^=\active\def^{\ifmmode\sp\else\^{}\fi}\catcode`\%=\active\def%{\%}90\% CI $= [-2.4,\,5.7]$}}%
\end{pgfscope}%
\begin{pgfscope}%
\definecolor{textcolor}{rgb}{0.000000,0.000000,0.000000}%
\pgfsetstrokecolor{textcolor}%
\pgfsetfillcolor{textcolor}%
\pgftext[x=2.266320in, y=1.328445in, left, base]{\color{textcolor}{\rmfamily\fontsize{7.000000}{8.400000}\selectfont\catcode`\^=\active\def^{\ifmmode\sp\else\^{}\fi}\catcode`\%=\active\def%{\%}$1.09\,\sigma$ from zero}}%
\end{pgfscope}%
\begin{pgfscope}%
\definecolor{textcolor}{rgb}{0.000000,0.000000,0.000000}%
\pgfsetstrokecolor{textcolor}%
\pgfsetfillcolor{textcolor}%
\pgftext[x=1.804970in,y=3.600904in,,base]{\color{textcolor}{\rmfamily\fontsize{10.000000}{12.000000}\selectfont\catcode`\^=\active\def^{\ifmmode\sp\else\^{}\fi}\catcode`\%=\active\def%{\%}(a)}}%
\end{pgfscope}%
\begin{pgfscope}%
\pgfsetbuttcap%
\pgfsetmiterjoin%
\definecolor{currentfill}{rgb}{1.000000,1.000000,1.000000}%
\pgfsetfillcolor{currentfill}%
\pgfsetlinewidth{0.000000pt}%
\definecolor{currentstroke}{rgb}{0.000000,0.000000,0.000000}%
\pgfsetstrokecolor{currentstroke}%
\pgfsetstrokeopacity{0.000000}%
\pgfsetdash{}{0pt}%
\pgfpathmoveto{\pgfqpoint{3.659940in}{1.184369in}}%
\pgfpathlineto{\pgfqpoint{6.100000in}{1.184369in}}%
\pgfpathlineto{\pgfqpoint{6.100000in}{3.517571in}}%
\pgfpathlineto{\pgfqpoint{3.659940in}{3.517571in}}%
\pgfpathlineto{\pgfqpoint{3.659940in}{1.184369in}}%
\pgfpathclose%
\pgfusepath{fill}%
\end{pgfscope}%
\begin{pgfscope}%
\pgfpathrectangle{\pgfqpoint{3.659940in}{1.184369in}}{\pgfqpoint{2.440060in}{2.333202in}}%
\pgfusepath{clip}%
\pgfsetbuttcap%
\pgfsetroundjoin%
\definecolor{currentfill}{rgb}{0.501961,0.501961,0.501961}%
\pgfsetfillcolor{currentfill}%
\pgfsetfillopacity{0.150000}%
\pgfsetlinewidth{1.003750pt}%
\definecolor{currentstroke}{rgb}{0.501961,0.501961,0.501961}%
\pgfsetstrokecolor{currentstroke}%
\pgfsetstrokeopacity{0.150000}%
\pgfsetdash{}{0pt}%
\pgfsys@defobject{currentmarker}{\pgfqpoint{3.659940in}{0.783058in}}{\pgfqpoint{6.100000in}{4.096407in}}{%
\pgfpathmoveto{\pgfqpoint{3.659940in}{2.678524in}}%
\pgfpathlineto{\pgfqpoint{3.659940in}{0.783058in}}%
\pgfpathlineto{\pgfqpoint{3.672202in}{0.799708in}}%
\pgfpathlineto{\pgfqpoint{3.684463in}{0.816358in}}%
\pgfpathlineto{\pgfqpoint{3.696725in}{0.833008in}}%
\pgfpathlineto{\pgfqpoint{3.708986in}{0.849658in}}%
\pgfpathlineto{\pgfqpoint{3.721248in}{0.866308in}}%
\pgfpathlineto{\pgfqpoint{3.733510in}{0.882958in}}%
\pgfpathlineto{\pgfqpoint{3.745771in}{0.899608in}}%
\pgfpathlineto{\pgfqpoint{3.758033in}{0.916258in}}%
\pgfpathlineto{\pgfqpoint{3.770294in}{0.932908in}}%
\pgfpathlineto{\pgfqpoint{3.782556in}{0.949558in}}%
\pgfpathlineto{\pgfqpoint{3.794818in}{0.966208in}}%
\pgfpathlineto{\pgfqpoint{3.807079in}{0.982858in}}%
\pgfpathlineto{\pgfqpoint{3.819341in}{0.999508in}}%
\pgfpathlineto{\pgfqpoint{3.831603in}{1.016158in}}%
\pgfpathlineto{\pgfqpoint{3.843864in}{1.032808in}}%
\pgfpathlineto{\pgfqpoint{3.856126in}{1.049458in}}%
\pgfpathlineto{\pgfqpoint{3.868387in}{1.066108in}}%
\pgfpathlineto{\pgfqpoint{3.880649in}{1.082758in}}%
\pgfpathlineto{\pgfqpoint{3.892911in}{1.099408in}}%
\pgfpathlineto{\pgfqpoint{3.905172in}{1.116058in}}%
\pgfpathlineto{\pgfqpoint{3.917434in}{1.132708in}}%
\pgfpathlineto{\pgfqpoint{3.929695in}{1.149358in}}%
\pgfpathlineto{\pgfqpoint{3.941957in}{1.166008in}}%
\pgfpathlineto{\pgfqpoint{3.954219in}{1.182658in}}%
\pgfpathlineto{\pgfqpoint{3.966480in}{1.199308in}}%
\pgfpathlineto{\pgfqpoint{3.978742in}{1.215958in}}%
\pgfpathlineto{\pgfqpoint{3.991003in}{1.232608in}}%
\pgfpathlineto{\pgfqpoint{4.003265in}{1.249258in}}%
\pgfpathlineto{\pgfqpoint{4.015527in}{1.265908in}}%
\pgfpathlineto{\pgfqpoint{4.027788in}{1.282558in}}%
\pgfpathlineto{\pgfqpoint{4.040050in}{1.299208in}}%
\pgfpathlineto{\pgfqpoint{4.052311in}{1.315858in}}%
\pgfpathlineto{\pgfqpoint{4.064573in}{1.332508in}}%
\pgfpathlineto{\pgfqpoint{4.076835in}{1.349158in}}%
\pgfpathlineto{\pgfqpoint{4.089096in}{1.365808in}}%
\pgfpathlineto{\pgfqpoint{4.101358in}{1.382458in}}%
\pgfpathlineto{\pgfqpoint{4.113619in}{1.399108in}}%
\pgfpathlineto{\pgfqpoint{4.125881in}{1.415758in}}%
\pgfpathlineto{\pgfqpoint{4.138143in}{1.432408in}}%
\pgfpathlineto{\pgfqpoint{4.150404in}{1.449058in}}%
\pgfpathlineto{\pgfqpoint{4.162666in}{1.465708in}}%
\pgfpathlineto{\pgfqpoint{4.174928in}{1.482358in}}%
\pgfpathlineto{\pgfqpoint{4.187189in}{1.499008in}}%
\pgfpathlineto{\pgfqpoint{4.199451in}{1.515658in}}%
\pgfpathlineto{\pgfqpoint{4.211712in}{1.532308in}}%
\pgfpathlineto{\pgfqpoint{4.223974in}{1.548958in}}%
\pgfpathlineto{\pgfqpoint{4.236236in}{1.565608in}}%
\pgfpathlineto{\pgfqpoint{4.248497in}{1.582258in}}%
\pgfpathlineto{\pgfqpoint{4.260759in}{1.598908in}}%
\pgfpathlineto{\pgfqpoint{4.273020in}{1.615558in}}%
\pgfpathlineto{\pgfqpoint{4.285282in}{1.632208in}}%
\pgfpathlineto{\pgfqpoint{4.297544in}{1.648858in}}%
\pgfpathlineto{\pgfqpoint{4.309805in}{1.665508in}}%
\pgfpathlineto{\pgfqpoint{4.322067in}{1.682158in}}%
\pgfpathlineto{\pgfqpoint{4.334328in}{1.698808in}}%
\pgfpathlineto{\pgfqpoint{4.346590in}{1.715458in}}%
\pgfpathlineto{\pgfqpoint{4.358852in}{1.732108in}}%
\pgfpathlineto{\pgfqpoint{4.371113in}{1.748758in}}%
\pgfpathlineto{\pgfqpoint{4.383375in}{1.765408in}}%
\pgfpathlineto{\pgfqpoint{4.395636in}{1.782058in}}%
\pgfpathlineto{\pgfqpoint{4.407898in}{1.798708in}}%
\pgfpathlineto{\pgfqpoint{4.420160in}{1.815358in}}%
\pgfpathlineto{\pgfqpoint{4.432421in}{1.832008in}}%
\pgfpathlineto{\pgfqpoint{4.444683in}{1.848658in}}%
\pgfpathlineto{\pgfqpoint{4.456945in}{1.865308in}}%
\pgfpathlineto{\pgfqpoint{4.469206in}{1.881958in}}%
\pgfpathlineto{\pgfqpoint{4.481468in}{1.898608in}}%
\pgfpathlineto{\pgfqpoint{4.493729in}{1.915258in}}%
\pgfpathlineto{\pgfqpoint{4.505991in}{1.931908in}}%
\pgfpathlineto{\pgfqpoint{4.518253in}{1.948558in}}%
\pgfpathlineto{\pgfqpoint{4.530514in}{1.965208in}}%
\pgfpathlineto{\pgfqpoint{4.542776in}{1.981858in}}%
\pgfpathlineto{\pgfqpoint{4.555037in}{1.998508in}}%
\pgfpathlineto{\pgfqpoint{4.567299in}{2.015158in}}%
\pgfpathlineto{\pgfqpoint{4.579561in}{2.031808in}}%
\pgfpathlineto{\pgfqpoint{4.591822in}{2.048458in}}%
\pgfpathlineto{\pgfqpoint{4.604084in}{2.065108in}}%
\pgfpathlineto{\pgfqpoint{4.616345in}{2.081758in}}%
\pgfpathlineto{\pgfqpoint{4.628607in}{2.098408in}}%
\pgfpathlineto{\pgfqpoint{4.640869in}{2.115058in}}%
\pgfpathlineto{\pgfqpoint{4.653130in}{2.131708in}}%
\pgfpathlineto{\pgfqpoint{4.665392in}{2.142146in}}%
\pgfpathlineto{\pgfqpoint{4.677653in}{2.135605in}}%
\pgfpathlineto{\pgfqpoint{4.689915in}{2.129063in}}%
\pgfpathlineto{\pgfqpoint{4.702177in}{2.122522in}}%
\pgfpathlineto{\pgfqpoint{4.714438in}{2.115981in}}%
\pgfpathlineto{\pgfqpoint{4.726700in}{2.109440in}}%
\pgfpathlineto{\pgfqpoint{4.738962in}{2.102899in}}%
\pgfpathlineto{\pgfqpoint{4.751223in}{2.096358in}}%
\pgfpathlineto{\pgfqpoint{4.763485in}{2.089816in}}%
\pgfpathlineto{\pgfqpoint{4.775746in}{2.083275in}}%
\pgfpathlineto{\pgfqpoint{4.788008in}{2.076734in}}%
\pgfpathlineto{\pgfqpoint{4.800270in}{2.070193in}}%
\pgfpathlineto{\pgfqpoint{4.812531in}{2.063652in}}%
\pgfpathlineto{\pgfqpoint{4.824793in}{2.057110in}}%
\pgfpathlineto{\pgfqpoint{4.837054in}{2.050569in}}%
\pgfpathlineto{\pgfqpoint{4.849316in}{2.044028in}}%
\pgfpathlineto{\pgfqpoint{4.861578in}{2.037487in}}%
\pgfpathlineto{\pgfqpoint{4.873839in}{2.030946in}}%
\pgfpathlineto{\pgfqpoint{4.886101in}{2.024404in}}%
\pgfpathlineto{\pgfqpoint{4.898362in}{2.017863in}}%
\pgfpathlineto{\pgfqpoint{4.910624in}{2.011322in}}%
\pgfpathlineto{\pgfqpoint{4.922886in}{2.004781in}}%
\pgfpathlineto{\pgfqpoint{4.935147in}{1.998240in}}%
\pgfpathlineto{\pgfqpoint{4.947409in}{1.991698in}}%
\pgfpathlineto{\pgfqpoint{4.959670in}{1.985157in}}%
\pgfpathlineto{\pgfqpoint{4.971932in}{1.978616in}}%
\pgfpathlineto{\pgfqpoint{4.984194in}{1.972075in}}%
\pgfpathlineto{\pgfqpoint{4.996455in}{1.965534in}}%
\pgfpathlineto{\pgfqpoint{5.008717in}{1.958992in}}%
\pgfpathlineto{\pgfqpoint{5.020978in}{1.952451in}}%
\pgfpathlineto{\pgfqpoint{5.033240in}{1.945910in}}%
\pgfpathlineto{\pgfqpoint{5.045502in}{1.939369in}}%
\pgfpathlineto{\pgfqpoint{5.057763in}{1.932828in}}%
\pgfpathlineto{\pgfqpoint{5.070025in}{1.926286in}}%
\pgfpathlineto{\pgfqpoint{5.082287in}{1.919745in}}%
\pgfpathlineto{\pgfqpoint{5.094548in}{1.913204in}}%
\pgfpathlineto{\pgfqpoint{5.106810in}{1.906663in}}%
\pgfpathlineto{\pgfqpoint{5.119071in}{1.900122in}}%
\pgfpathlineto{\pgfqpoint{5.131333in}{1.893580in}}%
\pgfpathlineto{\pgfqpoint{5.143595in}{1.887039in}}%
\pgfpathlineto{\pgfqpoint{5.155856in}{1.880498in}}%
\pgfpathlineto{\pgfqpoint{5.168118in}{1.873957in}}%
\pgfpathlineto{\pgfqpoint{5.180379in}{1.867416in}}%
\pgfpathlineto{\pgfqpoint{5.192641in}{1.860874in}}%
\pgfpathlineto{\pgfqpoint{5.204903in}{1.854333in}}%
\pgfpathlineto{\pgfqpoint{5.217164in}{1.847792in}}%
\pgfpathlineto{\pgfqpoint{5.229426in}{1.841251in}}%
\pgfpathlineto{\pgfqpoint{5.241687in}{1.834710in}}%
\pgfpathlineto{\pgfqpoint{5.253949in}{1.828168in}}%
\pgfpathlineto{\pgfqpoint{5.266211in}{1.821627in}}%
\pgfpathlineto{\pgfqpoint{5.278472in}{1.815086in}}%
\pgfpathlineto{\pgfqpoint{5.290734in}{1.808545in}}%
\pgfpathlineto{\pgfqpoint{5.302995in}{1.802004in}}%
\pgfpathlineto{\pgfqpoint{5.315257in}{1.795462in}}%
\pgfpathlineto{\pgfqpoint{5.327519in}{1.788921in}}%
\pgfpathlineto{\pgfqpoint{5.339780in}{1.782380in}}%
\pgfpathlineto{\pgfqpoint{5.352042in}{1.775839in}}%
\pgfpathlineto{\pgfqpoint{5.364304in}{1.769298in}}%
\pgfpathlineto{\pgfqpoint{5.376565in}{1.762756in}}%
\pgfpathlineto{\pgfqpoint{5.388827in}{1.756215in}}%
\pgfpathlineto{\pgfqpoint{5.401088in}{1.749674in}}%
\pgfpathlineto{\pgfqpoint{5.413350in}{1.743133in}}%
\pgfpathlineto{\pgfqpoint{5.425612in}{1.736592in}}%
\pgfpathlineto{\pgfqpoint{5.437873in}{1.730050in}}%
\pgfpathlineto{\pgfqpoint{5.450135in}{1.723509in}}%
\pgfpathlineto{\pgfqpoint{5.462396in}{1.716968in}}%
\pgfpathlineto{\pgfqpoint{5.474658in}{1.710427in}}%
\pgfpathlineto{\pgfqpoint{5.486920in}{1.703886in}}%
\pgfpathlineto{\pgfqpoint{5.499181in}{1.697344in}}%
\pgfpathlineto{\pgfqpoint{5.511443in}{1.690803in}}%
\pgfpathlineto{\pgfqpoint{5.523704in}{1.684262in}}%
\pgfpathlineto{\pgfqpoint{5.535966in}{1.677721in}}%
\pgfpathlineto{\pgfqpoint{5.548228in}{1.671180in}}%
\pgfpathlineto{\pgfqpoint{5.560489in}{1.664638in}}%
\pgfpathlineto{\pgfqpoint{5.572751in}{1.658097in}}%
\pgfpathlineto{\pgfqpoint{5.585012in}{1.651556in}}%
\pgfpathlineto{\pgfqpoint{5.597274in}{1.645015in}}%
\pgfpathlineto{\pgfqpoint{5.609536in}{1.638474in}}%
\pgfpathlineto{\pgfqpoint{5.621797in}{1.631932in}}%
\pgfpathlineto{\pgfqpoint{5.634059in}{1.625391in}}%
\pgfpathlineto{\pgfqpoint{5.646321in}{1.618850in}}%
\pgfpathlineto{\pgfqpoint{5.658582in}{1.612309in}}%
\pgfpathlineto{\pgfqpoint{5.670844in}{1.605768in}}%
\pgfpathlineto{\pgfqpoint{5.683105in}{1.599226in}}%
\pgfpathlineto{\pgfqpoint{5.695367in}{1.592685in}}%
\pgfpathlineto{\pgfqpoint{5.707629in}{1.586144in}}%
\pgfpathlineto{\pgfqpoint{5.719890in}{1.579603in}}%
\pgfpathlineto{\pgfqpoint{5.732152in}{1.573062in}}%
\pgfpathlineto{\pgfqpoint{5.744413in}{1.566520in}}%
\pgfpathlineto{\pgfqpoint{5.756675in}{1.559979in}}%
\pgfpathlineto{\pgfqpoint{5.768937in}{1.553438in}}%
\pgfpathlineto{\pgfqpoint{5.781198in}{1.546897in}}%
\pgfpathlineto{\pgfqpoint{5.793460in}{1.540356in}}%
\pgfpathlineto{\pgfqpoint{5.805721in}{1.533815in}}%
\pgfpathlineto{\pgfqpoint{5.817983in}{1.527273in}}%
\pgfpathlineto{\pgfqpoint{5.830245in}{1.520732in}}%
\pgfpathlineto{\pgfqpoint{5.842506in}{1.514191in}}%
\pgfpathlineto{\pgfqpoint{5.854768in}{1.507650in}}%
\pgfpathlineto{\pgfqpoint{5.867029in}{1.501109in}}%
\pgfpathlineto{\pgfqpoint{5.879291in}{1.494567in}}%
\pgfpathlineto{\pgfqpoint{5.891553in}{1.488026in}}%
\pgfpathlineto{\pgfqpoint{5.903814in}{1.481485in}}%
\pgfpathlineto{\pgfqpoint{5.916076in}{1.474944in}}%
\pgfpathlineto{\pgfqpoint{5.928337in}{1.468403in}}%
\pgfpathlineto{\pgfqpoint{5.940599in}{1.461861in}}%
\pgfpathlineto{\pgfqpoint{5.952861in}{1.455320in}}%
\pgfpathlineto{\pgfqpoint{5.965122in}{1.448779in}}%
\pgfpathlineto{\pgfqpoint{5.977384in}{1.442238in}}%
\pgfpathlineto{\pgfqpoint{5.989646in}{1.435697in}}%
\pgfpathlineto{\pgfqpoint{6.001907in}{1.429155in}}%
\pgfpathlineto{\pgfqpoint{6.014169in}{1.422614in}}%
\pgfpathlineto{\pgfqpoint{6.026430in}{1.416073in}}%
\pgfpathlineto{\pgfqpoint{6.038692in}{1.409532in}}%
\pgfpathlineto{\pgfqpoint{6.050954in}{1.402991in}}%
\pgfpathlineto{\pgfqpoint{6.063215in}{1.396449in}}%
\pgfpathlineto{\pgfqpoint{6.075477in}{1.389908in}}%
\pgfpathlineto{\pgfqpoint{6.087738in}{1.383367in}}%
\pgfpathlineto{\pgfqpoint{6.100000in}{1.376826in}}%
\pgfpathlineto{\pgfqpoint{6.100000in}{4.096407in}}%
\pgfpathlineto{\pgfqpoint{6.100000in}{4.096407in}}%
\pgfpathlineto{\pgfqpoint{6.087738in}{4.079757in}}%
\pgfpathlineto{\pgfqpoint{6.075477in}{4.063107in}}%
\pgfpathlineto{\pgfqpoint{6.063215in}{4.046457in}}%
\pgfpathlineto{\pgfqpoint{6.050954in}{4.029807in}}%
\pgfpathlineto{\pgfqpoint{6.038692in}{4.013157in}}%
\pgfpathlineto{\pgfqpoint{6.026430in}{3.996507in}}%
\pgfpathlineto{\pgfqpoint{6.014169in}{3.979857in}}%
\pgfpathlineto{\pgfqpoint{6.001907in}{3.963207in}}%
\pgfpathlineto{\pgfqpoint{5.989646in}{3.946557in}}%
\pgfpathlineto{\pgfqpoint{5.977384in}{3.929907in}}%
\pgfpathlineto{\pgfqpoint{5.965122in}{3.913257in}}%
\pgfpathlineto{\pgfqpoint{5.952861in}{3.896607in}}%
\pgfpathlineto{\pgfqpoint{5.940599in}{3.879957in}}%
\pgfpathlineto{\pgfqpoint{5.928337in}{3.863307in}}%
\pgfpathlineto{\pgfqpoint{5.916076in}{3.846657in}}%
\pgfpathlineto{\pgfqpoint{5.903814in}{3.830007in}}%
\pgfpathlineto{\pgfqpoint{5.891553in}{3.813357in}}%
\pgfpathlineto{\pgfqpoint{5.879291in}{3.796707in}}%
\pgfpathlineto{\pgfqpoint{5.867029in}{3.780057in}}%
\pgfpathlineto{\pgfqpoint{5.854768in}{3.763407in}}%
\pgfpathlineto{\pgfqpoint{5.842506in}{3.746757in}}%
\pgfpathlineto{\pgfqpoint{5.830245in}{3.730107in}}%
\pgfpathlineto{\pgfqpoint{5.817983in}{3.713457in}}%
\pgfpathlineto{\pgfqpoint{5.805721in}{3.696807in}}%
\pgfpathlineto{\pgfqpoint{5.793460in}{3.680157in}}%
\pgfpathlineto{\pgfqpoint{5.781198in}{3.663507in}}%
\pgfpathlineto{\pgfqpoint{5.768937in}{3.646857in}}%
\pgfpathlineto{\pgfqpoint{5.756675in}{3.630207in}}%
\pgfpathlineto{\pgfqpoint{5.744413in}{3.613557in}}%
\pgfpathlineto{\pgfqpoint{5.732152in}{3.596907in}}%
\pgfpathlineto{\pgfqpoint{5.719890in}{3.580257in}}%
\pgfpathlineto{\pgfqpoint{5.707629in}{3.563607in}}%
\pgfpathlineto{\pgfqpoint{5.695367in}{3.546957in}}%
\pgfpathlineto{\pgfqpoint{5.683105in}{3.530307in}}%
\pgfpathlineto{\pgfqpoint{5.670844in}{3.513657in}}%
\pgfpathlineto{\pgfqpoint{5.658582in}{3.497007in}}%
\pgfpathlineto{\pgfqpoint{5.646321in}{3.480357in}}%
\pgfpathlineto{\pgfqpoint{5.634059in}{3.463707in}}%
\pgfpathlineto{\pgfqpoint{5.621797in}{3.447057in}}%
\pgfpathlineto{\pgfqpoint{5.609536in}{3.430407in}}%
\pgfpathlineto{\pgfqpoint{5.597274in}{3.413757in}}%
\pgfpathlineto{\pgfqpoint{5.585012in}{3.397107in}}%
\pgfpathlineto{\pgfqpoint{5.572751in}{3.380457in}}%
\pgfpathlineto{\pgfqpoint{5.560489in}{3.363807in}}%
\pgfpathlineto{\pgfqpoint{5.548228in}{3.347158in}}%
\pgfpathlineto{\pgfqpoint{5.535966in}{3.330508in}}%
\pgfpathlineto{\pgfqpoint{5.523704in}{3.313858in}}%
\pgfpathlineto{\pgfqpoint{5.511443in}{3.297208in}}%
\pgfpathlineto{\pgfqpoint{5.499181in}{3.280558in}}%
\pgfpathlineto{\pgfqpoint{5.486920in}{3.263908in}}%
\pgfpathlineto{\pgfqpoint{5.474658in}{3.247258in}}%
\pgfpathlineto{\pgfqpoint{5.462396in}{3.230608in}}%
\pgfpathlineto{\pgfqpoint{5.450135in}{3.213958in}}%
\pgfpathlineto{\pgfqpoint{5.437873in}{3.197308in}}%
\pgfpathlineto{\pgfqpoint{5.425612in}{3.180658in}}%
\pgfpathlineto{\pgfqpoint{5.413350in}{3.164008in}}%
\pgfpathlineto{\pgfqpoint{5.401088in}{3.147358in}}%
\pgfpathlineto{\pgfqpoint{5.388827in}{3.130708in}}%
\pgfpathlineto{\pgfqpoint{5.376565in}{3.114058in}}%
\pgfpathlineto{\pgfqpoint{5.364304in}{3.097408in}}%
\pgfpathlineto{\pgfqpoint{5.352042in}{3.080758in}}%
\pgfpathlineto{\pgfqpoint{5.339780in}{3.064108in}}%
\pgfpathlineto{\pgfqpoint{5.327519in}{3.047458in}}%
\pgfpathlineto{\pgfqpoint{5.315257in}{3.030808in}}%
\pgfpathlineto{\pgfqpoint{5.302995in}{3.014158in}}%
\pgfpathlineto{\pgfqpoint{5.290734in}{2.997508in}}%
\pgfpathlineto{\pgfqpoint{5.278472in}{2.980858in}}%
\pgfpathlineto{\pgfqpoint{5.266211in}{2.964208in}}%
\pgfpathlineto{\pgfqpoint{5.253949in}{2.947558in}}%
\pgfpathlineto{\pgfqpoint{5.241687in}{2.930908in}}%
\pgfpathlineto{\pgfqpoint{5.229426in}{2.914258in}}%
\pgfpathlineto{\pgfqpoint{5.217164in}{2.897608in}}%
\pgfpathlineto{\pgfqpoint{5.204903in}{2.880958in}}%
\pgfpathlineto{\pgfqpoint{5.192641in}{2.864308in}}%
\pgfpathlineto{\pgfqpoint{5.180379in}{2.847658in}}%
\pgfpathlineto{\pgfqpoint{5.168118in}{2.831008in}}%
\pgfpathlineto{\pgfqpoint{5.155856in}{2.814358in}}%
\pgfpathlineto{\pgfqpoint{5.143595in}{2.797708in}}%
\pgfpathlineto{\pgfqpoint{5.131333in}{2.781058in}}%
\pgfpathlineto{\pgfqpoint{5.119071in}{2.764408in}}%
\pgfpathlineto{\pgfqpoint{5.106810in}{2.747758in}}%
\pgfpathlineto{\pgfqpoint{5.094548in}{2.731108in}}%
\pgfpathlineto{\pgfqpoint{5.082287in}{2.714458in}}%
\pgfpathlineto{\pgfqpoint{5.070025in}{2.697808in}}%
\pgfpathlineto{\pgfqpoint{5.057763in}{2.681158in}}%
\pgfpathlineto{\pgfqpoint{5.045502in}{2.664508in}}%
\pgfpathlineto{\pgfqpoint{5.033240in}{2.647858in}}%
\pgfpathlineto{\pgfqpoint{5.020978in}{2.631208in}}%
\pgfpathlineto{\pgfqpoint{5.008717in}{2.614558in}}%
\pgfpathlineto{\pgfqpoint{4.996455in}{2.597908in}}%
\pgfpathlineto{\pgfqpoint{4.984194in}{2.581258in}}%
\pgfpathlineto{\pgfqpoint{4.971932in}{2.564608in}}%
\pgfpathlineto{\pgfqpoint{4.959670in}{2.547958in}}%
\pgfpathlineto{\pgfqpoint{4.947409in}{2.531308in}}%
\pgfpathlineto{\pgfqpoint{4.935147in}{2.514658in}}%
\pgfpathlineto{\pgfqpoint{4.922886in}{2.498008in}}%
\pgfpathlineto{\pgfqpoint{4.910624in}{2.481358in}}%
\pgfpathlineto{\pgfqpoint{4.898362in}{2.464708in}}%
\pgfpathlineto{\pgfqpoint{4.886101in}{2.448058in}}%
\pgfpathlineto{\pgfqpoint{4.873839in}{2.431408in}}%
\pgfpathlineto{\pgfqpoint{4.861578in}{2.414758in}}%
\pgfpathlineto{\pgfqpoint{4.849316in}{2.398108in}}%
\pgfpathlineto{\pgfqpoint{4.837054in}{2.381458in}}%
\pgfpathlineto{\pgfqpoint{4.824793in}{2.364808in}}%
\pgfpathlineto{\pgfqpoint{4.812531in}{2.348158in}}%
\pgfpathlineto{\pgfqpoint{4.800270in}{2.331508in}}%
\pgfpathlineto{\pgfqpoint{4.788008in}{2.314858in}}%
\pgfpathlineto{\pgfqpoint{4.775746in}{2.298208in}}%
\pgfpathlineto{\pgfqpoint{4.763485in}{2.281558in}}%
\pgfpathlineto{\pgfqpoint{4.751223in}{2.264908in}}%
\pgfpathlineto{\pgfqpoint{4.738962in}{2.248258in}}%
\pgfpathlineto{\pgfqpoint{4.726700in}{2.231608in}}%
\pgfpathlineto{\pgfqpoint{4.714438in}{2.214958in}}%
\pgfpathlineto{\pgfqpoint{4.702177in}{2.198308in}}%
\pgfpathlineto{\pgfqpoint{4.689915in}{2.181658in}}%
\pgfpathlineto{\pgfqpoint{4.677653in}{2.165008in}}%
\pgfpathlineto{\pgfqpoint{4.665392in}{2.148358in}}%
\pgfpathlineto{\pgfqpoint{4.653130in}{2.148687in}}%
\pgfpathlineto{\pgfqpoint{4.640869in}{2.155228in}}%
\pgfpathlineto{\pgfqpoint{4.628607in}{2.161769in}}%
\pgfpathlineto{\pgfqpoint{4.616345in}{2.168311in}}%
\pgfpathlineto{\pgfqpoint{4.604084in}{2.174852in}}%
\pgfpathlineto{\pgfqpoint{4.591822in}{2.181393in}}%
\pgfpathlineto{\pgfqpoint{4.579561in}{2.187934in}}%
\pgfpathlineto{\pgfqpoint{4.567299in}{2.194475in}}%
\pgfpathlineto{\pgfqpoint{4.555037in}{2.201017in}}%
\pgfpathlineto{\pgfqpoint{4.542776in}{2.207558in}}%
\pgfpathlineto{\pgfqpoint{4.530514in}{2.214099in}}%
\pgfpathlineto{\pgfqpoint{4.518253in}{2.220640in}}%
\pgfpathlineto{\pgfqpoint{4.505991in}{2.227181in}}%
\pgfpathlineto{\pgfqpoint{4.493729in}{2.233723in}}%
\pgfpathlineto{\pgfqpoint{4.481468in}{2.240264in}}%
\pgfpathlineto{\pgfqpoint{4.469206in}{2.246805in}}%
\pgfpathlineto{\pgfqpoint{4.456945in}{2.253346in}}%
\pgfpathlineto{\pgfqpoint{4.444683in}{2.259887in}}%
\pgfpathlineto{\pgfqpoint{4.432421in}{2.266429in}}%
\pgfpathlineto{\pgfqpoint{4.420160in}{2.272970in}}%
\pgfpathlineto{\pgfqpoint{4.407898in}{2.279511in}}%
\pgfpathlineto{\pgfqpoint{4.395636in}{2.286052in}}%
\pgfpathlineto{\pgfqpoint{4.383375in}{2.292593in}}%
\pgfpathlineto{\pgfqpoint{4.371113in}{2.299135in}}%
\pgfpathlineto{\pgfqpoint{4.358852in}{2.305676in}}%
\pgfpathlineto{\pgfqpoint{4.346590in}{2.312217in}}%
\pgfpathlineto{\pgfqpoint{4.334328in}{2.318758in}}%
\pgfpathlineto{\pgfqpoint{4.322067in}{2.325299in}}%
\pgfpathlineto{\pgfqpoint{4.309805in}{2.331841in}}%
\pgfpathlineto{\pgfqpoint{4.297544in}{2.338382in}}%
\pgfpathlineto{\pgfqpoint{4.285282in}{2.344923in}}%
\pgfpathlineto{\pgfqpoint{4.273020in}{2.351464in}}%
\pgfpathlineto{\pgfqpoint{4.260759in}{2.358005in}}%
\pgfpathlineto{\pgfqpoint{4.248497in}{2.364547in}}%
\pgfpathlineto{\pgfqpoint{4.236236in}{2.371088in}}%
\pgfpathlineto{\pgfqpoint{4.223974in}{2.377629in}}%
\pgfpathlineto{\pgfqpoint{4.211712in}{2.384170in}}%
\pgfpathlineto{\pgfqpoint{4.199451in}{2.390711in}}%
\pgfpathlineto{\pgfqpoint{4.187189in}{2.397253in}}%
\pgfpathlineto{\pgfqpoint{4.174928in}{2.403794in}}%
\pgfpathlineto{\pgfqpoint{4.162666in}{2.410335in}}%
\pgfpathlineto{\pgfqpoint{4.150404in}{2.416876in}}%
\pgfpathlineto{\pgfqpoint{4.138143in}{2.423417in}}%
\pgfpathlineto{\pgfqpoint{4.125881in}{2.429959in}}%
\pgfpathlineto{\pgfqpoint{4.113619in}{2.436500in}}%
\pgfpathlineto{\pgfqpoint{4.101358in}{2.443041in}}%
\pgfpathlineto{\pgfqpoint{4.089096in}{2.449582in}}%
\pgfpathlineto{\pgfqpoint{4.076835in}{2.456123in}}%
\pgfpathlineto{\pgfqpoint{4.064573in}{2.462665in}}%
\pgfpathlineto{\pgfqpoint{4.052311in}{2.469206in}}%
\pgfpathlineto{\pgfqpoint{4.040050in}{2.475747in}}%
\pgfpathlineto{\pgfqpoint{4.027788in}{2.482288in}}%
\pgfpathlineto{\pgfqpoint{4.015527in}{2.488829in}}%
\pgfpathlineto{\pgfqpoint{4.003265in}{2.495371in}}%
\pgfpathlineto{\pgfqpoint{3.991003in}{2.501912in}}%
\pgfpathlineto{\pgfqpoint{3.978742in}{2.508453in}}%
\pgfpathlineto{\pgfqpoint{3.966480in}{2.514994in}}%
\pgfpathlineto{\pgfqpoint{3.954219in}{2.521535in}}%
\pgfpathlineto{\pgfqpoint{3.941957in}{2.528077in}}%
\pgfpathlineto{\pgfqpoint{3.929695in}{2.534618in}}%
\pgfpathlineto{\pgfqpoint{3.917434in}{2.541159in}}%
\pgfpathlineto{\pgfqpoint{3.905172in}{2.547700in}}%
\pgfpathlineto{\pgfqpoint{3.892911in}{2.554241in}}%
\pgfpathlineto{\pgfqpoint{3.880649in}{2.560783in}}%
\pgfpathlineto{\pgfqpoint{3.868387in}{2.567324in}}%
\pgfpathlineto{\pgfqpoint{3.856126in}{2.573865in}}%
\pgfpathlineto{\pgfqpoint{3.843864in}{2.580406in}}%
\pgfpathlineto{\pgfqpoint{3.831603in}{2.586947in}}%
\pgfpathlineto{\pgfqpoint{3.819341in}{2.593489in}}%
\pgfpathlineto{\pgfqpoint{3.807079in}{2.600030in}}%
\pgfpathlineto{\pgfqpoint{3.794818in}{2.606571in}}%
\pgfpathlineto{\pgfqpoint{3.782556in}{2.613112in}}%
\pgfpathlineto{\pgfqpoint{3.770294in}{2.619653in}}%
\pgfpathlineto{\pgfqpoint{3.758033in}{2.626195in}}%
\pgfpathlineto{\pgfqpoint{3.745771in}{2.632736in}}%
\pgfpathlineto{\pgfqpoint{3.733510in}{2.639277in}}%
\pgfpathlineto{\pgfqpoint{3.721248in}{2.645818in}}%
\pgfpathlineto{\pgfqpoint{3.708986in}{2.652359in}}%
\pgfpathlineto{\pgfqpoint{3.696725in}{2.658900in}}%
\pgfpathlineto{\pgfqpoint{3.684463in}{2.665442in}}%
\pgfpathlineto{\pgfqpoint{3.672202in}{2.671983in}}%
\pgfpathlineto{\pgfqpoint{3.659940in}{2.678524in}}%
\pgfpathlineto{\pgfqpoint{3.659940in}{2.678524in}}%
\pgfpathclose%
\pgfusepath{stroke,fill}%
}%
\begin{pgfscope}%
\pgfsys@transformshift{0.000000in}{0.000000in}%
\pgfsys@useobject{currentmarker}{}%
\end{pgfscope}%
\end{pgfscope}%
\begin{pgfscope}%
\pgfsetbuttcap%
\pgfsetroundjoin%
\definecolor{currentfill}{rgb}{0.000000,0.000000,0.000000}%
\pgfsetfillcolor{currentfill}%
\pgfsetlinewidth{0.803000pt}%
\definecolor{currentstroke}{rgb}{0.000000,0.000000,0.000000}%
\pgfsetstrokecolor{currentstroke}%
\pgfsetdash{}{0pt}%
\pgfsys@defobject{currentmarker}{\pgfqpoint{0.000000in}{-0.048611in}}{\pgfqpoint{0.000000in}{0.000000in}}{%
\pgfpathmoveto{\pgfqpoint{0.000000in}{0.000000in}}%
\pgfpathlineto{\pgfqpoint{0.000000in}{-0.048611in}}%
\pgfusepath{stroke,fill}%
}%
\begin{pgfscope}%
\pgfsys@transformshift{4.226382in}{1.184369in}%
\pgfsys@useobject{currentmarker}{}%
\end{pgfscope}%
\end{pgfscope}%
\begin{pgfscope}%
\definecolor{textcolor}{rgb}{0.000000,0.000000,0.000000}%
\pgfsetstrokecolor{textcolor}%
\pgfsetfillcolor{textcolor}%
\pgftext[x=4.226382in,y=1.087147in,,top]{\color{textcolor}{\rmfamily\fontsize{10.000000}{12.000000}\selectfont\catcode`\^=\active\def^{\ifmmode\sp\else\^{}\fi}\catcode`\%=\active\def%{\%}$\mathdefault{1.4}$}}%
\end{pgfscope}%
\begin{pgfscope}%
\pgfsetbuttcap%
\pgfsetroundjoin%
\definecolor{currentfill}{rgb}{0.000000,0.000000,0.000000}%
\pgfsetfillcolor{currentfill}%
\pgfsetlinewidth{0.803000pt}%
\definecolor{currentstroke}{rgb}{0.000000,0.000000,0.000000}%
\pgfsetstrokecolor{currentstroke}%
\pgfsetdash{}{0pt}%
\pgfsys@defobject{currentmarker}{\pgfqpoint{0.000000in}{-0.048611in}}{\pgfqpoint{0.000000in}{0.000000in}}{%
\pgfpathmoveto{\pgfqpoint{0.000000in}{0.000000in}}%
\pgfpathlineto{\pgfqpoint{0.000000in}{-0.048611in}}%
\pgfusepath{stroke,fill}%
}%
\begin{pgfscope}%
\pgfsys@transformshift{5.097832in}{1.184369in}%
\pgfsys@useobject{currentmarker}{}%
\end{pgfscope}%
\end{pgfscope}%
\begin{pgfscope}%
\definecolor{textcolor}{rgb}{0.000000,0.000000,0.000000}%
\pgfsetstrokecolor{textcolor}%
\pgfsetfillcolor{textcolor}%
\pgftext[x=5.097832in,y=1.087147in,,top]{\color{textcolor}{\rmfamily\fontsize{10.000000}{12.000000}\selectfont\catcode`\^=\active\def^{\ifmmode\sp\else\^{}\fi}\catcode`\%=\active\def%{\%}$\mathdefault{1.6}$}}%
\end{pgfscope}%
\begin{pgfscope}%
\pgfsetbuttcap%
\pgfsetroundjoin%
\definecolor{currentfill}{rgb}{0.000000,0.000000,0.000000}%
\pgfsetfillcolor{currentfill}%
\pgfsetlinewidth{0.803000pt}%
\definecolor{currentstroke}{rgb}{0.000000,0.000000,0.000000}%
\pgfsetstrokecolor{currentstroke}%
\pgfsetdash{}{0pt}%
\pgfsys@defobject{currentmarker}{\pgfqpoint{0.000000in}{-0.048611in}}{\pgfqpoint{0.000000in}{0.000000in}}{%
\pgfpathmoveto{\pgfqpoint{0.000000in}{0.000000in}}%
\pgfpathlineto{\pgfqpoint{0.000000in}{-0.048611in}}%
\pgfusepath{stroke,fill}%
}%
\begin{pgfscope}%
\pgfsys@transformshift{5.969282in}{1.184369in}%
\pgfsys@useobject{currentmarker}{}%
\end{pgfscope}%
\end{pgfscope}%
\begin{pgfscope}%
\definecolor{textcolor}{rgb}{0.000000,0.000000,0.000000}%
\pgfsetstrokecolor{textcolor}%
\pgfsetfillcolor{textcolor}%
\pgftext[x=5.969282in,y=1.087147in,,top]{\color{textcolor}{\rmfamily\fontsize{10.000000}{12.000000}\selectfont\catcode`\^=\active\def^{\ifmmode\sp\else\^{}\fi}\catcode`\%=\active\def%{\%}$\mathdefault{1.8}$}}%
\end{pgfscope}%
\begin{pgfscope}%
\definecolor{textcolor}{rgb}{0.000000,0.000000,0.000000}%
\pgfsetstrokecolor{textcolor}%
\pgfsetfillcolor{textcolor}%
\pgftext[x=4.879970in,y=0.891940in,,top]{\color{textcolor}{\rmfamily\fontsize{10.000000}{12.000000}\selectfont\catcode`\^=\active\def^{\ifmmode\sp\else\^{}\fi}\catcode`\%=\active\def%{\%}$C_{DA}$ (m$^2$)}}%
\end{pgfscope}%
\begin{pgfscope}%
\pgfsetbuttcap%
\pgfsetroundjoin%
\definecolor{currentfill}{rgb}{0.000000,0.000000,0.000000}%
\pgfsetfillcolor{currentfill}%
\pgfsetlinewidth{0.803000pt}%
\definecolor{currentstroke}{rgb}{0.000000,0.000000,0.000000}%
\pgfsetstrokecolor{currentstroke}%
\pgfsetdash{}{0pt}%
\pgfsys@defobject{currentmarker}{\pgfqpoint{-0.048611in}{0.000000in}}{\pgfqpoint{-0.000000in}{0.000000in}}{%
\pgfpathmoveto{\pgfqpoint{-0.000000in}{0.000000in}}%
\pgfpathlineto{\pgfqpoint{-0.048611in}{0.000000in}}%
\pgfusepath{stroke,fill}%
}%
\begin{pgfscope}%
\pgfsys@transformshift{3.659940in}{1.252992in}%
\pgfsys@useobject{currentmarker}{}%
\end{pgfscope}%
\end{pgfscope}%
\begin{pgfscope}%
\definecolor{textcolor}{rgb}{0.000000,0.000000,0.000000}%
\pgfsetstrokecolor{textcolor}%
\pgfsetfillcolor{textcolor}%
\pgftext[x=3.389107in, y=1.202750in, left, base]{\color{textcolor}{\rmfamily\fontsize{10.000000}{12.000000}\selectfont\catcode`\^=\active\def^{\ifmmode\sp\else\^{}\fi}\catcode`\%=\active\def%{\%}$\mathdefault{2.0}$}}%
\end{pgfscope}%
\begin{pgfscope}%
\pgfsetbuttcap%
\pgfsetroundjoin%
\definecolor{currentfill}{rgb}{0.000000,0.000000,0.000000}%
\pgfsetfillcolor{currentfill}%
\pgfsetlinewidth{0.803000pt}%
\definecolor{currentstroke}{rgb}{0.000000,0.000000,0.000000}%
\pgfsetstrokecolor{currentstroke}%
\pgfsetdash{}{0pt}%
\pgfsys@defobject{currentmarker}{\pgfqpoint{-0.048611in}{0.000000in}}{\pgfqpoint{-0.000000in}{0.000000in}}{%
\pgfpathmoveto{\pgfqpoint{-0.000000in}{0.000000in}}%
\pgfpathlineto{\pgfqpoint{-0.048611in}{0.000000in}}%
\pgfusepath{stroke,fill}%
}%
\begin{pgfscope}%
\pgfsys@transformshift{3.659940in}{1.596110in}%
\pgfsys@useobject{currentmarker}{}%
\end{pgfscope}%
\end{pgfscope}%
\begin{pgfscope}%
\definecolor{textcolor}{rgb}{0.000000,0.000000,0.000000}%
\pgfsetstrokecolor{textcolor}%
\pgfsetfillcolor{textcolor}%
\pgftext[x=3.389107in, y=1.545868in, left, base]{\color{textcolor}{\rmfamily\fontsize{10.000000}{12.000000}\selectfont\catcode`\^=\active\def^{\ifmmode\sp\else\^{}\fi}\catcode`\%=\active\def%{\%}$\mathdefault{2.5}$}}%
\end{pgfscope}%
\begin{pgfscope}%
\pgfsetbuttcap%
\pgfsetroundjoin%
\definecolor{currentfill}{rgb}{0.000000,0.000000,0.000000}%
\pgfsetfillcolor{currentfill}%
\pgfsetlinewidth{0.803000pt}%
\definecolor{currentstroke}{rgb}{0.000000,0.000000,0.000000}%
\pgfsetstrokecolor{currentstroke}%
\pgfsetdash{}{0pt}%
\pgfsys@defobject{currentmarker}{\pgfqpoint{-0.048611in}{0.000000in}}{\pgfqpoint{-0.000000in}{0.000000in}}{%
\pgfpathmoveto{\pgfqpoint{-0.000000in}{0.000000in}}%
\pgfpathlineto{\pgfqpoint{-0.048611in}{0.000000in}}%
\pgfusepath{stroke,fill}%
}%
\begin{pgfscope}%
\pgfsys@transformshift{3.659940in}{1.939228in}%
\pgfsys@useobject{currentmarker}{}%
\end{pgfscope}%
\end{pgfscope}%
\begin{pgfscope}%
\definecolor{textcolor}{rgb}{0.000000,0.000000,0.000000}%
\pgfsetstrokecolor{textcolor}%
\pgfsetfillcolor{textcolor}%
\pgftext[x=3.389107in, y=1.888986in, left, base]{\color{textcolor}{\rmfamily\fontsize{10.000000}{12.000000}\selectfont\catcode`\^=\active\def^{\ifmmode\sp\else\^{}\fi}\catcode`\%=\active\def%{\%}$\mathdefault{3.0}$}}%
\end{pgfscope}%
\begin{pgfscope}%
\pgfsetbuttcap%
\pgfsetroundjoin%
\definecolor{currentfill}{rgb}{0.000000,0.000000,0.000000}%
\pgfsetfillcolor{currentfill}%
\pgfsetlinewidth{0.803000pt}%
\definecolor{currentstroke}{rgb}{0.000000,0.000000,0.000000}%
\pgfsetstrokecolor{currentstroke}%
\pgfsetdash{}{0pt}%
\pgfsys@defobject{currentmarker}{\pgfqpoint{-0.048611in}{0.000000in}}{\pgfqpoint{-0.000000in}{0.000000in}}{%
\pgfpathmoveto{\pgfqpoint{-0.000000in}{0.000000in}}%
\pgfpathlineto{\pgfqpoint{-0.048611in}{0.000000in}}%
\pgfusepath{stroke,fill}%
}%
\begin{pgfscope}%
\pgfsys@transformshift{3.659940in}{2.282346in}%
\pgfsys@useobject{currentmarker}{}%
\end{pgfscope}%
\end{pgfscope}%
\begin{pgfscope}%
\definecolor{textcolor}{rgb}{0.000000,0.000000,0.000000}%
\pgfsetstrokecolor{textcolor}%
\pgfsetfillcolor{textcolor}%
\pgftext[x=3.389107in, y=2.232104in, left, base]{\color{textcolor}{\rmfamily\fontsize{10.000000}{12.000000}\selectfont\catcode`\^=\active\def^{\ifmmode\sp\else\^{}\fi}\catcode`\%=\active\def%{\%}$\mathdefault{3.5}$}}%
\end{pgfscope}%
\begin{pgfscope}%
\pgfsetbuttcap%
\pgfsetroundjoin%
\definecolor{currentfill}{rgb}{0.000000,0.000000,0.000000}%
\pgfsetfillcolor{currentfill}%
\pgfsetlinewidth{0.803000pt}%
\definecolor{currentstroke}{rgb}{0.000000,0.000000,0.000000}%
\pgfsetstrokecolor{currentstroke}%
\pgfsetdash{}{0pt}%
\pgfsys@defobject{currentmarker}{\pgfqpoint{-0.048611in}{0.000000in}}{\pgfqpoint{-0.000000in}{0.000000in}}{%
\pgfpathmoveto{\pgfqpoint{-0.000000in}{0.000000in}}%
\pgfpathlineto{\pgfqpoint{-0.048611in}{0.000000in}}%
\pgfusepath{stroke,fill}%
}%
\begin{pgfscope}%
\pgfsys@transformshift{3.659940in}{2.625464in}%
\pgfsys@useobject{currentmarker}{}%
\end{pgfscope}%
\end{pgfscope}%
\begin{pgfscope}%
\definecolor{textcolor}{rgb}{0.000000,0.000000,0.000000}%
\pgfsetstrokecolor{textcolor}%
\pgfsetfillcolor{textcolor}%
\pgftext[x=3.389107in, y=2.575221in, left, base]{\color{textcolor}{\rmfamily\fontsize{10.000000}{12.000000}\selectfont\catcode`\^=\active\def^{\ifmmode\sp\else\^{}\fi}\catcode`\%=\active\def%{\%}$\mathdefault{4.0}$}}%
\end{pgfscope}%
\begin{pgfscope}%
\pgfsetbuttcap%
\pgfsetroundjoin%
\definecolor{currentfill}{rgb}{0.000000,0.000000,0.000000}%
\pgfsetfillcolor{currentfill}%
\pgfsetlinewidth{0.803000pt}%
\definecolor{currentstroke}{rgb}{0.000000,0.000000,0.000000}%
\pgfsetstrokecolor{currentstroke}%
\pgfsetdash{}{0pt}%
\pgfsys@defobject{currentmarker}{\pgfqpoint{-0.048611in}{0.000000in}}{\pgfqpoint{-0.000000in}{0.000000in}}{%
\pgfpathmoveto{\pgfqpoint{-0.000000in}{0.000000in}}%
\pgfpathlineto{\pgfqpoint{-0.048611in}{0.000000in}}%
\pgfusepath{stroke,fill}%
}%
\begin{pgfscope}%
\pgfsys@transformshift{3.659940in}{2.968582in}%
\pgfsys@useobject{currentmarker}{}%
\end{pgfscope}%
\end{pgfscope}%
\begin{pgfscope}%
\definecolor{textcolor}{rgb}{0.000000,0.000000,0.000000}%
\pgfsetstrokecolor{textcolor}%
\pgfsetfillcolor{textcolor}%
\pgftext[x=3.389107in, y=2.918339in, left, base]{\color{textcolor}{\rmfamily\fontsize{10.000000}{12.000000}\selectfont\catcode`\^=\active\def^{\ifmmode\sp\else\^{}\fi}\catcode`\%=\active\def%{\%}$\mathdefault{4.5}$}}%
\end{pgfscope}%
\begin{pgfscope}%
\pgfsetbuttcap%
\pgfsetroundjoin%
\definecolor{currentfill}{rgb}{0.000000,0.000000,0.000000}%
\pgfsetfillcolor{currentfill}%
\pgfsetlinewidth{0.803000pt}%
\definecolor{currentstroke}{rgb}{0.000000,0.000000,0.000000}%
\pgfsetstrokecolor{currentstroke}%
\pgfsetdash{}{0pt}%
\pgfsys@defobject{currentmarker}{\pgfqpoint{-0.048611in}{0.000000in}}{\pgfqpoint{-0.000000in}{0.000000in}}{%
\pgfpathmoveto{\pgfqpoint{-0.000000in}{0.000000in}}%
\pgfpathlineto{\pgfqpoint{-0.048611in}{0.000000in}}%
\pgfusepath{stroke,fill}%
}%
\begin{pgfscope}%
\pgfsys@transformshift{3.659940in}{3.311700in}%
\pgfsys@useobject{currentmarker}{}%
\end{pgfscope}%
\end{pgfscope}%
\begin{pgfscope}%
\definecolor{textcolor}{rgb}{0.000000,0.000000,0.000000}%
\pgfsetstrokecolor{textcolor}%
\pgfsetfillcolor{textcolor}%
\pgftext[x=3.389107in, y=3.261457in, left, base]{\color{textcolor}{\rmfamily\fontsize{10.000000}{12.000000}\selectfont\catcode`\^=\active\def^{\ifmmode\sp\else\^{}\fi}\catcode`\%=\active\def%{\%}$\mathdefault{5.0}$}}%
\end{pgfscope}%
\begin{pgfscope}%
\definecolor{textcolor}{rgb}{0.000000,0.000000,0.000000}%
\pgfsetstrokecolor{textcolor}%
\pgfsetfillcolor{textcolor}%
\pgftext[x=3.333551in,y=2.350970in,,bottom,rotate=90.000000]{\color{textcolor}{\rmfamily\fontsize{10.000000}{12.000000}\selectfont\catcode`\^=\active\def^{\ifmmode\sp\else\^{}\fi}\catcode`\%=\active\def%{\%}$C_{LA}$ (m$^2$)}}%
\end{pgfscope}%
\begin{pgfscope}%
\pgfpathrectangle{\pgfqpoint{3.659940in}{1.184369in}}{\pgfqpoint{2.440060in}{2.333202in}}%
\pgfusepath{clip}%
\pgfsetbuttcap%
\pgfsetroundjoin%
\pgfsetlinewidth{0.803000pt}%
\definecolor{currentstroke}{rgb}{0.501961,0.501961,0.501961}%
\pgfsetstrokecolor{currentstroke}%
\pgfsetstrokeopacity{0.600000}%
\pgfsetdash{{2.960000pt}{1.280000pt}}{0.000000pt}%
\pgfpathmoveto{\pgfqpoint{3.948114in}{1.174369in}}%
\pgfpathlineto{\pgfqpoint{4.653130in}{2.131708in}}%
\pgfpathlineto{\pgfqpoint{4.665392in}{2.142146in}}%
\pgfpathlineto{\pgfqpoint{6.100000in}{1.376826in}}%
\pgfpathlineto{\pgfqpoint{6.100000in}{1.376826in}}%
\pgfusepath{stroke}%
\end{pgfscope}%
\begin{pgfscope}%
\pgfpathrectangle{\pgfqpoint{3.659940in}{1.184369in}}{\pgfqpoint{2.440060in}{2.333202in}}%
\pgfusepath{clip}%
\pgfsetbuttcap%
\pgfsetroundjoin%
\pgfsetlinewidth{0.803000pt}%
\definecolor{currentstroke}{rgb}{0.501961,0.501961,0.501961}%
\pgfsetstrokecolor{currentstroke}%
\pgfsetstrokeopacity{0.600000}%
\pgfsetdash{{2.960000pt}{1.280000pt}}{0.000000pt}%
\pgfpathmoveto{\pgfqpoint{3.659940in}{2.678524in}}%
\pgfpathlineto{\pgfqpoint{4.653130in}{2.148687in}}%
\pgfpathlineto{\pgfqpoint{4.665392in}{2.148358in}}%
\pgfpathlineto{\pgfqpoint{5.681090in}{3.527571in}}%
\pgfpathlineto{\pgfqpoint{5.681090in}{3.527571in}}%
\pgfusepath{stroke}%
\end{pgfscope}%
\begin{pgfscope}%
\pgfpathrectangle{\pgfqpoint{3.659940in}{1.184369in}}{\pgfqpoint{2.440060in}{2.333202in}}%
\pgfusepath{clip}%
\pgfsetrectcap%
\pgfsetroundjoin%
\pgfsetlinewidth{1.505625pt}%
\definecolor{currentstroke}{rgb}{0.000000,0.000000,0.000000}%
\pgfsetstrokecolor{currentstroke}%
\pgfsetdash{}{0pt}%
\pgfpathmoveto{\pgfqpoint{3.659940in}{1.385788in}}%
\pgfpathlineto{\pgfqpoint{6.100000in}{3.231621in}}%
\pgfpathlineto{\pgfqpoint{6.100000in}{3.231621in}}%
\pgfusepath{stroke}%
\end{pgfscope}%
\begin{pgfscope}%
\pgfpathrectangle{\pgfqpoint{3.659940in}{1.184369in}}{\pgfqpoint{2.440060in}{2.333202in}}%
\pgfusepath{clip}%
\pgfsetbuttcap%
\pgfsetroundjoin%
\pgfsetlinewidth{1.505625pt}%
\definecolor{currentstroke}{rgb}{0.129412,0.400000,0.674510}%
\pgfsetstrokecolor{currentstroke}%
\pgfsetdash{}{0pt}%
\pgfpathmoveto{\pgfqpoint{4.014397in}{1.707280in}}%
\pgfpathlineto{\pgfqpoint{4.856664in}{1.707280in}}%
\pgfusepath{stroke}%
\end{pgfscope}%
\begin{pgfscope}%
\pgfpathrectangle{\pgfqpoint{3.659940in}{1.184369in}}{\pgfqpoint{2.440060in}{2.333202in}}%
\pgfusepath{clip}%
\pgfsetbuttcap%
\pgfsetroundjoin%
\pgfsetlinewidth{1.505625pt}%
\definecolor{currentstroke}{rgb}{0.129412,0.400000,0.674510}%
\pgfsetstrokecolor{currentstroke}%
\pgfsetdash{}{0pt}%
\pgfpathmoveto{\pgfqpoint{4.435530in}{1.418248in}}%
\pgfpathlineto{\pgfqpoint{4.435530in}{1.996312in}}%
\pgfusepath{stroke}%
\end{pgfscope}%
\begin{pgfscope}%
\pgfpathrectangle{\pgfqpoint{3.659940in}{1.184369in}}{\pgfqpoint{2.440060in}{2.333202in}}%
\pgfusepath{clip}%
\pgfsetbuttcap%
\pgfsetroundjoin%
\definecolor{currentfill}{rgb}{0.129412,0.400000,0.674510}%
\pgfsetfillcolor{currentfill}%
\pgfsetlinewidth{1.003750pt}%
\definecolor{currentstroke}{rgb}{0.129412,0.400000,0.674510}%
\pgfsetstrokecolor{currentstroke}%
\pgfsetdash{}{0pt}%
\pgfsys@defobject{currentmarker}{\pgfqpoint{0.000000in}{-0.069444in}}{\pgfqpoint{0.000000in}{0.069444in}}{%
\pgfpathmoveto{\pgfqpoint{0.000000in}{-0.069444in}}%
\pgfpathlineto{\pgfqpoint{0.000000in}{0.069444in}}%
\pgfusepath{stroke,fill}%
}%
\begin{pgfscope}%
\pgfsys@transformshift{4.014397in}{1.707280in}%
\pgfsys@useobject{currentmarker}{}%
\end{pgfscope}%
\end{pgfscope}%
\begin{pgfscope}%
\pgfpathrectangle{\pgfqpoint{3.659940in}{1.184369in}}{\pgfqpoint{2.440060in}{2.333202in}}%
\pgfusepath{clip}%
\pgfsetbuttcap%
\pgfsetroundjoin%
\definecolor{currentfill}{rgb}{0.129412,0.400000,0.674510}%
\pgfsetfillcolor{currentfill}%
\pgfsetlinewidth{1.003750pt}%
\definecolor{currentstroke}{rgb}{0.129412,0.400000,0.674510}%
\pgfsetstrokecolor{currentstroke}%
\pgfsetdash{}{0pt}%
\pgfsys@defobject{currentmarker}{\pgfqpoint{0.000000in}{-0.069444in}}{\pgfqpoint{0.000000in}{0.069444in}}{%
\pgfpathmoveto{\pgfqpoint{0.000000in}{-0.069444in}}%
\pgfpathlineto{\pgfqpoint{0.000000in}{0.069444in}}%
\pgfusepath{stroke,fill}%
}%
\begin{pgfscope}%
\pgfsys@transformshift{4.856664in}{1.707280in}%
\pgfsys@useobject{currentmarker}{}%
\end{pgfscope}%
\end{pgfscope}%
\begin{pgfscope}%
\pgfpathrectangle{\pgfqpoint{3.659940in}{1.184369in}}{\pgfqpoint{2.440060in}{2.333202in}}%
\pgfusepath{clip}%
\pgfsetbuttcap%
\pgfsetroundjoin%
\definecolor{currentfill}{rgb}{0.129412,0.400000,0.674510}%
\pgfsetfillcolor{currentfill}%
\pgfsetlinewidth{1.003750pt}%
\definecolor{currentstroke}{rgb}{0.129412,0.400000,0.674510}%
\pgfsetstrokecolor{currentstroke}%
\pgfsetdash{}{0pt}%
\pgfsys@defobject{currentmarker}{\pgfqpoint{-0.069444in}{-0.000000in}}{\pgfqpoint{0.069444in}{0.000000in}}{%
\pgfpathmoveto{\pgfqpoint{0.069444in}{-0.000000in}}%
\pgfpathlineto{\pgfqpoint{-0.069444in}{0.000000in}}%
\pgfusepath{stroke,fill}%
}%
\begin{pgfscope}%
\pgfsys@transformshift{4.435530in}{1.418248in}%
\pgfsys@useobject{currentmarker}{}%
\end{pgfscope}%
\end{pgfscope}%
\begin{pgfscope}%
\pgfpathrectangle{\pgfqpoint{3.659940in}{1.184369in}}{\pgfqpoint{2.440060in}{2.333202in}}%
\pgfusepath{clip}%
\pgfsetbuttcap%
\pgfsetroundjoin%
\definecolor{currentfill}{rgb}{0.129412,0.400000,0.674510}%
\pgfsetfillcolor{currentfill}%
\pgfsetlinewidth{1.003750pt}%
\definecolor{currentstroke}{rgb}{0.129412,0.400000,0.674510}%
\pgfsetstrokecolor{currentstroke}%
\pgfsetdash{}{0pt}%
\pgfsys@defobject{currentmarker}{\pgfqpoint{-0.069444in}{-0.000000in}}{\pgfqpoint{0.069444in}{0.000000in}}{%
\pgfpathmoveto{\pgfqpoint{0.069444in}{-0.000000in}}%
\pgfpathlineto{\pgfqpoint{-0.069444in}{0.000000in}}%
\pgfusepath{stroke,fill}%
}%
\begin{pgfscope}%
\pgfsys@transformshift{4.435530in}{1.996312in}%
\pgfsys@useobject{currentmarker}{}%
\end{pgfscope}%
\end{pgfscope}%
\begin{pgfscope}%
\pgfpathrectangle{\pgfqpoint{3.659940in}{1.184369in}}{\pgfqpoint{2.440060in}{2.333202in}}%
\pgfusepath{clip}%
\pgfsetbuttcap%
\pgfsetroundjoin%
\pgfsetlinewidth{1.505625pt}%
\definecolor{currentstroke}{rgb}{0.454902,0.678431,0.819608}%
\pgfsetstrokecolor{currentstroke}%
\pgfsetdash{}{0pt}%
\pgfpathmoveto{\pgfqpoint{3.649940in}{1.983147in}}%
\pgfpathlineto{\pgfqpoint{4.845098in}{1.983147in}}%
\pgfusepath{stroke}%
\end{pgfscope}%
\begin{pgfscope}%
\pgfpathrectangle{\pgfqpoint{3.659940in}{1.184369in}}{\pgfqpoint{2.440060in}{2.333202in}}%
\pgfusepath{clip}%
\pgfsetbuttcap%
\pgfsetroundjoin%
\pgfsetlinewidth{1.505625pt}%
\definecolor{currentstroke}{rgb}{0.454902,0.678431,0.819608}%
\pgfsetstrokecolor{currentstroke}%
\pgfsetdash{}{0pt}%
\pgfpathmoveto{\pgfqpoint{4.121808in}{1.655565in}}%
\pgfpathlineto{\pgfqpoint{4.121808in}{2.310729in}}%
\pgfusepath{stroke}%
\end{pgfscope}%
\begin{pgfscope}%
\pgfpathrectangle{\pgfqpoint{3.659940in}{1.184369in}}{\pgfqpoint{2.440060in}{2.333202in}}%
\pgfusepath{clip}%
\pgfsetbuttcap%
\pgfsetroundjoin%
\definecolor{currentfill}{rgb}{0.454902,0.678431,0.819608}%
\pgfsetfillcolor{currentfill}%
\pgfsetlinewidth{1.003750pt}%
\definecolor{currentstroke}{rgb}{0.454902,0.678431,0.819608}%
\pgfsetstrokecolor{currentstroke}%
\pgfsetdash{}{0pt}%
\pgfsys@defobject{currentmarker}{\pgfqpoint{0.000000in}{-0.069444in}}{\pgfqpoint{0.000000in}{0.069444in}}{%
\pgfpathmoveto{\pgfqpoint{0.000000in}{-0.069444in}}%
\pgfpathlineto{\pgfqpoint{0.000000in}{0.069444in}}%
\pgfusepath{stroke,fill}%
}%
\begin{pgfscope}%
\pgfsys@transformshift{3.398519in}{1.983147in}%
\pgfsys@useobject{currentmarker}{}%
\end{pgfscope}%
\end{pgfscope}%
\begin{pgfscope}%
\pgfpathrectangle{\pgfqpoint{3.659940in}{1.184369in}}{\pgfqpoint{2.440060in}{2.333202in}}%
\pgfusepath{clip}%
\pgfsetbuttcap%
\pgfsetroundjoin%
\definecolor{currentfill}{rgb}{0.454902,0.678431,0.819608}%
\pgfsetfillcolor{currentfill}%
\pgfsetlinewidth{1.003750pt}%
\definecolor{currentstroke}{rgb}{0.454902,0.678431,0.819608}%
\pgfsetstrokecolor{currentstroke}%
\pgfsetdash{}{0pt}%
\pgfsys@defobject{currentmarker}{\pgfqpoint{0.000000in}{-0.069444in}}{\pgfqpoint{0.000000in}{0.069444in}}{%
\pgfpathmoveto{\pgfqpoint{0.000000in}{-0.069444in}}%
\pgfpathlineto{\pgfqpoint{0.000000in}{0.069444in}}%
\pgfusepath{stroke,fill}%
}%
\begin{pgfscope}%
\pgfsys@transformshift{4.845098in}{1.983147in}%
\pgfsys@useobject{currentmarker}{}%
\end{pgfscope}%
\end{pgfscope}%
\begin{pgfscope}%
\pgfpathrectangle{\pgfqpoint{3.659940in}{1.184369in}}{\pgfqpoint{2.440060in}{2.333202in}}%
\pgfusepath{clip}%
\pgfsetbuttcap%
\pgfsetroundjoin%
\definecolor{currentfill}{rgb}{0.454902,0.678431,0.819608}%
\pgfsetfillcolor{currentfill}%
\pgfsetlinewidth{1.003750pt}%
\definecolor{currentstroke}{rgb}{0.454902,0.678431,0.819608}%
\pgfsetstrokecolor{currentstroke}%
\pgfsetdash{}{0pt}%
\pgfsys@defobject{currentmarker}{\pgfqpoint{-0.069444in}{-0.000000in}}{\pgfqpoint{0.069444in}{0.000000in}}{%
\pgfpathmoveto{\pgfqpoint{0.069444in}{-0.000000in}}%
\pgfpathlineto{\pgfqpoint{-0.069444in}{0.000000in}}%
\pgfusepath{stroke,fill}%
}%
\begin{pgfscope}%
\pgfsys@transformshift{4.121808in}{1.655565in}%
\pgfsys@useobject{currentmarker}{}%
\end{pgfscope}%
\end{pgfscope}%
\begin{pgfscope}%
\pgfpathrectangle{\pgfqpoint{3.659940in}{1.184369in}}{\pgfqpoint{2.440060in}{2.333202in}}%
\pgfusepath{clip}%
\pgfsetbuttcap%
\pgfsetroundjoin%
\definecolor{currentfill}{rgb}{0.454902,0.678431,0.819608}%
\pgfsetfillcolor{currentfill}%
\pgfsetlinewidth{1.003750pt}%
\definecolor{currentstroke}{rgb}{0.454902,0.678431,0.819608}%
\pgfsetstrokecolor{currentstroke}%
\pgfsetdash{}{0pt}%
\pgfsys@defobject{currentmarker}{\pgfqpoint{-0.069444in}{-0.000000in}}{\pgfqpoint{0.069444in}{0.000000in}}{%
\pgfpathmoveto{\pgfqpoint{0.069444in}{-0.000000in}}%
\pgfpathlineto{\pgfqpoint{-0.069444in}{0.000000in}}%
\pgfusepath{stroke,fill}%
}%
\begin{pgfscope}%
\pgfsys@transformshift{4.121808in}{2.310729in}%
\pgfsys@useobject{currentmarker}{}%
\end{pgfscope}%
\end{pgfscope}%
\begin{pgfscope}%
\pgfpathrectangle{\pgfqpoint{3.659940in}{1.184369in}}{\pgfqpoint{2.440060in}{2.333202in}}%
\pgfusepath{clip}%
\pgfsetbuttcap%
\pgfsetroundjoin%
\pgfsetlinewidth{1.505625pt}%
\definecolor{currentstroke}{rgb}{0.992157,0.682353,0.380392}%
\pgfsetstrokecolor{currentstroke}%
\pgfsetdash{}{0pt}%
\pgfpathmoveto{\pgfqpoint{4.188921in}{2.033242in}}%
\pgfpathlineto{\pgfqpoint{5.022005in}{2.033242in}}%
\pgfusepath{stroke}%
\end{pgfscope}%
\begin{pgfscope}%
\pgfpathrectangle{\pgfqpoint{3.659940in}{1.184369in}}{\pgfqpoint{2.440060in}{2.333202in}}%
\pgfusepath{clip}%
\pgfsetbuttcap%
\pgfsetroundjoin%
\pgfsetlinewidth{1.505625pt}%
\definecolor{currentstroke}{rgb}{0.992157,0.682353,0.380392}%
\pgfsetstrokecolor{currentstroke}%
\pgfsetdash{}{0pt}%
\pgfpathmoveto{\pgfqpoint{4.605463in}{1.684488in}}%
\pgfpathlineto{\pgfqpoint{4.605463in}{2.381997in}}%
\pgfusepath{stroke}%
\end{pgfscope}%
\begin{pgfscope}%
\pgfpathrectangle{\pgfqpoint{3.659940in}{1.184369in}}{\pgfqpoint{2.440060in}{2.333202in}}%
\pgfusepath{clip}%
\pgfsetbuttcap%
\pgfsetroundjoin%
\definecolor{currentfill}{rgb}{0.992157,0.682353,0.380392}%
\pgfsetfillcolor{currentfill}%
\pgfsetlinewidth{1.003750pt}%
\definecolor{currentstroke}{rgb}{0.992157,0.682353,0.380392}%
\pgfsetstrokecolor{currentstroke}%
\pgfsetdash{}{0pt}%
\pgfsys@defobject{currentmarker}{\pgfqpoint{0.000000in}{-0.069444in}}{\pgfqpoint{0.000000in}{0.069444in}}{%
\pgfpathmoveto{\pgfqpoint{0.000000in}{-0.069444in}}%
\pgfpathlineto{\pgfqpoint{0.000000in}{0.069444in}}%
\pgfusepath{stroke,fill}%
}%
\begin{pgfscope}%
\pgfsys@transformshift{4.188921in}{2.033242in}%
\pgfsys@useobject{currentmarker}{}%
\end{pgfscope}%
\end{pgfscope}%
\begin{pgfscope}%
\pgfpathrectangle{\pgfqpoint{3.659940in}{1.184369in}}{\pgfqpoint{2.440060in}{2.333202in}}%
\pgfusepath{clip}%
\pgfsetbuttcap%
\pgfsetroundjoin%
\definecolor{currentfill}{rgb}{0.992157,0.682353,0.380392}%
\pgfsetfillcolor{currentfill}%
\pgfsetlinewidth{1.003750pt}%
\definecolor{currentstroke}{rgb}{0.992157,0.682353,0.380392}%
\pgfsetstrokecolor{currentstroke}%
\pgfsetdash{}{0pt}%
\pgfsys@defobject{currentmarker}{\pgfqpoint{0.000000in}{-0.069444in}}{\pgfqpoint{0.000000in}{0.069444in}}{%
\pgfpathmoveto{\pgfqpoint{0.000000in}{-0.069444in}}%
\pgfpathlineto{\pgfqpoint{0.000000in}{0.069444in}}%
\pgfusepath{stroke,fill}%
}%
\begin{pgfscope}%
\pgfsys@transformshift{5.022005in}{2.033242in}%
\pgfsys@useobject{currentmarker}{}%
\end{pgfscope}%
\end{pgfscope}%
\begin{pgfscope}%
\pgfpathrectangle{\pgfqpoint{3.659940in}{1.184369in}}{\pgfqpoint{2.440060in}{2.333202in}}%
\pgfusepath{clip}%
\pgfsetbuttcap%
\pgfsetroundjoin%
\definecolor{currentfill}{rgb}{0.992157,0.682353,0.380392}%
\pgfsetfillcolor{currentfill}%
\pgfsetlinewidth{1.003750pt}%
\definecolor{currentstroke}{rgb}{0.992157,0.682353,0.380392}%
\pgfsetstrokecolor{currentstroke}%
\pgfsetdash{}{0pt}%
\pgfsys@defobject{currentmarker}{\pgfqpoint{-0.069444in}{-0.000000in}}{\pgfqpoint{0.069444in}{0.000000in}}{%
\pgfpathmoveto{\pgfqpoint{0.069444in}{-0.000000in}}%
\pgfpathlineto{\pgfqpoint{-0.069444in}{0.000000in}}%
\pgfusepath{stroke,fill}%
}%
\begin{pgfscope}%
\pgfsys@transformshift{4.605463in}{1.684488in}%
\pgfsys@useobject{currentmarker}{}%
\end{pgfscope}%
\end{pgfscope}%
\begin{pgfscope}%
\pgfpathrectangle{\pgfqpoint{3.659940in}{1.184369in}}{\pgfqpoint{2.440060in}{2.333202in}}%
\pgfusepath{clip}%
\pgfsetbuttcap%
\pgfsetroundjoin%
\definecolor{currentfill}{rgb}{0.992157,0.682353,0.380392}%
\pgfsetfillcolor{currentfill}%
\pgfsetlinewidth{1.003750pt}%
\definecolor{currentstroke}{rgb}{0.992157,0.682353,0.380392}%
\pgfsetstrokecolor{currentstroke}%
\pgfsetdash{}{0pt}%
\pgfsys@defobject{currentmarker}{\pgfqpoint{-0.069444in}{-0.000000in}}{\pgfqpoint{0.069444in}{0.000000in}}{%
\pgfpathmoveto{\pgfqpoint{0.069444in}{-0.000000in}}%
\pgfpathlineto{\pgfqpoint{-0.069444in}{0.000000in}}%
\pgfusepath{stroke,fill}%
}%
\begin{pgfscope}%
\pgfsys@transformshift{4.605463in}{2.381997in}%
\pgfsys@useobject{currentmarker}{}%
\end{pgfscope}%
\end{pgfscope}%
\begin{pgfscope}%
\pgfpathrectangle{\pgfqpoint{3.659940in}{1.184369in}}{\pgfqpoint{2.440060in}{2.333202in}}%
\pgfusepath{clip}%
\pgfsetbuttcap%
\pgfsetroundjoin%
\pgfsetlinewidth{1.505625pt}%
\definecolor{currentstroke}{rgb}{0.843137,0.188235,0.152941}%
\pgfsetstrokecolor{currentstroke}%
\pgfsetdash{}{0pt}%
\pgfpathmoveto{\pgfqpoint{4.673585in}{2.851922in}}%
\pgfpathlineto{\pgfqpoint{6.110000in}{2.851922in}}%
\pgfusepath{stroke}%
\end{pgfscope}%
\begin{pgfscope}%
\pgfpathrectangle{\pgfqpoint{3.659940in}{1.184369in}}{\pgfqpoint{2.440060in}{2.333202in}}%
\pgfusepath{clip}%
\pgfsetbuttcap%
\pgfsetroundjoin%
\pgfsetlinewidth{1.505625pt}%
\definecolor{currentstroke}{rgb}{0.843137,0.188235,0.152941}%
\pgfsetstrokecolor{currentstroke}%
\pgfsetdash{}{0pt}%
\pgfpathmoveto{\pgfqpoint{5.485628in}{2.392887in}}%
\pgfpathlineto{\pgfqpoint{5.485628in}{3.310957in}}%
\pgfusepath{stroke}%
\end{pgfscope}%
\begin{pgfscope}%
\pgfpathrectangle{\pgfqpoint{3.659940in}{1.184369in}}{\pgfqpoint{2.440060in}{2.333202in}}%
\pgfusepath{clip}%
\pgfsetbuttcap%
\pgfsetroundjoin%
\definecolor{currentfill}{rgb}{0.843137,0.188235,0.152941}%
\pgfsetfillcolor{currentfill}%
\pgfsetlinewidth{1.003750pt}%
\definecolor{currentstroke}{rgb}{0.843137,0.188235,0.152941}%
\pgfsetstrokecolor{currentstroke}%
\pgfsetdash{}{0pt}%
\pgfsys@defobject{currentmarker}{\pgfqpoint{0.000000in}{-0.069444in}}{\pgfqpoint{0.000000in}{0.069444in}}{%
\pgfpathmoveto{\pgfqpoint{0.000000in}{-0.069444in}}%
\pgfpathlineto{\pgfqpoint{0.000000in}{0.069444in}}%
\pgfusepath{stroke,fill}%
}%
\begin{pgfscope}%
\pgfsys@transformshift{4.673585in}{2.851922in}%
\pgfsys@useobject{currentmarker}{}%
\end{pgfscope}%
\end{pgfscope}%
\begin{pgfscope}%
\pgfpathrectangle{\pgfqpoint{3.659940in}{1.184369in}}{\pgfqpoint{2.440060in}{2.333202in}}%
\pgfusepath{clip}%
\pgfsetbuttcap%
\pgfsetroundjoin%
\definecolor{currentfill}{rgb}{0.843137,0.188235,0.152941}%
\pgfsetfillcolor{currentfill}%
\pgfsetlinewidth{1.003750pt}%
\definecolor{currentstroke}{rgb}{0.843137,0.188235,0.152941}%
\pgfsetstrokecolor{currentstroke}%
\pgfsetdash{}{0pt}%
\pgfsys@defobject{currentmarker}{\pgfqpoint{0.000000in}{-0.069444in}}{\pgfqpoint{0.000000in}{0.069444in}}{%
\pgfpathmoveto{\pgfqpoint{0.000000in}{-0.069444in}}%
\pgfpathlineto{\pgfqpoint{0.000000in}{0.069444in}}%
\pgfusepath{stroke,fill}%
}%
\begin{pgfscope}%
\pgfsys@transformshift{6.297670in}{2.851922in}%
\pgfsys@useobject{currentmarker}{}%
\end{pgfscope}%
\end{pgfscope}%
\begin{pgfscope}%
\pgfpathrectangle{\pgfqpoint{3.659940in}{1.184369in}}{\pgfqpoint{2.440060in}{2.333202in}}%
\pgfusepath{clip}%
\pgfsetbuttcap%
\pgfsetroundjoin%
\definecolor{currentfill}{rgb}{0.843137,0.188235,0.152941}%
\pgfsetfillcolor{currentfill}%
\pgfsetlinewidth{1.003750pt}%
\definecolor{currentstroke}{rgb}{0.843137,0.188235,0.152941}%
\pgfsetstrokecolor{currentstroke}%
\pgfsetdash{}{0pt}%
\pgfsys@defobject{currentmarker}{\pgfqpoint{-0.069444in}{-0.000000in}}{\pgfqpoint{0.069444in}{0.000000in}}{%
\pgfpathmoveto{\pgfqpoint{0.069444in}{-0.000000in}}%
\pgfpathlineto{\pgfqpoint{-0.069444in}{0.000000in}}%
\pgfusepath{stroke,fill}%
}%
\begin{pgfscope}%
\pgfsys@transformshift{5.485628in}{2.392887in}%
\pgfsys@useobject{currentmarker}{}%
\end{pgfscope}%
\end{pgfscope}%
\begin{pgfscope}%
\pgfpathrectangle{\pgfqpoint{3.659940in}{1.184369in}}{\pgfqpoint{2.440060in}{2.333202in}}%
\pgfusepath{clip}%
\pgfsetbuttcap%
\pgfsetroundjoin%
\definecolor{currentfill}{rgb}{0.843137,0.188235,0.152941}%
\pgfsetfillcolor{currentfill}%
\pgfsetlinewidth{1.003750pt}%
\definecolor{currentstroke}{rgb}{0.843137,0.188235,0.152941}%
\pgfsetstrokecolor{currentstroke}%
\pgfsetdash{}{0pt}%
\pgfsys@defobject{currentmarker}{\pgfqpoint{-0.069444in}{-0.000000in}}{\pgfqpoint{0.069444in}{0.000000in}}{%
\pgfpathmoveto{\pgfqpoint{0.069444in}{-0.000000in}}%
\pgfpathlineto{\pgfqpoint{-0.069444in}{0.000000in}}%
\pgfusepath{stroke,fill}%
}%
\begin{pgfscope}%
\pgfsys@transformshift{5.485628in}{3.310957in}%
\pgfsys@useobject{currentmarker}{}%
\end{pgfscope}%
\end{pgfscope}%
\begin{pgfscope}%
\pgfpathrectangle{\pgfqpoint{3.659940in}{1.184369in}}{\pgfqpoint{2.440060in}{2.333202in}}%
\pgfusepath{clip}%
\pgfsetbuttcap%
\pgfsetroundjoin%
\definecolor{currentfill}{rgb}{0.129412,0.400000,0.674510}%
\pgfsetfillcolor{currentfill}%
\pgfsetlinewidth{1.003750pt}%
\definecolor{currentstroke}{rgb}{0.129412,0.400000,0.674510}%
\pgfsetstrokecolor{currentstroke}%
\pgfsetdash{}{0pt}%
\pgfsys@defobject{currentmarker}{\pgfqpoint{-0.055556in}{-0.055556in}}{\pgfqpoint{0.055556in}{0.055556in}}{%
\pgfpathmoveto{\pgfqpoint{0.000000in}{-0.055556in}}%
\pgfpathcurveto{\pgfqpoint{0.014734in}{-0.055556in}}{\pgfqpoint{0.028866in}{-0.049702in}}{\pgfqpoint{0.039284in}{-0.039284in}}%
\pgfpathcurveto{\pgfqpoint{0.049702in}{-0.028866in}}{\pgfqpoint{0.055556in}{-0.014734in}}{\pgfqpoint{0.055556in}{0.000000in}}%
\pgfpathcurveto{\pgfqpoint{0.055556in}{0.014734in}}{\pgfqpoint{0.049702in}{0.028866in}}{\pgfqpoint{0.039284in}{0.039284in}}%
\pgfpathcurveto{\pgfqpoint{0.028866in}{0.049702in}}{\pgfqpoint{0.014734in}{0.055556in}}{\pgfqpoint{0.000000in}{0.055556in}}%
\pgfpathcurveto{\pgfqpoint{-0.014734in}{0.055556in}}{\pgfqpoint{-0.028866in}{0.049702in}}{\pgfqpoint{-0.039284in}{0.039284in}}%
\pgfpathcurveto{\pgfqpoint{-0.049702in}{0.028866in}}{\pgfqpoint{-0.055556in}{0.014734in}}{\pgfqpoint{-0.055556in}{0.000000in}}%
\pgfpathcurveto{\pgfqpoint{-0.055556in}{-0.014734in}}{\pgfqpoint{-0.049702in}{-0.028866in}}{\pgfqpoint{-0.039284in}{-0.039284in}}%
\pgfpathcurveto{\pgfqpoint{-0.028866in}{-0.049702in}}{\pgfqpoint{-0.014734in}{-0.055556in}}{\pgfqpoint{0.000000in}{-0.055556in}}%
\pgfpathlineto{\pgfqpoint{0.000000in}{-0.055556in}}%
\pgfpathclose%
\pgfusepath{stroke,fill}%
}%
\begin{pgfscope}%
\pgfsys@transformshift{4.435530in}{1.707280in}%
\pgfsys@useobject{currentmarker}{}%
\end{pgfscope}%
\end{pgfscope}%
\begin{pgfscope}%
\pgfpathrectangle{\pgfqpoint{3.659940in}{1.184369in}}{\pgfqpoint{2.440060in}{2.333202in}}%
\pgfusepath{clip}%
\pgfsetbuttcap%
\pgfsetroundjoin%
\definecolor{currentfill}{rgb}{0.454902,0.678431,0.819608}%
\pgfsetfillcolor{currentfill}%
\pgfsetlinewidth{1.003750pt}%
\definecolor{currentstroke}{rgb}{0.454902,0.678431,0.819608}%
\pgfsetstrokecolor{currentstroke}%
\pgfsetdash{}{0pt}%
\pgfsys@defobject{currentmarker}{\pgfqpoint{-0.055556in}{-0.055556in}}{\pgfqpoint{0.055556in}{0.055556in}}{%
\pgfpathmoveto{\pgfqpoint{0.000000in}{-0.055556in}}%
\pgfpathcurveto{\pgfqpoint{0.014734in}{-0.055556in}}{\pgfqpoint{0.028866in}{-0.049702in}}{\pgfqpoint{0.039284in}{-0.039284in}}%
\pgfpathcurveto{\pgfqpoint{0.049702in}{-0.028866in}}{\pgfqpoint{0.055556in}{-0.014734in}}{\pgfqpoint{0.055556in}{0.000000in}}%
\pgfpathcurveto{\pgfqpoint{0.055556in}{0.014734in}}{\pgfqpoint{0.049702in}{0.028866in}}{\pgfqpoint{0.039284in}{0.039284in}}%
\pgfpathcurveto{\pgfqpoint{0.028866in}{0.049702in}}{\pgfqpoint{0.014734in}{0.055556in}}{\pgfqpoint{0.000000in}{0.055556in}}%
\pgfpathcurveto{\pgfqpoint{-0.014734in}{0.055556in}}{\pgfqpoint{-0.028866in}{0.049702in}}{\pgfqpoint{-0.039284in}{0.039284in}}%
\pgfpathcurveto{\pgfqpoint{-0.049702in}{0.028866in}}{\pgfqpoint{-0.055556in}{0.014734in}}{\pgfqpoint{-0.055556in}{0.000000in}}%
\pgfpathcurveto{\pgfqpoint{-0.055556in}{-0.014734in}}{\pgfqpoint{-0.049702in}{-0.028866in}}{\pgfqpoint{-0.039284in}{-0.039284in}}%
\pgfpathcurveto{\pgfqpoint{-0.028866in}{-0.049702in}}{\pgfqpoint{-0.014734in}{-0.055556in}}{\pgfqpoint{0.000000in}{-0.055556in}}%
\pgfpathlineto{\pgfqpoint{0.000000in}{-0.055556in}}%
\pgfpathclose%
\pgfusepath{stroke,fill}%
}%
\begin{pgfscope}%
\pgfsys@transformshift{4.121808in}{1.983147in}%
\pgfsys@useobject{currentmarker}{}%
\end{pgfscope}%
\end{pgfscope}%
\begin{pgfscope}%
\pgfpathrectangle{\pgfqpoint{3.659940in}{1.184369in}}{\pgfqpoint{2.440060in}{2.333202in}}%
\pgfusepath{clip}%
\pgfsetbuttcap%
\pgfsetroundjoin%
\definecolor{currentfill}{rgb}{0.992157,0.682353,0.380392}%
\pgfsetfillcolor{currentfill}%
\pgfsetlinewidth{1.003750pt}%
\definecolor{currentstroke}{rgb}{0.992157,0.682353,0.380392}%
\pgfsetstrokecolor{currentstroke}%
\pgfsetdash{}{0pt}%
\pgfsys@defobject{currentmarker}{\pgfqpoint{-0.055556in}{-0.055556in}}{\pgfqpoint{0.055556in}{0.055556in}}{%
\pgfpathmoveto{\pgfqpoint{0.000000in}{-0.055556in}}%
\pgfpathcurveto{\pgfqpoint{0.014734in}{-0.055556in}}{\pgfqpoint{0.028866in}{-0.049702in}}{\pgfqpoint{0.039284in}{-0.039284in}}%
\pgfpathcurveto{\pgfqpoint{0.049702in}{-0.028866in}}{\pgfqpoint{0.055556in}{-0.014734in}}{\pgfqpoint{0.055556in}{0.000000in}}%
\pgfpathcurveto{\pgfqpoint{0.055556in}{0.014734in}}{\pgfqpoint{0.049702in}{0.028866in}}{\pgfqpoint{0.039284in}{0.039284in}}%
\pgfpathcurveto{\pgfqpoint{0.028866in}{0.049702in}}{\pgfqpoint{0.014734in}{0.055556in}}{\pgfqpoint{0.000000in}{0.055556in}}%
\pgfpathcurveto{\pgfqpoint{-0.014734in}{0.055556in}}{\pgfqpoint{-0.028866in}{0.049702in}}{\pgfqpoint{-0.039284in}{0.039284in}}%
\pgfpathcurveto{\pgfqpoint{-0.049702in}{0.028866in}}{\pgfqpoint{-0.055556in}{0.014734in}}{\pgfqpoint{-0.055556in}{0.000000in}}%
\pgfpathcurveto{\pgfqpoint{-0.055556in}{-0.014734in}}{\pgfqpoint{-0.049702in}{-0.028866in}}{\pgfqpoint{-0.039284in}{-0.039284in}}%
\pgfpathcurveto{\pgfqpoint{-0.028866in}{-0.049702in}}{\pgfqpoint{-0.014734in}{-0.055556in}}{\pgfqpoint{0.000000in}{-0.055556in}}%
\pgfpathlineto{\pgfqpoint{0.000000in}{-0.055556in}}%
\pgfpathclose%
\pgfusepath{stroke,fill}%
}%
\begin{pgfscope}%
\pgfsys@transformshift{4.605463in}{2.033242in}%
\pgfsys@useobject{currentmarker}{}%
\end{pgfscope}%
\end{pgfscope}%
\begin{pgfscope}%
\pgfpathrectangle{\pgfqpoint{3.659940in}{1.184369in}}{\pgfqpoint{2.440060in}{2.333202in}}%
\pgfusepath{clip}%
\pgfsetbuttcap%
\pgfsetroundjoin%
\definecolor{currentfill}{rgb}{0.843137,0.188235,0.152941}%
\pgfsetfillcolor{currentfill}%
\pgfsetlinewidth{1.003750pt}%
\definecolor{currentstroke}{rgb}{0.843137,0.188235,0.152941}%
\pgfsetstrokecolor{currentstroke}%
\pgfsetdash{}{0pt}%
\pgfsys@defobject{currentmarker}{\pgfqpoint{-0.055556in}{-0.055556in}}{\pgfqpoint{0.055556in}{0.055556in}}{%
\pgfpathmoveto{\pgfqpoint{0.000000in}{-0.055556in}}%
\pgfpathcurveto{\pgfqpoint{0.014734in}{-0.055556in}}{\pgfqpoint{0.028866in}{-0.049702in}}{\pgfqpoint{0.039284in}{-0.039284in}}%
\pgfpathcurveto{\pgfqpoint{0.049702in}{-0.028866in}}{\pgfqpoint{0.055556in}{-0.014734in}}{\pgfqpoint{0.055556in}{0.000000in}}%
\pgfpathcurveto{\pgfqpoint{0.055556in}{0.014734in}}{\pgfqpoint{0.049702in}{0.028866in}}{\pgfqpoint{0.039284in}{0.039284in}}%
\pgfpathcurveto{\pgfqpoint{0.028866in}{0.049702in}}{\pgfqpoint{0.014734in}{0.055556in}}{\pgfqpoint{0.000000in}{0.055556in}}%
\pgfpathcurveto{\pgfqpoint{-0.014734in}{0.055556in}}{\pgfqpoint{-0.028866in}{0.049702in}}{\pgfqpoint{-0.039284in}{0.039284in}}%
\pgfpathcurveto{\pgfqpoint{-0.049702in}{0.028866in}}{\pgfqpoint{-0.055556in}{0.014734in}}{\pgfqpoint{-0.055556in}{0.000000in}}%
\pgfpathcurveto{\pgfqpoint{-0.055556in}{-0.014734in}}{\pgfqpoint{-0.049702in}{-0.028866in}}{\pgfqpoint{-0.039284in}{-0.039284in}}%
\pgfpathcurveto{\pgfqpoint{-0.028866in}{-0.049702in}}{\pgfqpoint{-0.014734in}{-0.055556in}}{\pgfqpoint{0.000000in}{-0.055556in}}%
\pgfpathlineto{\pgfqpoint{0.000000in}{-0.055556in}}%
\pgfpathclose%
\pgfusepath{stroke,fill}%
}%
\begin{pgfscope}%
\pgfsys@transformshift{5.485628in}{2.851922in}%
\pgfsys@useobject{currentmarker}{}%
\end{pgfscope}%
\end{pgfscope}%
\begin{pgfscope}%
\pgfsetrectcap%
\pgfsetmiterjoin%
\pgfsetlinewidth{0.803000pt}%
\definecolor{currentstroke}{rgb}{0.000000,0.000000,0.000000}%
\pgfsetstrokecolor{currentstroke}%
\pgfsetdash{}{0pt}%
\pgfpathmoveto{\pgfqpoint{3.659940in}{1.184369in}}%
\pgfpathlineto{\pgfqpoint{3.659940in}{3.517571in}}%
\pgfusepath{stroke}%
\end{pgfscope}%
\begin{pgfscope}%
\pgfsetrectcap%
\pgfsetmiterjoin%
\pgfsetlinewidth{0.803000pt}%
\definecolor{currentstroke}{rgb}{0.000000,0.000000,0.000000}%
\pgfsetstrokecolor{currentstroke}%
\pgfsetdash{}{0pt}%
\pgfpathmoveto{\pgfqpoint{6.100000in}{1.184369in}}%
\pgfpathlineto{\pgfqpoint{6.100000in}{3.517571in}}%
\pgfusepath{stroke}%
\end{pgfscope}%
\begin{pgfscope}%
\pgfsetrectcap%
\pgfsetmiterjoin%
\pgfsetlinewidth{0.803000pt}%
\definecolor{currentstroke}{rgb}{0.000000,0.000000,0.000000}%
\pgfsetstrokecolor{currentstroke}%
\pgfsetdash{}{0pt}%
\pgfpathmoveto{\pgfqpoint{3.659940in}{1.184369in}}%
\pgfpathlineto{\pgfqpoint{6.100000in}{1.184369in}}%
\pgfusepath{stroke}%
\end{pgfscope}%
\begin{pgfscope}%
\pgfsetrectcap%
\pgfsetmiterjoin%
\pgfsetlinewidth{0.803000pt}%
\definecolor{currentstroke}{rgb}{0.000000,0.000000,0.000000}%
\pgfsetstrokecolor{currentstroke}%
\pgfsetdash{}{0pt}%
\pgfpathmoveto{\pgfqpoint{3.659940in}{3.517571in}}%
\pgfpathlineto{\pgfqpoint{6.100000in}{3.517571in}}%
\pgfusepath{stroke}%
\end{pgfscope}%
\begin{pgfscope}%
\pgfsetbuttcap%
\pgfsetmiterjoin%
\definecolor{currentfill}{rgb}{1.000000,1.000000,1.000000}%
\pgfsetfillcolor{currentfill}%
\pgfsetfillopacity{0.800000}%
\pgfsetlinewidth{1.003750pt}%
\definecolor{currentstroke}{rgb}{0.000000,0.000000,0.000000}%
\pgfsetstrokecolor{currentstroke}%
\pgfsetstrokeopacity{0.800000}%
\pgfsetdash{}{0pt}%
\pgfpathmoveto{\pgfqpoint{5.147417in}{1.271862in}}%
\pgfpathlineto{\pgfqpoint{6.026798in}{1.271862in}}%
\pgfpathquadraticcurveto{\pgfqpoint{6.055965in}{1.271862in}}{\pgfqpoint{6.055965in}{1.301029in}}%
\pgfpathlineto{\pgfqpoint{6.055965in}{1.626807in}}%
\pgfpathquadraticcurveto{\pgfqpoint{6.055965in}{1.655974in}}{\pgfqpoint{6.026798in}{1.655974in}}%
\pgfpathlineto{\pgfqpoint{5.147417in}{1.655974in}}%
\pgfpathquadraticcurveto{\pgfqpoint{5.118250in}{1.655974in}}{\pgfqpoint{5.118250in}{1.626807in}}%
\pgfpathlineto{\pgfqpoint{5.118250in}{1.301029in}}%
\pgfpathquadraticcurveto{\pgfqpoint{5.118250in}{1.271862in}}{\pgfqpoint{5.147417in}{1.271862in}}%
\pgfpathlineto{\pgfqpoint{5.147417in}{1.271862in}}%
\pgfpathclose%
\pgfusepath{stroke,fill}%
\end{pgfscope}%
\begin{pgfscope}%
\definecolor{textcolor}{rgb}{0.000000,0.000000,0.000000}%
\pgfsetstrokecolor{textcolor}%
\pgfsetfillcolor{textcolor}%
\pgftext[x=5.496575in, y=1.556468in, left, base]{\color{textcolor}{\rmfamily\fontsize{7.000000}{8.400000}\selectfont\catcode`\^=\active\def^{\ifmmode\sp\else\^{}\fi}\catcode`\%=\active\def%{\%}slope $= 4.80$}}%
\end{pgfscope}%
\begin{pgfscope}%
\definecolor{textcolor}{rgb}{0.000000,0.000000,0.000000}%
\pgfsetstrokecolor{textcolor}%
\pgfsetfillcolor{textcolor}%
\pgftext[x=5.147417in, y=1.440268in, left, base]{\color{textcolor}{\rmfamily\fontsize{7.000000}{8.400000}\selectfont\catcode`\^=\active\def^{\ifmmode\sp\else\^{}\fi}\catcode`\%=\active\def%{\%}90\% CI $= [-3.4,\,8.6]$}}%
\end{pgfscope}%
\begin{pgfscope}%
\definecolor{textcolor}{rgb}{0.000000,0.000000,0.000000}%
\pgfsetstrokecolor{textcolor}%
\pgfsetfillcolor{textcolor}%
\pgftext[x=5.341320in, y=1.328445in, left, base]{\color{textcolor}{\rmfamily\fontsize{7.000000}{8.400000}\selectfont\catcode`\^=\active\def^{\ifmmode\sp\else\^{}\fi}\catcode`\%=\active\def%{\%}$1.03\,\sigma$ from zero}}%
\end{pgfscope}%
\begin{pgfscope}%
\definecolor{textcolor}{rgb}{0.000000,0.000000,0.000000}%
\pgfsetstrokecolor{textcolor}%
\pgfsetfillcolor{textcolor}%
\pgftext[x=4.879970in,y=3.600904in,,base]{\color{textcolor}{\rmfamily\fontsize{10.000000}{12.000000}\selectfont\catcode`\^=\active\def^{\ifmmode\sp\else\^{}\fi}\catcode`\%=\active\def%{\%}(b)}}%
\end{pgfscope}%
\begin{pgfscope}%
\pgfsetbuttcap%
\pgfsetmiterjoin%
\definecolor{currentfill}{rgb}{1.000000,1.000000,1.000000}%
\pgfsetfillcolor{currentfill}%
\pgfsetfillopacity{0.800000}%
\pgfsetlinewidth{1.003750pt}%
\definecolor{currentstroke}{rgb}{0.800000,0.800000,0.800000}%
\pgfsetstrokecolor{currentstroke}%
\pgfsetstrokeopacity{0.800000}%
\pgfsetdash{}{0pt}%
\pgfpathmoveto{\pgfqpoint{1.199265in}{0.100000in}}%
\pgfpathlineto{\pgfqpoint{5.000735in}{0.100000in}}%
\pgfpathquadraticcurveto{\pgfqpoint{5.020179in}{0.100000in}}{\pgfqpoint{5.020179in}{0.119444in}}%
\pgfpathlineto{\pgfqpoint{5.020179in}{0.421900in}}%
\pgfpathquadraticcurveto{\pgfqpoint{5.020179in}{0.441344in}}{\pgfqpoint{5.000735in}{0.441344in}}%
\pgfpathlineto{\pgfqpoint{1.199265in}{0.441344in}}%
\pgfpathquadraticcurveto{\pgfqpoint{1.179821in}{0.441344in}}{\pgfqpoint{1.179821in}{0.421900in}}%
\pgfpathlineto{\pgfqpoint{1.179821in}{0.119444in}}%
\pgfpathquadraticcurveto{\pgfqpoint{1.179821in}{0.100000in}}{\pgfqpoint{1.199265in}{0.100000in}}%
\pgfpathlineto{\pgfqpoint{1.199265in}{0.100000in}}%
\pgfpathclose%
\pgfusepath{stroke,fill}%
\end{pgfscope}%
\begin{pgfscope}%
\pgfsetbuttcap%
\pgfsetmiterjoin%
\definecolor{currentfill}{rgb}{0.501961,0.501961,0.501961}%
\pgfsetfillcolor{currentfill}%
\pgfsetfillopacity{0.150000}%
\pgfsetlinewidth{1.003750pt}%
\definecolor{currentstroke}{rgb}{0.501961,0.501961,0.501961}%
\pgfsetstrokecolor{currentstroke}%
\pgfsetstrokeopacity{0.150000}%
\pgfsetdash{}{0pt}%
\pgfpathmoveto{\pgfqpoint{1.228432in}{0.322394in}}%
\pgfpathlineto{\pgfqpoint{1.374265in}{0.322394in}}%
\pgfpathlineto{\pgfqpoint{1.374265in}{0.390449in}}%
\pgfpathlineto{\pgfqpoint{1.228432in}{0.390449in}}%
\pgfpathlineto{\pgfqpoint{1.228432in}{0.322394in}}%
\pgfpathclose%
\pgfusepath{stroke,fill}%
\end{pgfscope}%
\begin{pgfscope}%
\definecolor{textcolor}{rgb}{0.000000,0.000000,0.000000}%
\pgfsetstrokecolor{textcolor}%
\pgfsetfillcolor{textcolor}%
\pgftext[x=1.413154in,y=0.322394in,left,base]{\color{textcolor}{\rmfamily\fontsize{7.000000}{8.400000}\selectfont\catcode`\^=\active\def^{\ifmmode\sp\else\^{}\fi}\catcode`\%=\active\def%{\%}MC 90\% CI on slope}}%
\end{pgfscope}%
\begin{pgfscope}%
\pgfsetrectcap%
\pgfsetroundjoin%
\pgfsetlinewidth{1.505625pt}%
\definecolor{currentstroke}{rgb}{0.000000,0.000000,0.000000}%
\pgfsetstrokecolor{currentstroke}%
\pgfsetdash{}{0pt}%
\pgfpathmoveto{\pgfqpoint{1.228432in}{0.210055in}}%
\pgfpathlineto{\pgfqpoint{1.301348in}{0.210055in}}%
\pgfpathlineto{\pgfqpoint{1.374265in}{0.210055in}}%
\pgfusepath{stroke}%
\end{pgfscope}%
\begin{pgfscope}%
\definecolor{textcolor}{rgb}{0.000000,0.000000,0.000000}%
\pgfsetstrokecolor{textcolor}%
\pgfsetfillcolor{textcolor}%
\pgftext[x=1.413154in,y=0.176027in,left,base]{\color{textcolor}{\rmfamily\fontsize{7.000000}{8.400000}\selectfont\catcode`\^=\active\def^{\ifmmode\sp\else\^{}\fi}\catcode`\%=\active\def%{\%}OLS fit}}%
\end{pgfscope}%
\begin{pgfscope}%
\pgfsetbuttcap%
\pgfsetroundjoin%
\pgfsetlinewidth{1.505625pt}%
\definecolor{currentstroke}{rgb}{0.129412,0.400000,0.674510}%
\pgfsetstrokecolor{currentstroke}%
\pgfsetdash{}{0pt}%
\pgfpathmoveto{\pgfqpoint{2.410646in}{0.356421in}}%
\pgfpathlineto{\pgfqpoint{2.507869in}{0.356421in}}%
\pgfusepath{stroke}%
\end{pgfscope}%
\begin{pgfscope}%
\pgfsetbuttcap%
\pgfsetroundjoin%
\pgfsetlinewidth{1.505625pt}%
\definecolor{currentstroke}{rgb}{0.129412,0.400000,0.674510}%
\pgfsetstrokecolor{currentstroke}%
\pgfsetdash{}{0pt}%
\pgfpathmoveto{\pgfqpoint{2.459258in}{0.307810in}}%
\pgfpathlineto{\pgfqpoint{2.459258in}{0.405032in}}%
\pgfusepath{stroke}%
\end{pgfscope}%
\begin{pgfscope}%
\pgfsetbuttcap%
\pgfsetroundjoin%
\definecolor{currentfill}{rgb}{0.129412,0.400000,0.674510}%
\pgfsetfillcolor{currentfill}%
\pgfsetlinewidth{1.003750pt}%
\definecolor{currentstroke}{rgb}{0.129412,0.400000,0.674510}%
\pgfsetstrokecolor{currentstroke}%
\pgfsetdash{}{0pt}%
\pgfsys@defobject{currentmarker}{\pgfqpoint{0.000000in}{-0.069444in}}{\pgfqpoint{0.000000in}{0.069444in}}{%
\pgfpathmoveto{\pgfqpoint{0.000000in}{-0.069444in}}%
\pgfpathlineto{\pgfqpoint{0.000000in}{0.069444in}}%
\pgfusepath{stroke,fill}%
}%
\begin{pgfscope}%
\pgfsys@transformshift{2.410646in}{0.356421in}%
\pgfsys@useobject{currentmarker}{}%
\end{pgfscope}%
\end{pgfscope}%
\begin{pgfscope}%
\pgfsetbuttcap%
\pgfsetroundjoin%
\definecolor{currentfill}{rgb}{0.129412,0.400000,0.674510}%
\pgfsetfillcolor{currentfill}%
\pgfsetlinewidth{1.003750pt}%
\definecolor{currentstroke}{rgb}{0.129412,0.400000,0.674510}%
\pgfsetstrokecolor{currentstroke}%
\pgfsetdash{}{0pt}%
\pgfsys@defobject{currentmarker}{\pgfqpoint{0.000000in}{-0.069444in}}{\pgfqpoint{0.000000in}{0.069444in}}{%
\pgfpathmoveto{\pgfqpoint{0.000000in}{-0.069444in}}%
\pgfpathlineto{\pgfqpoint{0.000000in}{0.069444in}}%
\pgfusepath{stroke,fill}%
}%
\begin{pgfscope}%
\pgfsys@transformshift{2.507869in}{0.356421in}%
\pgfsys@useobject{currentmarker}{}%
\end{pgfscope}%
\end{pgfscope}%
\begin{pgfscope}%
\pgfsetbuttcap%
\pgfsetroundjoin%
\definecolor{currentfill}{rgb}{0.129412,0.400000,0.674510}%
\pgfsetfillcolor{currentfill}%
\pgfsetlinewidth{1.003750pt}%
\definecolor{currentstroke}{rgb}{0.129412,0.400000,0.674510}%
\pgfsetstrokecolor{currentstroke}%
\pgfsetdash{}{0pt}%
\pgfsys@defobject{currentmarker}{\pgfqpoint{-0.069444in}{-0.000000in}}{\pgfqpoint{0.069444in}{0.000000in}}{%
\pgfpathmoveto{\pgfqpoint{0.069444in}{-0.000000in}}%
\pgfpathlineto{\pgfqpoint{-0.069444in}{0.000000in}}%
\pgfusepath{stroke,fill}%
}%
\begin{pgfscope}%
\pgfsys@transformshift{2.459258in}{0.307810in}%
\pgfsys@useobject{currentmarker}{}%
\end{pgfscope}%
\end{pgfscope}%
\begin{pgfscope}%
\pgfsetbuttcap%
\pgfsetroundjoin%
\definecolor{currentfill}{rgb}{0.129412,0.400000,0.674510}%
\pgfsetfillcolor{currentfill}%
\pgfsetlinewidth{1.003750pt}%
\definecolor{currentstroke}{rgb}{0.129412,0.400000,0.674510}%
\pgfsetstrokecolor{currentstroke}%
\pgfsetdash{}{0pt}%
\pgfsys@defobject{currentmarker}{\pgfqpoint{-0.069444in}{-0.000000in}}{\pgfqpoint{0.069444in}{0.000000in}}{%
\pgfpathmoveto{\pgfqpoint{0.069444in}{-0.000000in}}%
\pgfpathlineto{\pgfqpoint{-0.069444in}{0.000000in}}%
\pgfusepath{stroke,fill}%
}%
\begin{pgfscope}%
\pgfsys@transformshift{2.459258in}{0.405032in}%
\pgfsys@useobject{currentmarker}{}%
\end{pgfscope}%
\end{pgfscope}%
\begin{pgfscope}%
\pgfsetbuttcap%
\pgfsetroundjoin%
\definecolor{currentfill}{rgb}{0.129412,0.400000,0.674510}%
\pgfsetfillcolor{currentfill}%
\pgfsetlinewidth{1.003750pt}%
\definecolor{currentstroke}{rgb}{0.129412,0.400000,0.674510}%
\pgfsetstrokecolor{currentstroke}%
\pgfsetdash{}{0pt}%
\pgfsys@defobject{currentmarker}{\pgfqpoint{-0.055556in}{-0.055556in}}{\pgfqpoint{0.055556in}{0.055556in}}{%
\pgfpathmoveto{\pgfqpoint{0.000000in}{-0.055556in}}%
\pgfpathcurveto{\pgfqpoint{0.014734in}{-0.055556in}}{\pgfqpoint{0.028866in}{-0.049702in}}{\pgfqpoint{0.039284in}{-0.039284in}}%
\pgfpathcurveto{\pgfqpoint{0.049702in}{-0.028866in}}{\pgfqpoint{0.055556in}{-0.014734in}}{\pgfqpoint{0.055556in}{0.000000in}}%
\pgfpathcurveto{\pgfqpoint{0.055556in}{0.014734in}}{\pgfqpoint{0.049702in}{0.028866in}}{\pgfqpoint{0.039284in}{0.039284in}}%
\pgfpathcurveto{\pgfqpoint{0.028866in}{0.049702in}}{\pgfqpoint{0.014734in}{0.055556in}}{\pgfqpoint{0.000000in}{0.055556in}}%
\pgfpathcurveto{\pgfqpoint{-0.014734in}{0.055556in}}{\pgfqpoint{-0.028866in}{0.049702in}}{\pgfqpoint{-0.039284in}{0.039284in}}%
\pgfpathcurveto{\pgfqpoint{-0.049702in}{0.028866in}}{\pgfqpoint{-0.055556in}{0.014734in}}{\pgfqpoint{-0.055556in}{0.000000in}}%
\pgfpathcurveto{\pgfqpoint{-0.055556in}{-0.014734in}}{\pgfqpoint{-0.049702in}{-0.028866in}}{\pgfqpoint{-0.039284in}{-0.039284in}}%
\pgfpathcurveto{\pgfqpoint{-0.028866in}{-0.049702in}}{\pgfqpoint{-0.014734in}{-0.055556in}}{\pgfqpoint{0.000000in}{-0.055556in}}%
\pgfpathlineto{\pgfqpoint{0.000000in}{-0.055556in}}%
\pgfpathclose%
\pgfusepath{stroke,fill}%
}%
\begin{pgfscope}%
\pgfsys@transformshift{2.459258in}{0.356421in}%
\pgfsys@useobject{currentmarker}{}%
\end{pgfscope}%
\end{pgfscope}%
\begin{pgfscope}%
\definecolor{textcolor}{rgb}{0.000000,0.000000,0.000000}%
\pgfsetstrokecolor{textcolor}%
\pgfsetfillcolor{textcolor}%
\pgftext[x=2.571063in,y=0.322394in,left,base]{\color{textcolor}{\rmfamily\fontsize{7.000000}{8.400000}\selectfont\catcode`\^=\active\def^{\ifmmode\sp\else\^{}\fi}\catcode`\%=\active\def%{\%}Jeddah 2024}}%
\end{pgfscope}%
\begin{pgfscope}%
\pgfsetbuttcap%
\pgfsetroundjoin%
\pgfsetlinewidth{1.505625pt}%
\definecolor{currentstroke}{rgb}{0.454902,0.678431,0.819608}%
\pgfsetstrokecolor{currentstroke}%
\pgfsetdash{}{0pt}%
\pgfpathmoveto{\pgfqpoint{2.410646in}{0.210055in}}%
\pgfpathlineto{\pgfqpoint{2.507869in}{0.210055in}}%
\pgfusepath{stroke}%
\end{pgfscope}%
\begin{pgfscope}%
\pgfsetbuttcap%
\pgfsetroundjoin%
\pgfsetlinewidth{1.505625pt}%
\definecolor{currentstroke}{rgb}{0.454902,0.678431,0.819608}%
\pgfsetstrokecolor{currentstroke}%
\pgfsetdash{}{0pt}%
\pgfpathmoveto{\pgfqpoint{2.459258in}{0.161444in}}%
\pgfpathlineto{\pgfqpoint{2.459258in}{0.258666in}}%
\pgfusepath{stroke}%
\end{pgfscope}%
\begin{pgfscope}%
\pgfsetbuttcap%
\pgfsetroundjoin%
\definecolor{currentfill}{rgb}{0.454902,0.678431,0.819608}%
\pgfsetfillcolor{currentfill}%
\pgfsetlinewidth{1.003750pt}%
\definecolor{currentstroke}{rgb}{0.454902,0.678431,0.819608}%
\pgfsetstrokecolor{currentstroke}%
\pgfsetdash{}{0pt}%
\pgfsys@defobject{currentmarker}{\pgfqpoint{0.000000in}{-0.069444in}}{\pgfqpoint{0.000000in}{0.069444in}}{%
\pgfpathmoveto{\pgfqpoint{0.000000in}{-0.069444in}}%
\pgfpathlineto{\pgfqpoint{0.000000in}{0.069444in}}%
\pgfusepath{stroke,fill}%
}%
\begin{pgfscope}%
\pgfsys@transformshift{2.410646in}{0.210055in}%
\pgfsys@useobject{currentmarker}{}%
\end{pgfscope}%
\end{pgfscope}%
\begin{pgfscope}%
\pgfsetbuttcap%
\pgfsetroundjoin%
\definecolor{currentfill}{rgb}{0.454902,0.678431,0.819608}%
\pgfsetfillcolor{currentfill}%
\pgfsetlinewidth{1.003750pt}%
\definecolor{currentstroke}{rgb}{0.454902,0.678431,0.819608}%
\pgfsetstrokecolor{currentstroke}%
\pgfsetdash{}{0pt}%
\pgfsys@defobject{currentmarker}{\pgfqpoint{0.000000in}{-0.069444in}}{\pgfqpoint{0.000000in}{0.069444in}}{%
\pgfpathmoveto{\pgfqpoint{0.000000in}{-0.069444in}}%
\pgfpathlineto{\pgfqpoint{0.000000in}{0.069444in}}%
\pgfusepath{stroke,fill}%
}%
\begin{pgfscope}%
\pgfsys@transformshift{2.507869in}{0.210055in}%
\pgfsys@useobject{currentmarker}{}%
\end{pgfscope}%
\end{pgfscope}%
\begin{pgfscope}%
\pgfsetbuttcap%
\pgfsetroundjoin%
\definecolor{currentfill}{rgb}{0.454902,0.678431,0.819608}%
\pgfsetfillcolor{currentfill}%
\pgfsetlinewidth{1.003750pt}%
\definecolor{currentstroke}{rgb}{0.454902,0.678431,0.819608}%
\pgfsetstrokecolor{currentstroke}%
\pgfsetdash{}{0pt}%
\pgfsys@defobject{currentmarker}{\pgfqpoint{-0.069444in}{-0.000000in}}{\pgfqpoint{0.069444in}{0.000000in}}{%
\pgfpathmoveto{\pgfqpoint{0.069444in}{-0.000000in}}%
\pgfpathlineto{\pgfqpoint{-0.069444in}{0.000000in}}%
\pgfusepath{stroke,fill}%
}%
\begin{pgfscope}%
\pgfsys@transformshift{2.459258in}{0.161444in}%
\pgfsys@useobject{currentmarker}{}%
\end{pgfscope}%
\end{pgfscope}%
\begin{pgfscope}%
\pgfsetbuttcap%
\pgfsetroundjoin%
\definecolor{currentfill}{rgb}{0.454902,0.678431,0.819608}%
\pgfsetfillcolor{currentfill}%
\pgfsetlinewidth{1.003750pt}%
\definecolor{currentstroke}{rgb}{0.454902,0.678431,0.819608}%
\pgfsetstrokecolor{currentstroke}%
\pgfsetdash{}{0pt}%
\pgfsys@defobject{currentmarker}{\pgfqpoint{-0.069444in}{-0.000000in}}{\pgfqpoint{0.069444in}{0.000000in}}{%
\pgfpathmoveto{\pgfqpoint{0.069444in}{-0.000000in}}%
\pgfpathlineto{\pgfqpoint{-0.069444in}{0.000000in}}%
\pgfusepath{stroke,fill}%
}%
\begin{pgfscope}%
\pgfsys@transformshift{2.459258in}{0.258666in}%
\pgfsys@useobject{currentmarker}{}%
\end{pgfscope}%
\end{pgfscope}%
\begin{pgfscope}%
\pgfsetbuttcap%
\pgfsetroundjoin%
\definecolor{currentfill}{rgb}{0.454902,0.678431,0.819608}%
\pgfsetfillcolor{currentfill}%
\pgfsetlinewidth{1.003750pt}%
\definecolor{currentstroke}{rgb}{0.454902,0.678431,0.819608}%
\pgfsetstrokecolor{currentstroke}%
\pgfsetdash{}{0pt}%
\pgfsys@defobject{currentmarker}{\pgfqpoint{-0.055556in}{-0.055556in}}{\pgfqpoint{0.055556in}{0.055556in}}{%
\pgfpathmoveto{\pgfqpoint{0.000000in}{-0.055556in}}%
\pgfpathcurveto{\pgfqpoint{0.014734in}{-0.055556in}}{\pgfqpoint{0.028866in}{-0.049702in}}{\pgfqpoint{0.039284in}{-0.039284in}}%
\pgfpathcurveto{\pgfqpoint{0.049702in}{-0.028866in}}{\pgfqpoint{0.055556in}{-0.014734in}}{\pgfqpoint{0.055556in}{0.000000in}}%
\pgfpathcurveto{\pgfqpoint{0.055556in}{0.014734in}}{\pgfqpoint{0.049702in}{0.028866in}}{\pgfqpoint{0.039284in}{0.039284in}}%
\pgfpathcurveto{\pgfqpoint{0.028866in}{0.049702in}}{\pgfqpoint{0.014734in}{0.055556in}}{\pgfqpoint{0.000000in}{0.055556in}}%
\pgfpathcurveto{\pgfqpoint{-0.014734in}{0.055556in}}{\pgfqpoint{-0.028866in}{0.049702in}}{\pgfqpoint{-0.039284in}{0.039284in}}%
\pgfpathcurveto{\pgfqpoint{-0.049702in}{0.028866in}}{\pgfqpoint{-0.055556in}{0.014734in}}{\pgfqpoint{-0.055556in}{0.000000in}}%
\pgfpathcurveto{\pgfqpoint{-0.055556in}{-0.014734in}}{\pgfqpoint{-0.049702in}{-0.028866in}}{\pgfqpoint{-0.039284in}{-0.039284in}}%
\pgfpathcurveto{\pgfqpoint{-0.028866in}{-0.049702in}}{\pgfqpoint{-0.014734in}{-0.055556in}}{\pgfqpoint{0.000000in}{-0.055556in}}%
\pgfpathlineto{\pgfqpoint{0.000000in}{-0.055556in}}%
\pgfpathclose%
\pgfusepath{stroke,fill}%
}%
\begin{pgfscope}%
\pgfsys@transformshift{2.459258in}{0.210055in}%
\pgfsys@useobject{currentmarker}{}%
\end{pgfscope}%
\end{pgfscope}%
\begin{pgfscope}%
\definecolor{textcolor}{rgb}{0.000000,0.000000,0.000000}%
\pgfsetstrokecolor{textcolor}%
\pgfsetfillcolor{textcolor}%
\pgftext[x=2.571063in,y=0.176027in,left,base]{\color{textcolor}{\rmfamily\fontsize{7.000000}{8.400000}\selectfont\catcode`\^=\active\def^{\ifmmode\sp\else\^{}\fi}\catcode`\%=\active\def%{\%}Spa 2024}}%
\end{pgfscope}%
\begin{pgfscope}%
\pgfsetbuttcap%
\pgfsetroundjoin%
\pgfsetlinewidth{1.505625pt}%
\definecolor{currentstroke}{rgb}{0.992157,0.682353,0.380392}%
\pgfsetstrokecolor{currentstroke}%
\pgfsetdash{}{0pt}%
\pgfpathmoveto{\pgfqpoint{3.214283in}{0.356421in}}%
\pgfpathlineto{\pgfqpoint{3.311505in}{0.356421in}}%
\pgfusepath{stroke}%
\end{pgfscope}%
\begin{pgfscope}%
\pgfsetbuttcap%
\pgfsetroundjoin%
\pgfsetlinewidth{1.505625pt}%
\definecolor{currentstroke}{rgb}{0.992157,0.682353,0.380392}%
\pgfsetstrokecolor{currentstroke}%
\pgfsetdash{}{0pt}%
\pgfpathmoveto{\pgfqpoint{3.262894in}{0.307810in}}%
\pgfpathlineto{\pgfqpoint{3.262894in}{0.405032in}}%
\pgfusepath{stroke}%
\end{pgfscope}%
\begin{pgfscope}%
\pgfsetbuttcap%
\pgfsetroundjoin%
\definecolor{currentfill}{rgb}{0.992157,0.682353,0.380392}%
\pgfsetfillcolor{currentfill}%
\pgfsetlinewidth{1.003750pt}%
\definecolor{currentstroke}{rgb}{0.992157,0.682353,0.380392}%
\pgfsetstrokecolor{currentstroke}%
\pgfsetdash{}{0pt}%
\pgfsys@defobject{currentmarker}{\pgfqpoint{0.000000in}{-0.069444in}}{\pgfqpoint{0.000000in}{0.069444in}}{%
\pgfpathmoveto{\pgfqpoint{0.000000in}{-0.069444in}}%
\pgfpathlineto{\pgfqpoint{0.000000in}{0.069444in}}%
\pgfusepath{stroke,fill}%
}%
\begin{pgfscope}%
\pgfsys@transformshift{3.214283in}{0.356421in}%
\pgfsys@useobject{currentmarker}{}%
\end{pgfscope}%
\end{pgfscope}%
\begin{pgfscope}%
\pgfsetbuttcap%
\pgfsetroundjoin%
\definecolor{currentfill}{rgb}{0.992157,0.682353,0.380392}%
\pgfsetfillcolor{currentfill}%
\pgfsetlinewidth{1.003750pt}%
\definecolor{currentstroke}{rgb}{0.992157,0.682353,0.380392}%
\pgfsetstrokecolor{currentstroke}%
\pgfsetdash{}{0pt}%
\pgfsys@defobject{currentmarker}{\pgfqpoint{0.000000in}{-0.069444in}}{\pgfqpoint{0.000000in}{0.069444in}}{%
\pgfpathmoveto{\pgfqpoint{0.000000in}{-0.069444in}}%
\pgfpathlineto{\pgfqpoint{0.000000in}{0.069444in}}%
\pgfusepath{stroke,fill}%
}%
\begin{pgfscope}%
\pgfsys@transformshift{3.311505in}{0.356421in}%
\pgfsys@useobject{currentmarker}{}%
\end{pgfscope}%
\end{pgfscope}%
\begin{pgfscope}%
\pgfsetbuttcap%
\pgfsetroundjoin%
\definecolor{currentfill}{rgb}{0.992157,0.682353,0.380392}%
\pgfsetfillcolor{currentfill}%
\pgfsetlinewidth{1.003750pt}%
\definecolor{currentstroke}{rgb}{0.992157,0.682353,0.380392}%
\pgfsetstrokecolor{currentstroke}%
\pgfsetdash{}{0pt}%
\pgfsys@defobject{currentmarker}{\pgfqpoint{-0.069444in}{-0.000000in}}{\pgfqpoint{0.069444in}{0.000000in}}{%
\pgfpathmoveto{\pgfqpoint{0.069444in}{-0.000000in}}%
\pgfpathlineto{\pgfqpoint{-0.069444in}{0.000000in}}%
\pgfusepath{stroke,fill}%
}%
\begin{pgfscope}%
\pgfsys@transformshift{3.262894in}{0.307810in}%
\pgfsys@useobject{currentmarker}{}%
\end{pgfscope}%
\end{pgfscope}%
\begin{pgfscope}%
\pgfsetbuttcap%
\pgfsetroundjoin%
\definecolor{currentfill}{rgb}{0.992157,0.682353,0.380392}%
\pgfsetfillcolor{currentfill}%
\pgfsetlinewidth{1.003750pt}%
\definecolor{currentstroke}{rgb}{0.992157,0.682353,0.380392}%
\pgfsetstrokecolor{currentstroke}%
\pgfsetdash{}{0pt}%
\pgfsys@defobject{currentmarker}{\pgfqpoint{-0.069444in}{-0.000000in}}{\pgfqpoint{0.069444in}{0.000000in}}{%
\pgfpathmoveto{\pgfqpoint{0.069444in}{-0.000000in}}%
\pgfpathlineto{\pgfqpoint{-0.069444in}{0.000000in}}%
\pgfusepath{stroke,fill}%
}%
\begin{pgfscope}%
\pgfsys@transformshift{3.262894in}{0.405032in}%
\pgfsys@useobject{currentmarker}{}%
\end{pgfscope}%
\end{pgfscope}%
\begin{pgfscope}%
\pgfsetbuttcap%
\pgfsetroundjoin%
\definecolor{currentfill}{rgb}{0.992157,0.682353,0.380392}%
\pgfsetfillcolor{currentfill}%
\pgfsetlinewidth{1.003750pt}%
\definecolor{currentstroke}{rgb}{0.992157,0.682353,0.380392}%
\pgfsetstrokecolor{currentstroke}%
\pgfsetdash{}{0pt}%
\pgfsys@defobject{currentmarker}{\pgfqpoint{-0.055556in}{-0.055556in}}{\pgfqpoint{0.055556in}{0.055556in}}{%
\pgfpathmoveto{\pgfqpoint{0.000000in}{-0.055556in}}%
\pgfpathcurveto{\pgfqpoint{0.014734in}{-0.055556in}}{\pgfqpoint{0.028866in}{-0.049702in}}{\pgfqpoint{0.039284in}{-0.039284in}}%
\pgfpathcurveto{\pgfqpoint{0.049702in}{-0.028866in}}{\pgfqpoint{0.055556in}{-0.014734in}}{\pgfqpoint{0.055556in}{0.000000in}}%
\pgfpathcurveto{\pgfqpoint{0.055556in}{0.014734in}}{\pgfqpoint{0.049702in}{0.028866in}}{\pgfqpoint{0.039284in}{0.039284in}}%
\pgfpathcurveto{\pgfqpoint{0.028866in}{0.049702in}}{\pgfqpoint{0.014734in}{0.055556in}}{\pgfqpoint{0.000000in}{0.055556in}}%
\pgfpathcurveto{\pgfqpoint{-0.014734in}{0.055556in}}{\pgfqpoint{-0.028866in}{0.049702in}}{\pgfqpoint{-0.039284in}{0.039284in}}%
\pgfpathcurveto{\pgfqpoint{-0.049702in}{0.028866in}}{\pgfqpoint{-0.055556in}{0.014734in}}{\pgfqpoint{-0.055556in}{0.000000in}}%
\pgfpathcurveto{\pgfqpoint{-0.055556in}{-0.014734in}}{\pgfqpoint{-0.049702in}{-0.028866in}}{\pgfqpoint{-0.039284in}{-0.039284in}}%
\pgfpathcurveto{\pgfqpoint{-0.028866in}{-0.049702in}}{\pgfqpoint{-0.014734in}{-0.055556in}}{\pgfqpoint{0.000000in}{-0.055556in}}%
\pgfpathlineto{\pgfqpoint{0.000000in}{-0.055556in}}%
\pgfpathclose%
\pgfusepath{stroke,fill}%
}%
\begin{pgfscope}%
\pgfsys@transformshift{3.262894in}{0.356421in}%
\pgfsys@useobject{currentmarker}{}%
\end{pgfscope}%
\end{pgfscope}%
\begin{pgfscope}%
\definecolor{textcolor}{rgb}{0.000000,0.000000,0.000000}%
\pgfsetstrokecolor{textcolor}%
\pgfsetfillcolor{textcolor}%
\pgftext[x=3.374700in,y=0.322394in,left,base]{\color{textcolor}{\rmfamily\fontsize{7.000000}{8.400000}\selectfont\catcode`\^=\active\def^{\ifmmode\sp\else\^{}\fi}\catcode`\%=\active\def%{\%}Silverstone 2024}}%
\end{pgfscope}%
\begin{pgfscope}%
\pgfsetbuttcap%
\pgfsetroundjoin%
\pgfsetlinewidth{1.505625pt}%
\definecolor{currentstroke}{rgb}{0.956863,0.427451,0.262745}%
\pgfsetstrokecolor{currentstroke}%
\pgfsetdash{}{0pt}%
\pgfpathmoveto{\pgfqpoint{3.214283in}{0.210055in}}%
\pgfpathlineto{\pgfqpoint{3.311505in}{0.210055in}}%
\pgfusepath{stroke}%
\end{pgfscope}%
\begin{pgfscope}%
\pgfsetbuttcap%
\pgfsetroundjoin%
\pgfsetlinewidth{1.505625pt}%
\definecolor{currentstroke}{rgb}{0.956863,0.427451,0.262745}%
\pgfsetstrokecolor{currentstroke}%
\pgfsetdash{}{0pt}%
\pgfpathmoveto{\pgfqpoint{3.262894in}{0.161444in}}%
\pgfpathlineto{\pgfqpoint{3.262894in}{0.258666in}}%
\pgfusepath{stroke}%
\end{pgfscope}%
\begin{pgfscope}%
\pgfsetbuttcap%
\pgfsetroundjoin%
\definecolor{currentfill}{rgb}{0.956863,0.427451,0.262745}%
\pgfsetfillcolor{currentfill}%
\pgfsetlinewidth{2.007500pt}%
\definecolor{currentstroke}{rgb}{0.956863,0.427451,0.262745}%
\pgfsetstrokecolor{currentstroke}%
\pgfsetdash{}{0pt}%
\pgfsys@defobject{currentmarker}{\pgfqpoint{0.000000in}{-0.069444in}}{\pgfqpoint{0.000000in}{0.069444in}}{%
\pgfpathmoveto{\pgfqpoint{0.000000in}{-0.069444in}}%
\pgfpathlineto{\pgfqpoint{0.000000in}{0.069444in}}%
\pgfusepath{stroke,fill}%
}%
\begin{pgfscope}%
\pgfsys@transformshift{3.214283in}{0.210055in}%
\pgfsys@useobject{currentmarker}{}%
\end{pgfscope}%
\end{pgfscope}%
\begin{pgfscope}%
\pgfsetbuttcap%
\pgfsetroundjoin%
\definecolor{currentfill}{rgb}{0.956863,0.427451,0.262745}%
\pgfsetfillcolor{currentfill}%
\pgfsetlinewidth{2.007500pt}%
\definecolor{currentstroke}{rgb}{0.956863,0.427451,0.262745}%
\pgfsetstrokecolor{currentstroke}%
\pgfsetdash{}{0pt}%
\pgfsys@defobject{currentmarker}{\pgfqpoint{0.000000in}{-0.069444in}}{\pgfqpoint{0.000000in}{0.069444in}}{%
\pgfpathmoveto{\pgfqpoint{0.000000in}{-0.069444in}}%
\pgfpathlineto{\pgfqpoint{0.000000in}{0.069444in}}%
\pgfusepath{stroke,fill}%
}%
\begin{pgfscope}%
\pgfsys@transformshift{3.311505in}{0.210055in}%
\pgfsys@useobject{currentmarker}{}%
\end{pgfscope}%
\end{pgfscope}%
\begin{pgfscope}%
\pgfsetbuttcap%
\pgfsetroundjoin%
\definecolor{currentfill}{rgb}{0.956863,0.427451,0.262745}%
\pgfsetfillcolor{currentfill}%
\pgfsetlinewidth{2.007500pt}%
\definecolor{currentstroke}{rgb}{0.956863,0.427451,0.262745}%
\pgfsetstrokecolor{currentstroke}%
\pgfsetdash{}{0pt}%
\pgfsys@defobject{currentmarker}{\pgfqpoint{-0.069444in}{-0.000000in}}{\pgfqpoint{0.069444in}{0.000000in}}{%
\pgfpathmoveto{\pgfqpoint{0.069444in}{-0.000000in}}%
\pgfpathlineto{\pgfqpoint{-0.069444in}{0.000000in}}%
\pgfusepath{stroke,fill}%
}%
\begin{pgfscope}%
\pgfsys@transformshift{3.262894in}{0.161444in}%
\pgfsys@useobject{currentmarker}{}%
\end{pgfscope}%
\end{pgfscope}%
\begin{pgfscope}%
\pgfsetbuttcap%
\pgfsetroundjoin%
\definecolor{currentfill}{rgb}{0.956863,0.427451,0.262745}%
\pgfsetfillcolor{currentfill}%
\pgfsetlinewidth{2.007500pt}%
\definecolor{currentstroke}{rgb}{0.956863,0.427451,0.262745}%
\pgfsetstrokecolor{currentstroke}%
\pgfsetdash{}{0pt}%
\pgfsys@defobject{currentmarker}{\pgfqpoint{-0.069444in}{-0.000000in}}{\pgfqpoint{0.069444in}{0.000000in}}{%
\pgfpathmoveto{\pgfqpoint{0.069444in}{-0.000000in}}%
\pgfpathlineto{\pgfqpoint{-0.069444in}{0.000000in}}%
\pgfusepath{stroke,fill}%
}%
\begin{pgfscope}%
\pgfsys@transformshift{3.262894in}{0.258666in}%
\pgfsys@useobject{currentmarker}{}%
\end{pgfscope}%
\end{pgfscope}%
\begin{pgfscope}%
\pgfsetbuttcap%
\pgfsetmiterjoin%
\definecolor{currentfill}{rgb}{0.000000,0.000000,0.000000}%
\pgfsetfillcolor{currentfill}%
\pgfsetfillopacity{0.000000}%
\pgfsetlinewidth{2.007500pt}%
\definecolor{currentstroke}{rgb}{0.956863,0.427451,0.262745}%
\pgfsetstrokecolor{currentstroke}%
\pgfsetdash{}{0pt}%
\pgfsys@defobject{currentmarker}{\pgfqpoint{-0.078567in}{-0.078567in}}{\pgfqpoint{0.078567in}{0.078567in}}{%
\pgfpathmoveto{\pgfqpoint{-0.000000in}{-0.078567in}}%
\pgfpathlineto{\pgfqpoint{0.078567in}{0.000000in}}%
\pgfpathlineto{\pgfqpoint{0.000000in}{0.078567in}}%
\pgfpathlineto{\pgfqpoint{-0.078567in}{0.000000in}}%
\pgfpathlineto{\pgfqpoint{-0.000000in}{-0.078567in}}%
\pgfpathclose%
\pgfusepath{stroke,fill}%
}%
\begin{pgfscope}%
\pgfsys@transformshift{3.262894in}{0.210055in}%
\pgfsys@useobject{currentmarker}{}%
\end{pgfscope}%
\end{pgfscope}%
\begin{pgfscope}%
\definecolor{textcolor}{rgb}{0.000000,0.000000,0.000000}%
\pgfsetstrokecolor{textcolor}%
\pgfsetfillcolor{textcolor}%
\pgftext[x=3.374700in,y=0.176027in,left,base]{\color{textcolor}{\rmfamily\fontsize{7.000000}{8.400000}\selectfont\catcode`\^=\active\def^{\ifmmode\sp\else\^{}\fi}\catcode`\%=\active\def%{\%}Yas Marina 2024 $\dagger$}}%
\end{pgfscope}%
\begin{pgfscope}%
\pgfsetbuttcap%
\pgfsetroundjoin%
\pgfsetlinewidth{1.505625pt}%
\definecolor{currentstroke}{rgb}{0.843137,0.188235,0.152941}%
\pgfsetstrokecolor{currentstroke}%
\pgfsetdash{}{0pt}%
\pgfpathmoveto{\pgfqpoint{4.272835in}{0.356421in}}%
\pgfpathlineto{\pgfqpoint{4.370058in}{0.356421in}}%
\pgfusepath{stroke}%
\end{pgfscope}%
\begin{pgfscope}%
\pgfsetbuttcap%
\pgfsetroundjoin%
\pgfsetlinewidth{1.505625pt}%
\definecolor{currentstroke}{rgb}{0.843137,0.188235,0.152941}%
\pgfsetstrokecolor{currentstroke}%
\pgfsetdash{}{0pt}%
\pgfpathmoveto{\pgfqpoint{4.321447in}{0.307810in}}%
\pgfpathlineto{\pgfqpoint{4.321447in}{0.405032in}}%
\pgfusepath{stroke}%
\end{pgfscope}%
\begin{pgfscope}%
\pgfsetbuttcap%
\pgfsetroundjoin%
\definecolor{currentfill}{rgb}{0.843137,0.188235,0.152941}%
\pgfsetfillcolor{currentfill}%
\pgfsetlinewidth{1.003750pt}%
\definecolor{currentstroke}{rgb}{0.843137,0.188235,0.152941}%
\pgfsetstrokecolor{currentstroke}%
\pgfsetdash{}{0pt}%
\pgfsys@defobject{currentmarker}{\pgfqpoint{0.000000in}{-0.069444in}}{\pgfqpoint{0.000000in}{0.069444in}}{%
\pgfpathmoveto{\pgfqpoint{0.000000in}{-0.069444in}}%
\pgfpathlineto{\pgfqpoint{0.000000in}{0.069444in}}%
\pgfusepath{stroke,fill}%
}%
\begin{pgfscope}%
\pgfsys@transformshift{4.272835in}{0.356421in}%
\pgfsys@useobject{currentmarker}{}%
\end{pgfscope}%
\end{pgfscope}%
\begin{pgfscope}%
\pgfsetbuttcap%
\pgfsetroundjoin%
\definecolor{currentfill}{rgb}{0.843137,0.188235,0.152941}%
\pgfsetfillcolor{currentfill}%
\pgfsetlinewidth{1.003750pt}%
\definecolor{currentstroke}{rgb}{0.843137,0.188235,0.152941}%
\pgfsetstrokecolor{currentstroke}%
\pgfsetdash{}{0pt}%
\pgfsys@defobject{currentmarker}{\pgfqpoint{0.000000in}{-0.069444in}}{\pgfqpoint{0.000000in}{0.069444in}}{%
\pgfpathmoveto{\pgfqpoint{0.000000in}{-0.069444in}}%
\pgfpathlineto{\pgfqpoint{0.000000in}{0.069444in}}%
\pgfusepath{stroke,fill}%
}%
\begin{pgfscope}%
\pgfsys@transformshift{4.370058in}{0.356421in}%
\pgfsys@useobject{currentmarker}{}%
\end{pgfscope}%
\end{pgfscope}%
\begin{pgfscope}%
\pgfsetbuttcap%
\pgfsetroundjoin%
\definecolor{currentfill}{rgb}{0.843137,0.188235,0.152941}%
\pgfsetfillcolor{currentfill}%
\pgfsetlinewidth{1.003750pt}%
\definecolor{currentstroke}{rgb}{0.843137,0.188235,0.152941}%
\pgfsetstrokecolor{currentstroke}%
\pgfsetdash{}{0pt}%
\pgfsys@defobject{currentmarker}{\pgfqpoint{-0.069444in}{-0.000000in}}{\pgfqpoint{0.069444in}{0.000000in}}{%
\pgfpathmoveto{\pgfqpoint{0.069444in}{-0.000000in}}%
\pgfpathlineto{\pgfqpoint{-0.069444in}{0.000000in}}%
\pgfusepath{stroke,fill}%
}%
\begin{pgfscope}%
\pgfsys@transformshift{4.321447in}{0.307810in}%
\pgfsys@useobject{currentmarker}{}%
\end{pgfscope}%
\end{pgfscope}%
\begin{pgfscope}%
\pgfsetbuttcap%
\pgfsetroundjoin%
\definecolor{currentfill}{rgb}{0.843137,0.188235,0.152941}%
\pgfsetfillcolor{currentfill}%
\pgfsetlinewidth{1.003750pt}%
\definecolor{currentstroke}{rgb}{0.843137,0.188235,0.152941}%
\pgfsetstrokecolor{currentstroke}%
\pgfsetdash{}{0pt}%
\pgfsys@defobject{currentmarker}{\pgfqpoint{-0.069444in}{-0.000000in}}{\pgfqpoint{0.069444in}{0.000000in}}{%
\pgfpathmoveto{\pgfqpoint{0.069444in}{-0.000000in}}%
\pgfpathlineto{\pgfqpoint{-0.069444in}{0.000000in}}%
\pgfusepath{stroke,fill}%
}%
\begin{pgfscope}%
\pgfsys@transformshift{4.321447in}{0.405032in}%
\pgfsys@useobject{currentmarker}{}%
\end{pgfscope}%
\end{pgfscope}%
\begin{pgfscope}%
\pgfsetbuttcap%
\pgfsetroundjoin%
\definecolor{currentfill}{rgb}{0.843137,0.188235,0.152941}%
\pgfsetfillcolor{currentfill}%
\pgfsetlinewidth{1.003750pt}%
\definecolor{currentstroke}{rgb}{0.843137,0.188235,0.152941}%
\pgfsetstrokecolor{currentstroke}%
\pgfsetdash{}{0pt}%
\pgfsys@defobject{currentmarker}{\pgfqpoint{-0.055556in}{-0.055556in}}{\pgfqpoint{0.055556in}{0.055556in}}{%
\pgfpathmoveto{\pgfqpoint{0.000000in}{-0.055556in}}%
\pgfpathcurveto{\pgfqpoint{0.014734in}{-0.055556in}}{\pgfqpoint{0.028866in}{-0.049702in}}{\pgfqpoint{0.039284in}{-0.039284in}}%
\pgfpathcurveto{\pgfqpoint{0.049702in}{-0.028866in}}{\pgfqpoint{0.055556in}{-0.014734in}}{\pgfqpoint{0.055556in}{0.000000in}}%
\pgfpathcurveto{\pgfqpoint{0.055556in}{0.014734in}}{\pgfqpoint{0.049702in}{0.028866in}}{\pgfqpoint{0.039284in}{0.039284in}}%
\pgfpathcurveto{\pgfqpoint{0.028866in}{0.049702in}}{\pgfqpoint{0.014734in}{0.055556in}}{\pgfqpoint{0.000000in}{0.055556in}}%
\pgfpathcurveto{\pgfqpoint{-0.014734in}{0.055556in}}{\pgfqpoint{-0.028866in}{0.049702in}}{\pgfqpoint{-0.039284in}{0.039284in}}%
\pgfpathcurveto{\pgfqpoint{-0.049702in}{0.028866in}}{\pgfqpoint{-0.055556in}{0.014734in}}{\pgfqpoint{-0.055556in}{0.000000in}}%
\pgfpathcurveto{\pgfqpoint{-0.055556in}{-0.014734in}}{\pgfqpoint{-0.049702in}{-0.028866in}}{\pgfqpoint{-0.039284in}{-0.039284in}}%
\pgfpathcurveto{\pgfqpoint{-0.028866in}{-0.049702in}}{\pgfqpoint{-0.014734in}{-0.055556in}}{\pgfqpoint{0.000000in}{-0.055556in}}%
\pgfpathlineto{\pgfqpoint{0.000000in}{-0.055556in}}%
\pgfpathclose%
\pgfusepath{stroke,fill}%
}%
\begin{pgfscope}%
\pgfsys@transformshift{4.321447in}{0.356421in}%
\pgfsys@useobject{currentmarker}{}%
\end{pgfscope}%
\end{pgfscope}%
\begin{pgfscope}%
\definecolor{textcolor}{rgb}{0.000000,0.000000,0.000000}%
\pgfsetstrokecolor{textcolor}%
\pgfsetfillcolor{textcolor}%
\pgftext[x=4.433252in,y=0.322394in,left,base]{\color{textcolor}{\rmfamily\fontsize{7.000000}{8.400000}\selectfont\catcode`\^=\active\def^{\ifmmode\sp\else\^{}\fi}\catcode`\%=\active\def%{\%}Suzuka 2024}}%
\end{pgfscope}%
\end{pgfpicture}%
\makeatother%
\endgroup%

    \caption[Aero polar sensitivity for Aston Martin AMR24 (driver 14, 2024). Left: all five circuits with Yas Marina flagged anomalous (open diamond). Right: four-circuit robustness check excluding Yas Marina. Error bars show $1\sigma$ total uncertainty on $C_DA$ (bootstrap-scaled) and $C_LA$ (statistical $\oplus$ $\mu$-systematic).]{Aero polar sensitivity for the Aston Martin AMR24 (driver~14, 2024). Error bars show $1\sigma$ total uncertainty on $C_DA$ (Jacobian $\times\,13\times$ bootstrap scale) and $C_LA$ (stat $\oplus$ $\mu$-systematic). \textbf{(a)}~All five circuits; Yas Marina (open diamond, $\dagger$) is a geometric outlier whose anomalous $C_DA$ ordering is unexplained (engine-braking contamination was falsified; see \S\ref{sec: Five Circuit Results}). OLS slope $= 3.02$; MC 90\% CI $= [-2.3, +6.0]$~m$^2$/m$^2$. \textbf{(b)}~Robustness check excluding Yas Marina. OLS slope $= 4.80$; MC 90\% CI $= [-2.9, +9.1]$~m$^2$/m$^2$. Both CIs include zero; neither fit constrains the slope to the positive physical regime.}
    \label{fig: aero polar}
\end{figure}

\subsection{DRS Drag Delta}
\label{sec: Results; DRS Drag Delta}

The DRS analysis is included to characterise the limits of the race-data method; as the confound analysis in \S\ref{sec: DRS Inseparable From Slipstream} shows, a clean DRS-only drag reduction is not isolatable from this data source.

The pooled Monza race-lap analysis yields 194 DRS-open and 733 DRS-closed classified car-laps. By convention, $\Delta C_DA < 0$ denotes reduced drag area (higher trap speed); $\Delta C_DA > 0$ would indicate increased drag. The pooled trap speed difference is:

\begin{equation}
    \Delta v_\text{trap} = 11.91 \pm 0.53~\text{km/h} \quad (90\%~\text{CI:}\ 10.87\text{--}12.95~\text{km/h})
\end{equation}

Converting via the energy balance gives $\Delta C_DA = -0.098$~m$^2$ ($90\%$ CI: $-0.26$ to $+0.05$~m$^2$). Note that the $\pm 0.53$~km/h standard error on $\Delta v_\text{trap}$ alone would map to a far tighter, strictly-negative $\Delta C_DA$ interval; the much wider CI quoted here is obtained by Monte Carlo over the full lap-to-lap trap-speed distributions of both groups (not from propagating the mean's SE), which is why its upper tail crosses zero. Physically, slipstream on DRS-open laps can inflate their trap speed enough to sign-flip the naive energy balance on individual laps, making DRS appear to add drag. The measured $-0.098$~m$^2$ is substantially smaller in magnitude than an unconfounded DRS drag reduction. The discrepancy is explained by slipstream confounding: every DRS-open lap in a race also has the following car within $1\,\text{s}$ of the car ahead, the same condition that guarantees aerodynamic wake exposure.

Stratifying by estimated gap to the car ahead (marginal-DRS group: $0.5-1.0\,\text{s}$, $n = 87$ laps) reduces the delta to $10.09$~km/h and $\Delta C_DA$ to $-0.079$~m$^2$, indicating some slipstream reduction but confirming that a pure DRS measurement is not achievable from race data alone.
